% Options for packages loaded elsewhere
\PassOptionsToPackage{unicode}{hyperref}
\PassOptionsToPackage{hyphens}{url}
\PassOptionsToPackage{dvipsnames,svgnames,x11names}{xcolor}
%
\documentclass[
  12pt]{article}

\usepackage{amsmath,amssymb,amsthm}
\usepackage{iftex}
\ifPDFTeX
  \usepackage[T1]{fontenc}
  \usepackage[utf8]{inputenc}
  \usepackage{textcomp} % provide euro and other symbols
\else % if luatex or xetex
  \usepackage{unicode-math}
  \defaultfontfeatures{Scale=MatchLowercase}
  \defaultfontfeatures[\rmfamily]{Ligatures=TeX,Scale=1}
\fi
\usepackage{lmodern}
\ifPDFTeX\else  
    % xetex/luatex font selection
\fi
% Use upquote if available, for straight quotes in verbatim environments
\IfFileExists{upquote.sty}{\usepackage{upquote}}{}
\IfFileExists{microtype.sty}{% use microtype if available
  \usepackage[]{microtype}
  \UseMicrotypeSet[protrusion]{basicmath} % disable protrusion for tt fonts
}{}
\makeatletter
\@ifundefined{KOMAClassName}{% if non-KOMA class
  \IfFileExists{parskip.sty}{%
    \usepackage{parskip}
  }{% else
    \setlength{\parindent}{0pt}
    \setlength{\parskip}{6pt plus 2pt minus 1pt}}
}{% if KOMA class
  \KOMAoptions{parskip=half}}
\makeatother
\usepackage{xcolor}
\setlength{\emergencystretch}{3em} % prevent overfull lines
\setcounter{secnumdepth}{5}
% Make \paragraph and \subparagraph free-standing
\makeatletter
\ifx\paragraph\undefined\else
  \let\oldparagraph\paragraph
  \renewcommand{\paragraph}{
    \@ifstar
      \xxxParagraphStar
      \xxxParagraphNoStar
  }
  \newcommand{\xxxParagraphStar}[1]{\oldparagraph*{#1}\mbox{}}
  \newcommand{\xxxParagraphNoStar}[1]{\oldparagraph{#1}\mbox{}}
\fi
\ifx\subparagraph\undefined\else
  \let\oldsubparagraph\subparagraph
  \renewcommand{\subparagraph}{
    \@ifstar
      \xxxSubParagraphStar
      \xxxSubParagraphNoStar
  }
  \newcommand{\xxxSubParagraphStar}[1]{\oldsubparagraph*{#1}\mbox{}}
  \newcommand{\xxxSubParagraphNoStar}[1]{\oldsubparagraph{#1}\mbox{}}
\fi
\makeatother


\providecommand{\tightlist}{%
  \setlength{\itemsep}{0pt}\setlength{\parskip}{0pt}}\usepackage{longtable,booktabs,array}
\usepackage{pdflscape}
\usepackage{calc} % for calculating minipage widths
% Correct order of tables after \paragraph or \subparagraph
\usepackage{etoolbox}
\makeatletter
\patchcmd\longtable{\par}{\if@noskipsec\mbox{}\fi\par}{}{}
\makeatother
% Allow footnotes in longtable head/foot
\IfFileExists{footnotehyper.sty}{\usepackage{footnotehyper}}{\usepackage{footnote}}
\makesavenoteenv{longtable}
\usepackage{graphicx}
\makeatletter
\def\maxwidth{\ifdim\Gin@nat@width>\linewidth\linewidth\else\Gin@nat@width\fi}
\def\maxheight{\ifdim\Gin@nat@height>\textheight\textheight\else\Gin@nat@height\fi}
\makeatother
% Scale images if necessary, so that they will not overflow the page
% margins by default, and it is still possible to overwrite the defaults
% using explicit options in \includegraphics[width, height, ...]{}
\setkeys{Gin}{width=\maxwidth,height=\maxheight,keepaspectratio}
% Set default figure placement to htbp
\makeatletter
\def\fps@figure{htbp}
\makeatother

\addtolength{\oddsidemargin}{-.5in}%
\addtolength{\evensidemargin}{-.1in}%
\addtolength{\textwidth}{1in}%
\addtolength{\textheight}{1.7in}%
\addtolength{\topmargin}{-1in}
\makeatletter
\@ifpackageloaded{caption}{}{\usepackage{caption}}
\AtBeginDocument{%
\ifdefined\contentsname
  \renewcommand*\contentsname{Table of contents}
\else
  \newcommand\contentsname{Table of contents}
\fi
\ifdefined\listfigurename
  \renewcommand*\listfigurename{List of Figures}
\else
  \newcommand\listfigurename{List of Figures}
\fi
\ifdefined\listtablename
  \renewcommand*\listtablename{List of Tables}
\else
  \newcommand\listtablename{List of Tables}
\fi
\ifdefined\figurename
  \renewcommand*\figurename{Figure}
\else
  \newcommand\figurename{Figure}
\fi
\ifdefined\tablename
  \renewcommand*\tablename{Table}
\else
  \newcommand\tablename{Table}
\fi
}
\@ifpackageloaded{float}{}{\usepackage{float}}
\floatstyle{ruled}
\@ifundefined{c@chapter}{\newfloat{codelisting}{h}{lop}}{\newfloat{codelisting}{h}{lop}[chapter]}
\floatname{codelisting}{Listing}
\newcommand*\listoflistings{\listof{codelisting}{List of Listings}}
\makeatother
\makeatletter
\makeatother
\makeatletter
\@ifpackageloaded{caption}{}{\usepackage{caption}}
\@ifpackageloaded{subcaption}{}{\usepackage{subcaption}}
\makeatother


\ifLuaTeX
  \usepackage{selnolig}  % disable illegal ligatures
\fi

% Additional packages and theorem environments transferred from the original imsart draft.
\usepackage{multirow}
\usepackage{mathtools}

\theoremstyle{plain}
\newtheorem{axiom}{Axiom}
\newtheorem{claim}[axiom]{Claim}
\newtheorem{theorem}{Theorem}[section]
\newtheorem{lemma}[theorem]{Lemma}
\newtheorem{corollary}[theorem]{Corollary}

\theoremstyle{definition}
\newtheorem{assumption}{Assumption}[section]
\newtheorem{condition}[assumption]{Condition}
\newtheorem{proposition}{Proposition}[section]
\newtheorem{definition}[theorem]{Definition}
\newtheorem*{example}{Example}
\newtheorem*{fact}{Fact}

\theoremstyle{remark}
\newtheorem{case}{Case}
\newtheorem{remark}[theorem]{Remark}

\usepackage[authoryear,round]{natbib}
\usepackage{bookmark}

\IfFileExists{xurl.sty}{\usepackage{xurl}}{} % add URL line breaks if available
\urlstyle{same} % disable monospaced font for URLs
\hypersetup{
  pdftitle={Simple Studentized l2 Statistics and Gaussian Bootstrap for Two-Sample Covariance Testing},
  pdfauthor={Haozhen Shu; Tianming Zhu},
  pdfkeywords={high-dimensional inference; Gaussian approximation; Gaussian bootstrap; studentization; transformed covariance statistics},
  colorlinks=true,
  linkcolor={blue},
  filecolor={Maroon},
  citecolor={Blue},
  urlcolor={Blue},
  pdfcreator={LaTeX via pandoc}}




\newcommand{\anon}{1}

%set the key \texttt{anon} to ``0'' to hide the authors and acknowledgements,
%  producing the required anonymized version. 
%Set the key \texttt{anon} to ``1'' to produce the manuscript with author details and
% acknowledgments. 


\begin{document}

\def\spacingset#1{\renewcommand{\baselinestretch}%
{#1}\small\normalsize} \spacingset{1}

%%%%%%%%%%%%%%%%%%%%%%%%%%%%%%%%%%%%%%%%%%%%%%%%%%%%%%%%%%%%%%%%%%%%%%%%%%%%%%

\if1\anon
{
  \title{\bf Simple Studentized \texorpdfstring{$\ell_2$}{l2} Statistics and Gaussian Bootstrap for Two-Sample Covariance Testing}
  \author{Haozhen Shu\thanks{Email: haozhen.shu@nie.edu.sg.}\hspace{.2cm}\\
Department, University or Company Name\\
and \\
Tianming Zhu\thanks{Email: tianming.zhu@nie.edu.sg.}\\
Department, University or Company Name}
  \maketitle
} \fi

\if0\anon
{
  \bigskip
  \bigskip
  \bigskip
  \begin{center}
    {\LARGE\bf Simple Studentized \texorpdfstring{$\ell_2$}{l2} Statistics and Gaussian Bootstrap for Two-Sample Covariance Testing}
\end{center}
  \medskip
} \fi

\bigskip
\begin{abstract}
We propose a simple studentized \(\ell_2\)-type statistic for high-dimensional
two-sample covariance testing.  The statistic vectorizes the difference between
the two sample covariance matrices, standardizes each covariance coordinate by
its marginal scale, and aggregates the squared standardized discrepancies. The theoretical analysis is based on a Gaussian approximation over centered
Euclidean balls, which matches the geometry of the proposed statistic and leads
to a direct Gaussian multiplier-bootstrap implementation.  We establish validity of the
feasible bootstrap test and power consistency under growing standardized
quadratic signal strength.  The method is easy to implement, avoids estimating
the full covariance matrix of covariance entries, and can incorporate
prespecified nonsingular transformations to improve finite-sample performance.
\end{abstract}

\noindent%
{\it Keywords:} high-dimensional inference; Gaussian approximation; Gaussian bootstrap; multiplier bootstrap; transformed covariance statistics
\vfill

\newpage
\spacingset{1.8} % DON'T change the spacing!

\section{Introduction and Review}\label{sec:intro}

Covariance comparison is a fundamental problem in multivariate statistics.  It
arises when the scientific question concerns not only marginal variation, but
also the dependence structure among variables.  Examples include
gene-expression studies, where genes may act jointly through co-expression
modules; neuroimaging studies, where functional connectivity is encoded by
covariance or correlation patterns across brain regions; finance, where
portfolio risk depends on the joint movement of asset returns; and network
analysis, where second-order dependence captures interaction patterns among
features.  In such settings, the covariance matrix is itself an object of
scientific interest.

Modern covariance testing is challenging because the dimension may be
comparable to, or larger than, the sample size.  Classical likelihood-ratio
procedures, including the likelihood criteria studied by
\citet{box_general_1949} and related covariance-matrix test criteria of
\citet{nagao_some_1973}, rely on fixed-dimensional approximations and may
become unstable when sample covariance matrices are singular or
ill-conditioned.  This difficulty motivated high-dimensional corrections for
covariance structure testing, including the large-dimensional tests of
\citet{ledoit_some_2002}, as well as the nonparametric identity and sphericity
tests of \citet{chen_tests_2010}.  It has also led to
a broad literature on tests that avoid direct likelihood inversion and compare
covariance matrices through entrywise, spectral, or quadratic functionals.

From the perspective of power against different alternatives, many procedures
are often discussed along a sparse-versus-dense contrast.  Max-type tests are
designed to detect alternatives where a small number of covariance entries
differ sharply.  This direction is connected to
the limiting behavior of sample correlations and coherence statistics
\citep{cai_limiting_2011}, and to entrywise maximum tests for covariance
comparison and support recovery
\citep{cai_two-sample_2013,chang_comparing_2017}.  Its modern calibration is
closely related to high-dimensional Gaussian approximation and multiplier
bootstrap theory for maxima and hyperrectangles
\citep{chernozhukov_gaussian_2013,chernozhukov_central_2017,
chernozhukov_high-dimensional_2023}.  In contrast, \(\ell_2\)-type procedures
aggregate discrepancies across many covariance coordinates and are therefore
natural for dense or moderately dense alternatives.  Early and influential
quadratic-form tests include those of \citet{schott_test_2007} and
\citet{srivastava_testing_2010}, while \citet{li_two_2012} proposed an
unbiased U-statistic estimator of the squared Frobenius-norm covariance
difference.  More recent work has further developed normal-reference and
quadratic-form approaches for high-dimensional covariance testing
\citep{wang_two-sample_2024}.

U-statistics play an important role in this literature because they provide
unbiased estimators of quadratic functionals and remove the leading bias that
appears in naive plug-in estimators such as
\(\|\widehat\Sigma_1-\widehat\Sigma_2\|_F^2\).  The Li--Chen statistic is a
canonical example in two-sample covariance testing \citep{li_two_2012}, and
its debiasing idea is related to high-dimensional two-sample mean testing based
on quadratic U-statistics \citep{chen_qin_2010}.  More generally,
high-dimensional Gaussian and bootstrap approximations for non-degenerate
U-statistics have been studied by \citet{chen_gaussian_2018}, with related
exchangeable-pair bounds by \citet{cheng_gaussian_2022}.
These results are highly relevant, but they do not directly yield the
approximation needed in the present paper.  Standard U-statistic asymptotics
often distinguishes non-degenerate kernels, where a leading Hájek projection
drives the limit, from degenerate kernels, where higher-order components can
become central.  Our statistic has a different geometric target.

To see the distinction, consider the oracle version of our statistic under the
null.  After diagonal standardization, it can be written as the squared
Euclidean norm of a high-dimensional covariance-coordinate mean vector.  If one
expands this squared norm, the off-diagonal part has a degenerate
second-order-kernel form, while the diagonal part and the studentization remain
essential to the statistic.  Thus a direct U-statistic normal approximation is
not the right object: the null distribution is naturally described as a
quadratic form of a Gaussian vector, or equivalently through probabilities of
centered Euclidean balls.  This is the main reason we develop the theory through
a centered-ball Gaussian approximation for the covariance-coordinate vector,
rather than through a scalar U-statistic central limit theorem.

We propose a simple studentized \(\ell_2\)-type statistic for two-sample
covariance testing.  The statistic vectorizes the upper-triangular entries of
the two sample covariance matrices, forms their two-sample difference,
standardizes each covariance coordinate by its marginal scale, and aggregates
the squared standardized discrepancies.  This produces a scale-invariant global
measure of covariance difference.  The construction is particularly suited to
dense alternatives, where many covariance entries change moderately rather than
a few entries changing sharply.

The key theoretical idea is to analyze the standardized covariance-coordinate
vector before applying the squared-norm map.  Events of the form
\[
  \left\{ d^{-1}\|Z\|_2^2 \le t \right\}
\]
are exactly centered Euclidean balls; see the statistic in
\eqref{eq:proposed-statistic} and its acceptance-region representation in
\eqref{eq:acceptance-ball}.  Therefore a Gaussian approximation over centered
balls matches the geometry of the proposed statistic.  This approach leads to a
direct multiplier-bootstrap calibration: the bootstrap is applied to the same
covariance-coordinate contributions used in the statistic, and only a diagonal
studentizer is needed.  In particular, the procedure avoids estimating or
inverting the full covariance matrix of covariance entries.

Our contributions are as follows.  First, we introduce a simple and
interpretable studentized \(\ell_2\) statistic for high-dimensional two-sample
covariance testing.  Second, we establish a Gaussian approximation over
centered Euclidean balls for the oracle statistic, with explicit dependence on
dimension, moments, and covariance geometry.  Third, we prove that the feasible
statistic, based on empirical centering and empirical diagonal studentization,
is close to its oracle counterpart.  Fourth, we develop a multiplier-bootstrap
calibration and prove asymptotic size validity under explicit sufficient
conditions.  Fifth, we derive power consistency under alternatives with growing
standardized quadratic signal strength.  The framework can also incorporate
common nonsingular transformations, allowing external scientific information or
algorithmic representations to make covariance differences more visible after
studentization.  The formal theory treats such transformations as fixed and
deterministic, while fully adaptive same-sample selection is left as a separate
problem.

The remainder of the paper is organized as follows.
Section~\ref{sec:main_results} presents the problem setting, the proposed
studentized statistic, the centered-ball Gaussian approximation, the
sample-to-oracle comparison, bootstrap calibration, size control, and power
analysis.  Section~\ref{sec:simulation} reports simulation studies comparing
the proposed bootstrap test with max-type and quadratic-form benchmarks.
Section~\ref{sec:realdata} presents an application to TCGA-BRCA gene-expression
data.  Section~\ref{sec:discussion} concludes with discussion and open
directions.

\section{Main Results}
\label{sec:main_results}

This section develops the proposed studentized two-sample covariance test under
a deterministic linear transformation.  We first introduce the testing problem
\eqref{eq:hypothesis}, the null-preserving transformation
\eqref{eq:transformed-null}, and the statistic \eqref{eq:proposed-statistic}.
We then establish the oracle Gaussian approximation, the sample-to-oracle
transfer, the multiplier-bootstrap calibration, size validity, and power
consistency.  We close the section by highlighting how the simplicity of 
the proposed quadratic form leads to a transparent 
inferential interpretation. In highly unbalanced designs, 
the statistic can be naturally viewed as a one-sample covariance 
discrepancy, with the larger sample providing 
a stable reference covariance for assessing the smaller sample.


\subsection{Problem setting}
\label{subsec:problem_setting}

Let
\[
X_1,\ldots,X_{n_1}\in\mathbb R^p,
\qquad
Y_1,\ldots,Y_{n_2}\in\mathbb R^p
\]
be two independent samples, with observations independent and identically
distributed within each sample.  Denote
\[
\mathbb E X_i=\mu_1,\qquad \operatorname{Cov}(X_i)=\Sigma_1,
\]
and
\[
\mathbb E Y_j=\mu_2,\qquad \operatorname{Cov}(Y_j)=\Sigma_2.
\]
We consider the two-sample covariance testing problem
\begin{equation}
\label{eq:hypothesis}
H_0:\Sigma_1=\Sigma_2
\qquad\text{against}\qquad
H_1:\Sigma_1\neq\Sigma_2.
\end{equation}

Throughout the theoretical analysis, \(A\in\mathbb R^{p\times p}\) is treated
as a deterministic nonsingular matrix.  Define the transformed observations
\[
X_i^{(A)}=AX_i,\qquad Y_j^{(A)}=AY_j,
\]
with transformed covariance matrices
\[
\Sigma_1^{(A)}=A\Sigma_1A^\top,
\qquad
\Sigma_2^{(A)}=A\Sigma_2A^\top.
\]
Since \(A\) is nonsingular,
\begin{equation}
\label{eq:transformed-null}
\Sigma_1=\Sigma_2
\quad\Longleftrightarrow\quad
\Sigma_1^{(A)}=\Sigma_2^{(A)}.
\end{equation}
Thus the equivalence in \eqref{eq:transformed-null} shows that applying the
same deterministic transformation to both samples preserves the null
hypothesis.  The transformation changes the geometric representation of the
covariance difference but not the null itself.

The matrix \(A\) should therefore be viewed as an analysis-side representation
choice rather than as a modification of the hypothesis.  The baseline choice
\(A=I_p\) tests covariance equality in the original coordinates.  Other
nonsingular choices can encode external scientific knowledge, pre-existing
feature maps, or transformations suggested by machine-learning preprocessing.
In this sense, \(A\) offers a representation-augmentation mechanism: the same
null-preserving test can be applied after re-expressing the data in coordinates
where the anticipated covariance discrepancy may be denser, better scaled, or
more aligned with the \(\ell_2\) aggregation.  The validity results below are
stated for fixed deterministic \(A\).  If \(A\) is estimated or selected from
the same observations used for testing, additional sample-splitting,
conditioning, or post-selection analysis would be required.

After introducing the transformation, we suppress the superscript \((A)\) and
work entirely in the transformed model.  Thus \(X_i,Y_j,\mu_1,\mu_2,\Sigma_1\),
and \(\Sigma_2\) below denote their transformed counterparts.  Let
\begin{equation}
\label{eq:dimension-effective-n}
d=\frac{p(p+1)}2,
\qquad
n_{\mathrm{eff}}=\frac{n_1n_2}{n_1+n_2}.
\end{equation}
Here \(d\) is the number of upper-triangular covariance coordinates, and
\(\operatorname{vech}(\cdot)\) denotes the half-vectorization operator that
stacks the upper-triangular entries of a symmetric matrix, including the
diagonal.  The effective sample size \(n_{\mathrm{eff}}\) in
\eqref{eq:dimension-effective-n} is the scale that appears in the covariance
difference statistic below.

We impose the following primitive regularity conditions in the transformed
coordinate system.  Throughout the theoretical analysis, write
\[
U_i=X_i-\mu_1,
\qquad
V_j=Y_j-\mu_2.
\]

\begin{assumption}[Coordinate high-moment control]
\label{ass:coordinate_high_moment}
There exists a constant \(K_8<\infty\), independent of \(n_1,n_2\), and \(p\),
such that
\[
\max_{1\le a\le p}\mathbb E|U_{ia}|^8
+
\max_{1\le a\le p}\mathbb E|V_{ja}|^8
\le K_8 .
\]
\end{assumption}

\begin{assumption}[Covariance-coordinate nondegeneracy]
\label{ass:coordinate_nondegeneracy}
There exists a constant \(c_H>0\), independent of \(n_1,n_2\), and \(p\), such
that for every \(1\le a\le b\le p\),
\[
\operatorname{Var}(U_{ia}U_{ib})\ge c_H,
\qquad
\operatorname{Var}(V_{ja}V_{jb})\ge c_H .
\]
\end{assumption}

Assumption~\ref{ass:coordinate_high_moment} is a coordinatewise moment
condition used to control covariance-coordinate estimators uniformly over the
\(d=p(p+1)/2\) entries.  It also implies the Euclidean moment and sample-mean
bounds needed for sample-centering arguments; the short derivation is given in
the supplement.  Assumption~\ref{ass:coordinate_nondegeneracy} excludes
degenerate covariance coordinates, so that the diagonal studentizer is well
defined and stable.  These assumptions are primitive because they are stated
directly in terms of transformed observations and covariance-coordinate
variances, rather than in terms of the limiting statistic itself.

\subsection{Statistic and Gaussian approximation}
\label{subsec:statistic_gaussian}

The sample means and sample covariance matrices are
\[
\bar X=\frac1{n_1}\sum_{i=1}^{n_1}X_i,
\qquad
\bar Y=\frac1{n_2}\sum_{j=1}^{n_2}Y_j,
\]
and
\[
\widehat\Sigma_1
=
\frac1{n_1}\sum_{i=1}^{n_1}(X_i-\bar X)(X_i-\bar X)^\top,
\]
\[
\widehat\Sigma_2
=
\frac1{n_2}\sum_{j=1}^{n_2}(Y_j-\bar Y)(Y_j-\bar Y)^\top.
\]
Define the scaled covariance-difference vector
\begin{equation}
\label{eq:scaled-difference}
\widehat W
=
\sqrt{n_{\mathrm{eff}}}\,
\operatorname{vech}(\widehat\Sigma_1-\widehat\Sigma_2)
\in\mathbb R^d.
\end{equation}

For \(k=k(a,b)\), \(1\le a\le b\le p\), define
\[
\widehat h_{i,k}^{(1)}
=
(X_{ia}-\bar X_a)(X_{ib}-\bar X_b)
-
\widehat\Sigma_{1,ab},
\]
and
\[
\widehat h_{j,k}^{(2)}
=
(Y_{ja}-\bar Y_a)(Y_{jb}-\bar Y_b)
-
\widehat\Sigma_{2,ab}.
\]
The diagonal variance estimator is
\begin{equation}
\label{eq:diagonal-studentizer}
\widehat D_{kk}
=
n_{\mathrm{eff}}
\left(
\frac1{n_1^2}
\sum_{i=1}^{n_1}
\bigl(\widehat h_{i,k}^{(1)}\bigr)^2
+
\frac1{n_2^2}
\sum_{j=1}^{n_2}
\bigl(\widehat h_{j,k}^{(2)}\bigr)^2
\right),
\qquad 1\le k\le d.
\end{equation}
Let
\[
\widehat D=\operatorname{diag}(\widehat D_{11},\ldots,\widehat D_{dd}).
\]
The proposed statistic is
\begin{equation}
\label{eq:proposed-statistic}
\widehat T_n
=
\frac1d
\widehat W^\top\widehat D^{-1}\widehat W
=
\frac1d
\left\|
\widehat D^{-1/2}\widehat W
\right\|_2^2.
\end{equation}
Large values of \(\widehat T_n\) in \eqref{eq:proposed-statistic} provide
evidence against \(H_0\).

The statistic is an \(\ell_2\)-type aggregation of standardized covariance
differences based on the scaled difference \eqref{eq:scaled-difference} and
diagonal studentizer \eqref{eq:diagonal-studentizer}.  Its acceptance region
satisfies
\begin{equation}
\label{eq:acceptance-ball}
\{\widehat T_n\le t\}
=
\left\{
\left\|
\widehat D^{-1/2}\widehat W
\right\|_2
\le \sqrt{dt}
\right\}.
\end{equation}
Thus, after diagonal standardization, the relevant geometry is Euclidean-ball
geometry.

For theoretical analysis, define the oracle covariance estimators
\[
\widetilde\Sigma_1
=
\frac1{n_1}\sum_{i=1}^{n_1}(X_i-\mu_1)(X_i-\mu_1)^\top,
\]
\[
\widetilde\Sigma_2
=
\frac1{n_2}\sum_{j=1}^{n_2}(Y_j-\mu_2)(Y_j-\mu_2)^\top,
\]
and
\[
\widetilde W
=
\sqrt{n_{\mathrm{eff}}}\,
\operatorname{vech}(\widetilde\Sigma_1-\widetilde\Sigma_2).
\]

It is useful to introduce the associated covariance-coordinate notation.  For
\(g=1,2\), define
\[
H_i^{(1)}
=
\operatorname{vech}\{U_iU_i^\top-\Sigma_1\},
\qquad
H_j^{(2)}
=
\operatorname{vech}\{V_jV_j^\top-\Sigma_2\},
\]
and write
\[
\Gamma_g=\operatorname{Cov}(H^{(g)}),
\qquad
D_g=\operatorname{diag}(\Gamma_g),
\]
where \(H^{(1)}\) and \(H^{(2)}\) denote generic copies of \(H_i^{(1)}\) and
\(H_j^{(2)}\), respectively.  Also set
\[
\bar H_1=\frac1{n_1}\sum_{i=1}^{n_1}H_i^{(1)},
\qquad
\bar H_2=\frac1{n_2}\sum_{j=1}^{n_2}H_j^{(2)}.
\]
Under \(H_0\),
\[
\widetilde W
=
\sqrt{n_{\mathrm{eff}}}(\bar H_1-\bar H_2).
\]

Let
\[
\widetilde D
=
\operatorname{diag}
\{
\operatorname{Var}(\widetilde W_1),\ldots,
\operatorname{Var}(\widetilde W_d)
\}.
\]
The oracle statistic is
\begin{equation}
\label{eq:oracle-statistic}
\widetilde T_n
=
\frac1d
\left\|
\widetilde D^{-1/2}\widetilde W
\right\|_2^2.
\end{equation}
Under \(H_0\), the standardized vector
\[
\widetilde D^{-1/2}\widetilde W
\]
is a sum of independent centered random vectors.  We collect its covariance
matrix and the associated oracle Gaussian reference statistic as
\begin{equation}
\label{eq:gaussian-reference}
\begin{aligned}
R
&=
\widetilde D^{-1/2}
\operatorname{Cov}(\widetilde W)
\widetilde D^{-1/2},\\
T_Z
&=
\frac1d\|Z\|_2^2,
\qquad
Z\sim N(0,R).
\end{aligned}
\end{equation}
Accordingly, \eqref{eq:gaussian-reference} defines the oracle Gaussian
reference statistic associated with the oracle statistic
\eqref{eq:oracle-statistic}.

Let
\[
\lambda_1\ge\lambda_2\ge\cdots\ge\lambda_d\ge0
\]
be the eigenvalues of \(R\).  For \(s=1,2\), define
\begin{equation}
\label{eq:reference-eigen-geometry}
\Lambda_s
=
\left(\sum_{\ell=s}^d\lambda_\ell^2\right)^{1/2},
\qquad
\kappa(R)=(\Lambda_1\Lambda_2)^{-1/2}
\quad\text{whenever }\Lambda_2>0.
\end{equation}

\begin{assumption}[Gaussian reference geometry]
\label{ass:gaussian_reference_geometry}
The oracle Gaussian reference matrix \(R\) satisfies
\[
\Lambda_2>0 .
\]
\end{assumption}

Assumption~\ref{ass:gaussian_reference_geometry} is a mild geometric condition
on the Gaussian quadratic reference, expressed through
\eqref{eq:reference-eigen-geometry}.  It rules out the degenerate case in which
the reference variation is concentrated in a single direction, where
anti-concentration for Euclidean balls becomes ill-posed.

The first theorem establishes that the oracle statistic
\eqref{eq:oracle-statistic} is approximated by the Gaussian quadratic reference
\eqref{eq:gaussian-reference}.

\begin{theorem}[Oracle Gaussian approximation]
\label{thm:oracle_gaussian}
Suppose \(H_0:\Sigma_1=\Sigma_2\) holds.  Suppose
Assumptions~\ref{ass:coordinate_high_moment},
\ref{ass:coordinate_nondegeneracy}, and
\ref{ass:gaussian_reference_geometry} hold.  Then there exists a constant
\(K>0\), depending only on the moment constants and the nondegeneracy constant
in the assumptions, such that
\[
\sup_{r\ge0}
\left|
\mathbb P(\widetilde T_n\le r)
-
\mathbb P(T_Z\le r)
\right|
\le
K
\left\{
\kappa(R)^{3/4}\delta_1^{1/4}
+
\kappa(R)^{2/3}\delta_2^{1/3}
\right\},
\]
where
\[
\delta_1
=
\|R\|_{\mathrm{HS}}
d^2
\left(n_1^{-1}+n_2^{-1}\right)
+
\|R\|_{\mathrm{op}}^{3/2}
d^{3/2}
\left(n_1^{-1/2}+n_2^{-1/2}\right),
\]
and
\[
\delta_2
=
\|R\|_{\mathrm{op}}^{1/2}
d^{3/2}
\left(n_1^{-1/2}+n_2^{-1/2}\right)
+
d^2
\left(n_1^{-1}+n_2^{-1}\right).
\]
\end{theorem}

The oracle approximation in Theorem~\ref{thm:oracle_gaussian} is stated for
\(\widetilde T_n\), which depends on the unknown population means and the
oracle diagonal variance matrix.  We next show that the feasible statistic
\(\widehat T_n\), obtained by replacing the population means and the oracle
studentizer by their empirical counterparts, remains close to \(\widetilde T_n\).
The proof relies on two ingredients.  The first is the sample-centering identity
\[
\widehat\Sigma_m-\widetilde\Sigma_m
=
-(\bar X_m-\mu_m)(\bar X_m-\mu_m)^\top,
\qquad m=1,2,
\]
where \((\bar X_1,\mu_1)=(\bar X,\mu_1)\) and
\((\bar X_2,\mu_2)=(\bar Y,\mu_2)\).  The second is the uniform consistency of
the diagonal studentizer:
\[
\max_{1\le k\le d}
\left|
\frac{\widehat D_{kk}}{\widetilde D_{kk}}-1
\right|
=
O_{\mathbb P}
\left(
\sqrt{\frac d{n_0}}
\right),
\qquad
n_0=n_1\wedge n_2.
\]
This studentizer consistency follows from the primitive moment and
covariance-coordinate nondegeneracy assumptions; its technical verification is
deferred to the appendix.

For \(\eta\in(0,1)\), define
\[
B_T(\eta)=K_0\eta^{-1/4},
\]
\[
B_W(\eta)
=
K_0
\frac{\sqrt{n_{\mathrm{eff}}}}{\sqrt d}
\eta^{-1/2}
\left(
\frac{p}{n_1}
+
\frac{p}{n_2}
\right),
\]
and
\[
B_D(\eta)
=
K_0
\left(
\frac{d}{n_0\eta}
\right)^{1/2},
\qquad
n_0=n_1\wedge n_2,
\]
where \(K_0>0\) denotes a generic constant depending only on the constants in the
assumptions.  Set
\[
\Delta(\eta)
=
2B_D(\eta)\{B_T(\eta)+B_W(\eta)\}^2
+
2B_T(\eta)B_W(\eta)
+
B_W^2(\eta).
\]

\begin{proposition}[Consistency of the feasible statistic]
\label{prop:sample_oracle_transfer}
Suppose \(H_0:\Sigma_1=\Sigma_2\) holds.  Suppose
Assumptions~\ref{ass:coordinate_high_moment} and
\ref{ass:coordinate_nondegeneracy} hold.  If \(B_D(\eta)<1/2\), then with probability at least
\(1-\eta\),
\[
|\widehat T_n-\widetilde T_n|
\le
\Delta(\eta).
\]
In particular, taking
\[
\eta=n_0^{-1/8},
\]
we have, with probability at least \(1-n_0^{-1/8}\),
\[
|\widehat T_n-\widetilde T_n|
\le
K_0
\left(\frac{d}{n_0^{3/4}}\right)^{1/2}.
\]
\end{proposition}

Combining Theorem~\ref{thm:oracle_gaussian} with
Proposition~\ref{prop:sample_oracle_transfer} gives a Gaussian approximation
for the feasible statistic.  Let
\[
\varepsilon_{\mathrm{or}}
=
K
\left\{
\kappa(R)^{3/4}\delta_1^{1/4}
+
\kappa(R)^{2/3}\delta_2^{1/3}
\right\}.
\]
Define the feasible approximation error
\begin{equation}
\label{eq:feasible-gaussian-error}
\varepsilon_{\mathrm{feas}}(\eta)
=
\varepsilon_{\mathrm{or}}
+
\eta
+
K\kappa(R)d\Delta(\eta).
\end{equation}

\begin{corollary}[Feasible Gaussian approximation]
\label{cor:feasible_gaussian}
Suppose \(H_0:\Sigma_1=\Sigma_2\) holds.  Suppose
Assumptions~\ref{ass:coordinate_high_moment},
\ref{ass:coordinate_nondegeneracy}, and
\ref{ass:gaussian_reference_geometry} hold.  If
\[
B_D(\eta)\le \frac12,
\]
then
\[
\sup_{r\ge0}
\left|
\mathbb P(\widehat T_n\le r)
-
\mathbb P(T_Z\le r)
\right|
\le
\varepsilon_{\mathrm{feas}}(\eta).
\]
\end{corollary}

\subsection{Bootstrap calibration}
\label{subsec:bootstrap_calibration}

The oracle Gaussian reference depends on the unknown covariance matrix \(R\).
We therefore calibrate the test by a Gaussian multiplier bootstrap.  Let
\[
e_1,\ldots,e_{n_1},
\qquad
e_1',\ldots,e_{n_2}'
\]
be i.i.d. \(N(0,1)\) variables independent of the data.  Define
\[
\widehat h_i^{(1)}
=
(\widehat h_{i,1}^{(1)},\ldots,\widehat h_{i,d}^{(1)})^\top,
\qquad
\widehat h_j^{(2)}
=
(\widehat h_{j,1}^{(2)},\ldots,\widehat h_{j,d}^{(2)})^\top.
\]
The multiplier bootstrap vector is
\begin{equation}
\label{eq:bootstrap-vector}
\widehat W^{\mathrm{mb}}
=
\sqrt{n_{\mathrm{eff}}}
\left\{
\frac1{n_1}\sum_{i=1}^{n_1}e_i\widehat h_i^{(1)}
-
\frac1{n_2}\sum_{j=1}^{n_2}e_j'\widehat h_j^{(2)}
\right\}.
\end{equation}
The bootstrap statistic is
\begin{equation}
\label{eq:bootstrap-statistic}
\widehat T_n^{\mathrm{mb}}
=
\frac1d
\left\|
\widehat D^{-1/2}\widehat W^{\mathrm{mb}}
\right\|_2^2.
\end{equation}
Conditionally on the data,
\[
\widehat D^{-1/2}\widehat W^{\mathrm{mb}}
\sim
N(0,\widehat R),
\]
where
\[
\widehat R
=
\widehat D^{-1/2}\widehat\Gamma\widehat D^{-1/2},
\]
and
\[
\widehat\Gamma
=
n_{\mathrm{eff}}
\left\{
\frac1{n_1^2}\sum_{i=1}^{n_1}
\widehat h_i^{(1)}(\widehat h_i^{(1)})^\top
+
\frac1{n_2^2}\sum_{j=1}^{n_2}
\widehat h_j^{(2)}(\widehat h_j^{(2)})^\top
\right\}.
\]

The bootstrap critical value is
\begin{equation}
\label{eq:bootstrap-critical-value}
\widehat c_{1-\alpha}^{\mathrm{mb}}
=
\inf
\left\{
t\ge0:
\mathbb P_\xi
\left(
\widehat T_n^{\mathrm{mb}}\le t\mid\mathcal X
\right)
\ge 1-\alpha
\right\}.
\end{equation}
The level-\(\alpha\) test rejects \(H_0\) when
\begin{equation}
\label{eq:bootstrap-test}
\widehat T_n>\widehat c_{1-\alpha}^{\mathrm{mb}}.
\end{equation}
In implementation, the conditional quantile is approximated by Monte Carlo
simulation.  For \(b=1,\ldots,B\), generate independent Gaussian multipliers,
compute the bootstrap statistic in \eqref{eq:bootstrap-statistic}, and take the
empirical \((1-\alpha)\)-quantile of
\[
\widehat T_{n,1}^{\mathrm{mb}},\ldots,\widehat T_{n,B}^{\mathrm{mb}}.
\]

The bootstrap construction in \eqref{eq:bootstrap-vector} and
\eqref{eq:bootstrap-statistic} is conditionally Gaussian given the data.  Thus
its validity reduces to comparing the conditional Gaussian quadratic form with
the oracle Gaussian reference.  The next theorem gives a non-asymptotic bound
for this comparison.  Its proof uses a Gaussian comparison inequality for
squared Euclidean norms and a nuclear-norm consistency bound for \(\widehat R\).

Let \(C_R>0\) and \(C_{\mathrm{mb}}>0\) be constants depending only on the
constants in Assumptions~\ref{ass:coordinate_high_moment} and
\ref{ass:coordinate_nondegeneracy}.  For \(\eta_2\in(0,1)\), define
\begin{equation}
\label{eq:bootstrap-error}
\rho_{\mathrm{mb}}(\eta_2)
=
C_{\mathrm{mb}}\kappa(R)
\sqrt{\frac{d^3}{n_0\eta_2}},
\qquad
n_0=n_1\wedge n_2 .
\end{equation}

\begin{theorem}[Multiplier bootstrap approximation]
\label{thm:bootstrap_explicit}
Suppose Assumptions~\ref{ass:coordinate_high_moment},
\ref{ass:coordinate_nondegeneracy}, and
\ref{ass:gaussian_reference_geometry} hold.  If
\[
C_R\sqrt{\frac{d^3}{n_0\eta_2}}
\le
\frac12\Lambda_2,
\]
then, with probability at least \(1-\eta_2\),
\[
\sup_{r\ge0}
\left|
\mathbb P_\xi
\left(
\widehat T_n^{\mathrm{mb}}\le r\mid\mathcal X
\right)
-
\mathbb P(T_Z\le r)
\right|
\le
\rho_{\mathrm{mb}}(\eta_2).
\]
\end{theorem}

Combining Corollary~\ref{cor:feasible_gaussian} and
Theorem~\ref{thm:bootstrap_explicit} gives the validity of the bootstrap
critical value \eqref{eq:bootstrap-critical-value} and the rejection rule
\eqref{eq:bootstrap-test}.

\begin{corollary}[Size validity]
\label{cor:size_validity}
Suppose \(H_0:\Sigma_1=\Sigma_2\) holds.  Suppose
Assumptions~\ref{ass:coordinate_high_moment},
\ref{ass:coordinate_nondegeneracy}, and
\ref{ass:gaussian_reference_geometry} hold.  Assume further that the
conditions of Corollary~\ref{cor:feasible_gaussian} and
Theorem~\ref{thm:bootstrap_explicit} are satisfied along the sequence, and
that the error bounds in \eqref{eq:feasible-gaussian-error} and
\eqref{eq:bootstrap-error}, together with the corresponding bootstrap failure
probability $\rho_{\mathrm{mb}}(\eta_2)$, converge to zero.  Then the multiplier-bootstrap test
has asymptotically correct size:
\[
\mathbb P_{H_0}
\left(
\widehat T_n>\widehat c_{1-\alpha}^{\mathrm{mb}}
\right)
\to \alpha.
\]
\end{corollary}

\subsection{Power}
\label{subsec:power}

Although the theoretical analysis is carried out after suppressing the
superscript \((A)\), we restore it in this subsection to emphasize how the
choice of transformation affects the alternative signal.  Under
\(H_1:\Sigma_1\neq\Sigma_2\), define
\begin{equation}
\label{eq:transformed-signal}
\begin{aligned}
\Delta_\Sigma^{(A)}
&=
\operatorname{vech}\{A(\Sigma_1-\Sigma_2)A^\top\},\\
\theta_n^{(A)}
&=
\widetilde D^{-1/2}\Delta_\Sigma^{(A)}
\in\mathbb R^d .
\end{aligned}
\end{equation}
Here \(\theta_n^{(A)}\) in \eqref{eq:transformed-signal} is the standardized
transformed covariance difference.

The following condition describes a primitive class of alternatives.  Suppose
there exist an index set \(\mathcal I_n\subset\{1,\ldots,d\}\), with
\(|\mathcal I_n|=m_n\), and a sequence \(a_n>0\) such that
\[
\min_{k\in\mathcal I_n}
\left|
\theta_{n,k}^{(A)}
\right|
\ge a_n.
\]
Thus at least \(m_n\) transformed covariance coordinates have standardized
separation no smaller than \(a_n\).  The relevant signal accumulation is then
\begin{equation}
\label{eq:power-signal}
g_n
=
n_{\mathrm{eff}}\frac{m_n}{d}a_n^2.
\end{equation}

\begin{theorem}[Power consistency]
\label{thm:power_consistency}
Suppose \(H_1:\Sigma_1\neq\Sigma_2\) holds.  Suppose
Assumptions~\ref{ass:coordinate_high_moment} and
\ref{ass:coordinate_nondegeneracy} hold under the alternative.  Assume that
\[
\frac{d}{n_0}\to0,
\qquad
n_0=n_1\wedge n_2,
\]
and
\[
\sqrt{\frac{n_{\mathrm{eff}}}{d}}
\left(
\frac{p}{n_1}+\frac{p}{n_2}
\right)
=
o(\sqrt{g_n}).
\]
If
\[
g_n
=
n_{\mathrm{eff}}\frac{m_n}{d}a_n^2
\to\infty,
\]
then
\[
\mathbb P_{H_1}
\left(
\widehat T_n>\widehat c_{1-\alpha}^{\mathrm{mb}}
\right)
\to1.
\]
Equivalently, the corresponding right-tail multiplier-bootstrap \(p\)-value
converges to zero in probability.
\end{theorem}

The theorem shows that the test is consistent when the transformed covariance
difference is sufficiently aggregated across coordinates, as measured by
\eqref{eq:power-signal}.  In particular, dense alternatives with
\(m_n\asymp d\) require only \(n_{\mathrm{eff}}a_n^2\to\infty\), whereas very
sparse alternatives require a larger per-coordinate separation because the
\(\ell_2\)-statistic averages over all \(d=p(p+1)/2\) covariance coordinates.
The role of \(A\) is therefore to increase the number or the magnitude of
standardized transformed differences in \eqref{eq:transformed-signal} without
creating excessive estimation error.

\subsection{Unbalanced sample-size effect and simplified implementation}
\label{subsec:unbalanced_effect}

The simple quadratic structure of the proposed statistic gives a direct
interpretation in highly unbalanced designs.  The preceding power result shows
that the signal is scaled by the effective sample size \(n_{\mathrm{eff}}\) in
\eqref{eq:dimension-effective-n}.
When \(n_1\gg n_2\), \(n_{\mathrm{eff}}\) is asymptotically equivalent to
\(n_2\), and the contribution of the larger sample to the stochastic fluctuation
becomes negligible.  Consequently, the full two-sample statistic reduces to a
one-sample covariance discrepancy statistic for the smaller sample, with the
larger sample serving as a covariance reference.

Under \(H_0\), recall that
\[
\widetilde W=\sqrt{n_{\mathrm{eff}}}(\bar H_1-\bar H_2).
\]
When \(n_1\) is much larger than \(n_2\), define the oracle one-sample reference
vector
\begin{equation}
\label{eq:unbalanced-reference}
W_{2|1}=-\sqrt{n_2}\,\bar H_2,
\end{equation}
and the associated oracle quadratic statistic
\[
T_{2|1}^{\mathrm{or}}
=
\frac1d
\left\|
D_2^{-1/2}W_{2|1}
\right\|_2^2 .
\]

\begin{proposition}[Unbalanced-sample reduction]
\label{prop:unbalanced_reduction}
Suppose \(H_0:\Sigma_1=\Sigma_2\) holds and
Assumptions~\ref{ass:coordinate_high_moment} and
\ref{ass:coordinate_nondegeneracy} are satisfied.  If
\[
\varepsilon_n=\frac{n_2}{n_1}\to0,
\]
then
\[
\frac{n_{\mathrm{eff}}}{n_2}=1+O(\varepsilon_n),
\qquad
\max_{1\le k\le d}
\left|
\frac{\widetilde D_k}{D_{2,k}}-1
\right|
=
O(\varepsilon_n),
\]
and
\[
\frac1{\sqrt d}
\left\|
\widetilde D^{-1/2}\widetilde W
-
D_2^{-1/2}W_{2|1}
\right\|_2
=
O_{\mathbb P}(\varepsilon_n^{1/2}).
\]
Consequently,
\[
\widetilde T_n-T_{2|1}^{\mathrm{or}}
=
O_{\mathbb P}(\varepsilon_n^{1/2}).
\]
\end{proposition}

Proposition~\ref{prop:unbalanced_reduction} shows that, when
\(n_1/n_2\to\infty\), the larger sample provides an asymptotically fixed
covariance reference, while the stochastic fluctuation and the critical region
are governed by the smaller-sample reference vector
\eqref{eq:unbalanced-reference}.  The rate
\(O_{\mathbb P}\{(n_2/n_1)^{1/2}\}\) quantifies how quickly this one-sample
reduction becomes accurate.  The symmetric conclusion holds when
\(n_2/n_1\to\infty\), with the roles of the two samples interchanged.

The preceding result also suggests a simplified implementation in highly
unbalanced settings.  Suppose \(n_1\gg n_2\).  Let \(\widehat\Sigma_1\) be the
covariance estimator from the larger sample and \(\widehat\Sigma_2\) the
covariance estimator from the smaller sample.  Define
\[
\widehat W_{2|1}
=
\sqrt{n_2}\,
\operatorname{vech}(\widehat\Sigma_1-\widehat\Sigma_2).
\]
For \(k=k(a,b)\), \(1\le a\le b\le p\), define the smaller-sample diagonal
studentizer
\[
\widehat D_{2,kk}
=
\frac1{n_2}
\sum_{j=1}^{n_2}
\bigl(\widehat h_{j,k}^{(2)}\bigr)^2,
\qquad
1\le k\le d,
\]
and set
\begin{equation}
\label{eq:unbalanced-statistic}
\widehat T_{2|1}
=
\frac1d
\left\|
\widehat D_2^{-1/2}\widehat W_{2|1}
\right\|_2^2 .
\end{equation}
The statistic in \eqref{eq:unbalanced-statistic} treats the larger sample as a
covariance reference and uses only the smaller sample to estimate the
stochastic fluctuation.

The ball geometry makes this reduction especially transparent.  For a critical
value \(c\), the simplified acceptance region can be written as
\begin{equation}
\label{eq:confidence-ball-interpretation}
\left\{
\widehat T_{2|1}\le c
\right\}
=
\left\{
\left\|
\widehat D_2^{-1/2}
\operatorname{vech}(\widehat\Sigma_1-\widehat\Sigma_2)
\right\|_2
\le
\sqrt{\frac{dc}{n_2}}
\right\}.
\end{equation}
Thus, in the studentized covariance-coordinate metric, the test accepts exactly
when the larger-sample covariance estimator \(\widehat\Sigma_1\) lies inside
the one-sample confidence ball in \eqref{eq:confidence-ball-interpretation},
centered at the smaller-sample estimator \(\widehat\Sigma_2\), with radius
determined by the one-sample bootstrap critical value.

\begin{corollary}[Feasible simplified statistic]
\label{cor:feasible_unbalanced}
Suppose \(H_0:\Sigma_1=\Sigma_2\) holds and
Assumptions~\ref{ass:coordinate_high_moment} and
\ref{ass:coordinate_nondegeneracy} are satisfied.  If
\[
\varepsilon_n=\frac{n_2}{n_1}\to0,
\qquad
\frac d{n_2}\to0,
\]
then
\[
\widehat T_{2|1}
-
T_{2|1}^{\mathrm{or}}
=
O_{\mathbb P}
\left(
\varepsilon_n^{1/2}
+
\sqrt{\frac d{n_2}}
\right).
\]
Moreover, if the full sample-to-oracle transfer satisfies
\[
\widehat T_n-\widetilde T_n=o_{\mathbb P}(1),
\]
then
\[
\widehat T_n-\widehat T_{2|1}=o_{\mathbb P}(1).
\]
More explicitly,
\[
|\widehat T_n-\widehat T_{2|1}|
\le
|\widehat T_n-\widetilde T_n|
+
O_{\mathbb P}
\left(
\varepsilon_n^{1/2}
+
\sqrt{\frac d{n_2}}
\right).
\]
\end{corollary}

The simplified statistic can be calibrated by a one-sample multiplier bootstrap.
Generate independent \(N(0,1)\) multipliers \(e_1',\ldots,e_{n_2}'\) and define
\[
\widehat W_{2|1}^{\mathrm{mb}}
=
\frac1{\sqrt{n_2}}
\sum_{j=1}^{n_2}
e_j'\widehat h_j^{(2)},
\]
with
\[
\widehat T_{2|1}^{\mathrm{mb}}
=
\frac1d
\left\|
\widehat D_2^{-1/2}\widehat W_{2|1}^{\mathrm{mb}}
\right\|_2^2 .
\]
The simplified bootstrap critical value is the conditional
\((1-\alpha)\)-quantile of \(\widehat T_{2|1}^{\mathrm{mb}}\).  This simplified
procedure should be viewed as an approximation to the full two-sample procedure
in highly unbalanced regimes.  The full statistic remains the default
implementation, while \(\widehat T_{2|1}\) provides a computationally and
conceptually transparent approximation when \(n_2/n_1\) is small.


\section{Simulation Studies}\label{sec:simulation}
The simulation studies are designed to evaluate the finite-sample behavior of
the proposed procedure in regimes where the covariance-coordinate dimension
\(d=p(p+1)/2\) is much larger than the sample sizes.  The numerical study is
organized around two questions.  Simulation 1 studies the proposed procedure
itself: we compare the identity implementation \(A=I_p\) with target
transformations \(A=G(\rho_{\rm tar})\) for
\(\rho_{\rm tar}\in\{0.3,0.5,0.7\}\).  This part asks whether the target
transformation improves finite-sample behavior and whether the procedure is
sensitive to the exact target correlation.  Simulation 2 then fixes the
representative target \(G(0.5)\) and compares it with existing covariance tests,
both under the same null designs and under power alternatives.  This separation
is important: the first simulation justifies the transformation and the choice
of a default target, while the second simulation evaluates the default method
against external benchmarks.

The proposed statistic can be implemented in the original coordinate system,
denoted by \(A=I_p\), or after a transformation \(A=G\).  In applications, a
transformation can be used to express the covariance difference in a geometry
where dense covariance-coordinate signals are more visible.  We implement this
idea by mapping the empirical pooled covariance to a dense equicorrelation
target.  Specifically, let
\[
\widehat\Sigma_g
=
\frac1{n_g}\sum_{i=1}^{n_g}
(X_{gi}-\bar X_g)(X_{gi}-\bar X_g)^\top,
\qquad g=1,2,
\]
and
\[
\widehat\Sigma_{\rm pool}
=
\frac{n_1\widehat\Sigma_1+n_2\widehat\Sigma_2}{n_1+n_2}.
\]
For a target correlation \(\rho_{\rm tar}\in\{0.3,0.5,0.7\}\), define the
dense equicorrelation target
\begin{equation}
\label{eq:simulation-target}
\Sigma_{\rm tar}(\rho_{\rm tar})
=
(1-\rho_{\rm tar})I_p+\rho_{\rm tar}{\bf 1}_p{\bf 1}_p^\top.
\end{equation}
The numerical transformation used in the simulations is
\begin{equation}
\label{eq:simulation-G}
G(\rho_{\rm tar})
=
\Sigma_{\rm tar}(\rho_{\rm tar})^{1/2}
\left(\widehat\Sigma_{\rm pool}+\lambda I_p\right)^{-1/2},
\qquad
\lambda
=
10^{-6}\frac{\operatorname{tr}(\widehat\Sigma_{\rm pool})}{p}.
\end{equation}
The ridge term in \eqref{eq:simulation-G} is used only to compute the inverse
square root stably.  The theoretical results are stated for a deterministic
transformation \(A\); \(G(\rho_{\rm tar})\) in \eqref{eq:simulation-G} is the
empirical implementation used to study how target geometry affects finite-sample
calibration and power.

The benchmark methods in Simulation 2 are a maximum-coordinate multiplier
bootstrap test, the normal-reference method of \citet{wang_two-sample_2024},
and the two-sample covariance test of \citet{li_two_2012}.  All simulations use
independent samples with sample covariance matrices computed with denominator
\(n\), as in the proposed statistic, and the nominal level is
\(\alpha=0.05\).  The number of Monte Carlo repetitions \(R\) and
multiplier-bootstrap draws \(B\) is stated in each simulation block.  In the
variance-heterogeneous power experiment, we also report a raw \(\ell_2\)
bootstrap statistic, defined by removing the diagonal studentizer
\(\widehat D^{-1/2}\) from \eqref{eq:proposed-statistic} while keeping the same
quadratic aggregation and bootstrap calibration.  This additional diagnostic
separates the effect of coordinate studentization from the effect of target
transformation.

We use the following covariance notation throughout the simulation section.
For \(\rho\in[0,1)\), let
\[
\Sigma_{\rm Eq}(\rho)=(1-\rho)I_p+\rho{\bf 1}_p{\bf 1}_p^\top
\]
denote the equicorrelation covariance matrix, and let
\[
\{\Sigma_{\rm Toe}(\rho)\}_{ab}=\rho^{|a-b|},\qquad 1\le a,b\le p,
\]
denote the Toeplitz covariance matrix.  In the tables, Eq and Toe abbreviate
\(\Sigma_{\rm Eq}(\rho_0)\) and \(\Sigma_{\rm Toe}(\rho_0)\), respectively.

\subsection{Simulation 1: target transformation and target sensitivity}
\label{sec:simulation_target}

We first examine calibration under the null hypothesis \(\Sigma_1=\Sigma_2\)
for the proposed family of tests.  The dimensions are \(p\in\{100,200\}\), and
the sample sizes are
\((n_1,n_2)\in\{(80,100),(120,150)\}\), with an additional small-sample check
at \(n_1=n_2=60\).  For each dimension and sample-size pair, the common
covariance is either \(\Sigma_{\rm Eq}(\rho_0)\) or
\(\Sigma_{\rm Toe}(\rho_0)\), with
\(\rho_0\in\{0.1,0.5,0.9\}\).  The innovations are Gaussian, standardized
centered \(\chi^2_1\), and standardized Laplace.  All formal size simulations
use \(R=1000\) Monte Carlo repetitions and \(B=500\) multiplier draws for
bootstrap methods.  Tables~\ref{tab:size_main_p100} and
\ref{tab:size_main_p200} report the empirical rejection probabilities for the
three target-transformed tests \(G(0.3)\), \(G(0.5)\), and \(G(0.7)\), together
with the identity implementation \(I_p\).  To keep the tables compact, each
slash-separated triplet is ordered as
\(\rho_0=0.1,0.5,0.9\); Gau, Chi, and Lap denote Gaussian, standardized
centered \(\chi^2_1\), and standardized Laplace innovations, while Eq and Toe
denote the two covariance families defined above.

\begin{table}[!htbp]
\centering
\begingroup
\spacingset{1}
\scriptsize
\captionsetup{font=footnotesize,skip=2pt}
\renewcommand{\arraystretch}{0.78}
\setlength{\tabcolsep}{1.8pt}
\setlength{\aboverulesep}{0.15ex}
\setlength{\belowrulesep}{0.15ex}
\caption{Empirical size of the proposed target-transformed and identity tests for $p=100$. Each entry reports the $\rho_0=0.1/0.5/0.9$ triplet; abbreviations are defined in the text.}\label{tab:size_main_p100}
\resizebox{\textwidth}{!}{%
\begin{tabular}{@{}lllcccc@{}}
\toprule
$(n_1,n_2)$ & Dist. & Cov. & $G(.3)$ & $G(.5)$ & $G(.7)$ & $I_p$ \\
\midrule
$(60,60)$ & Gau & Eq & 0.046/0.049/0.057 & 0.050/0.065/0.062 & 0.051/0.062/0.069 & 0.001/0.054/0.060 \\
 &  & Toe & 0.050/0.053/0.055 & 0.061/0.057/0.056 & 0.061/0.062/0.061 & 0.000/0.000/0.034 \\
 & Chi & Eq & 0.061/0.047/0.050 & 0.062/0.046/0.054 & 0.064/0.043/0.045 & 0.005/0.045/0.049 \\
 &  & Toe & 0.058/0.042/0.047 & 0.061/0.054/0.054 & 0.063/0.051/0.053 & 0.000/0.000/0.030 \\
 & Lap & Eq & 0.048/0.053/0.040 & 0.062/0.060/0.044 & 0.058/0.060/0.044 & 0.003/0.062/0.047 \\
 &  & Toe & 0.042/0.058/0.047 & 0.050/0.061/0.051 & 0.050/0.062/0.050 & 0.000/0.000/0.023 \\
\addlinespace[0.5pt]
$(80,100)$ & Gau & Eq & 0.046/0.054/0.044 & 0.050/0.054/0.053 & 0.047/0.052/0.049 & 0.005/0.050/0.054 \\
 &  & Toe & 0.059/0.054/0.056 & 0.063/0.062/0.063 & 0.065/0.059/0.056 & 0.000/0.000/0.039 \\
 & Chi & Eq & 0.058/0.049/0.059 & 0.058/0.045/0.062 & 0.062/0.054/0.061 & 0.010/0.048/0.061 \\
 &  & Toe & 0.060/0.053/0.050 & 0.061/0.058/0.055 & 0.063/0.055/0.058 & 0.000/0.000/0.027 \\
 & Lap & Eq & 0.062/0.056/0.054 & 0.066/0.060/0.057 & 0.063/0.057/0.056 & 0.008/0.054/0.058 \\
 &  & Toe & 0.047/0.059/0.042 & 0.051/0.067/0.052 & 0.050/0.062/0.052 & 0.000/0.000/0.042 \\
\addlinespace[0.5pt]
$(120,150)$ & Gau & Eq & 0.058/0.060/0.051 & 0.062/0.066/0.054 & 0.057/0.070/0.053 & 0.011/0.073/0.049 \\
 &  & Toe & 0.060/0.056/0.057 & 0.064/0.057/0.053 & 0.061/0.052/0.048 & 0.000/0.000/0.053 \\
 & Chi & Eq & 0.064/0.040/0.050 & 0.057/0.040/0.044 & 0.064/0.039/0.048 & 0.016/0.040/0.040 \\
 &  & Toe & 0.056/0.044/0.052 & 0.053/0.049/0.058 & 0.053/0.046/0.056 & 0.000/0.000/0.031 \\
 & Lap & Eq & 0.052/0.047/0.052 & 0.056/0.051/0.054 & 0.065/0.049/0.052 & 0.008/0.052/0.054 \\
 &  & Toe & 0.042/0.058/0.053 & 0.044/0.062/0.051 & 0.044/0.061/0.053 & 0.000/0.000/0.036 \\
\bottomrule
\end{tabular}%
}
\endgroup
\end{table}


Although the ambient dimensions \(p=100\) and \(p=200\) appear moderate, the
testing problem is high-dimensional because the statistic aggregates over
the covariance-coordinate dimension \(d=p(p+1)/2\).  Thus \(d=5050\) when
\(p=100\) and \(d=20100\) when \(p=200\), both far exceeding the sample sizes
used in the simulation, including the small-sample setting \(n_1=n_2=60\)
that satisfies the prespecified calibration check and is therefore included in
the main table.  This regime is precisely where direct estimation of
the full covariance matrix of covariance entries is unstable, and where the
proposed diagonal studentization and multiplier-bootstrap calibration are
intended to be useful.

\begin{table}[!htbp]
\centering
\begingroup
\spacingset{1}
\scriptsize
\captionsetup{font=footnotesize,skip=2pt}
\renewcommand{\arraystretch}{0.78}
\setlength{\tabcolsep}{1.8pt}
\setlength{\aboverulesep}{0.15ex}
\setlength{\belowrulesep}{0.15ex}
\caption{Empirical size of the proposed target-transformed and identity tests for $p=200$. Each entry reports the $\rho_0=0.1/0.5/0.9$ triplet; abbreviations are defined in the text.}\label{tab:size_main_p200}
\resizebox{\textwidth}{!}{%
\begin{tabular}{@{}lllcccc@{}}
\toprule
$(n_1,n_2)$ & Dist. & Cov. & $G(.3)$ & $G(.5)$ & $G(.7)$ & $I_p$ \\
\midrule
$(60,60)$ & Gau & Eq & 0.050/0.062/0.042 & 0.051/0.068/0.044 & 0.056/0.066/0.041 & 0.001/0.063/0.043 \\
 &  & Toe & 0.052/0.042/0.054 & 0.056/0.046/0.056 & 0.061/0.046/0.057 & 0.000/0.000/0.011 \\
 & Chi & Eq & 0.042/0.056/0.047 & 0.049/0.056/0.046 & 0.045/0.054/0.044 & 0.008/0.059/0.044 \\
 &  & Toe & 0.045/0.059/0.043 & 0.045/0.059/0.044 & 0.049/0.055/0.047 & 0.000/0.000/0.005 \\
 & Lap & Eq & 0.053/0.046/0.051 & 0.056/0.046/0.052 & 0.055/0.046/0.051 & 0.001/0.045/0.049 \\
 &  & Toe & 0.054/0.060/0.059 & 0.052/0.056/0.058 & 0.056/0.060/0.054 & 0.000/0.000/0.007 \\
\addlinespace[0.5pt]
$(80,100)$ & Gau & Eq & 0.059/0.053/0.057 & 0.060/0.050/0.061 & 0.060/0.053/0.056 & 0.004/0.053/0.058 \\
 &  & Toe & 0.053/0.050/0.052 & 0.053/0.049/0.054 & 0.056/0.051/0.056 & 0.000/0.000/0.015 \\
 & Chi & Eq & 0.050/0.069/0.052 & 0.053/0.064/0.053 & 0.048/0.061/0.057 & 0.006/0.073/0.053 \\
 &  & Toe & 0.044/0.054/0.056 & 0.048/0.063/0.057 & 0.048/0.056/0.056 & 0.000/0.000/0.011 \\
 & Lap & Eq & 0.061/0.061/0.051 & 0.067/0.062/0.055 & 0.066/0.072/0.053 & 0.005/0.065/0.058 \\
 &  & Toe & 0.047/0.049/0.053 & 0.052/0.048/0.049 & 0.056/0.051/0.057 & 0.000/0.000/0.016 \\
\addlinespace[0.5pt]
$(120,150)$ & Gau & Eq & 0.049/0.054/0.058 & 0.053/0.053/0.060 & 0.053/0.058/0.056 & 0.005/0.051/0.061 \\
 &  & Toe & 0.066/0.051/0.046 & 0.070/0.057/0.056 & 0.066/0.057/0.054 & 0.000/0.000/0.024 \\
 & Chi & Eq & 0.049/0.054/0.043 & 0.053/0.052/0.047 & 0.046/0.049/0.050 & 0.013/0.046/0.044 \\
 &  & Toe & 0.070/0.066/0.061 & 0.075/0.068/0.061 & 0.070/0.066/0.061 & 0.000/0.000/0.020 \\
 & Lap & Eq & 0.044/0.050/0.039 & 0.046/0.051/0.042 & 0.052/0.050/0.042 & 0.006/0.048/0.041 \\
 &  & Toe & 0.050/0.048/0.049 & 0.054/0.053/0.052 & 0.051/0.052/0.052 & 0.000/0.000/0.025 \\
\bottomrule
\end{tabular}%
}
\endgroup
\end{table}


The main conclusion from the size study is that the target transformation
stabilizes the proposed test.  Across the 108 main null settings, the mean
empirical sizes of \(G(0.3)\), \(G(0.5)\), and \(G(0.7)\) are 0.052, 0.055,
and 0.055, respectively, with all setting-level rejection rates lying between
0.039 and 0.075.  Thus the transformed tests remain close to the nominal level
over changes in \(p\), sample size, covariance family, dependence strength, and
innovation distribution.  The three target choices also produce very similar
calibration patterns, indicating that the transfer to a dense target geometry
is not especially sensitive to the precise value of
\(\rho_{\rm tar}\in\{0.3,0.5,0.7\}\).  In contrast, the identity implementation
has mean empirical size 0.023 and is often conservative, especially in weakly
dependent Toeplitz designs.  This comparison shows that the gain from
\(G(\rho_{\rm tar})\) is not a fragile tuning effect; it comes from
re-expressing the covariance difference in a geometry that is better aligned
with dense \(\ell_2\) aggregation.

We next examine power within the proposed family under dense equicorrelation
alternatives.  The power studies use \(p=200\), \(n_1=120\), \(n_2=150\),
Gaussian samples, nominal level \(0.05\), \(R=1000\) Monte Carlo repetitions,
and \(B=500\) multiplier draws for bootstrap methods.  The first experiment
sets \(\Sigma_1=\Sigma_{\rm Eq}(0)\) and
\(\Sigma_2=\Sigma_{\rm Eq}(\rho_1)\), where \(\rho_1\) increases from zero.
The second sets \(\Sigma_1=\Sigma_{\rm Eq}(0.5)\) and
\(\Sigma_2=\Sigma_{\rm Eq}(\rho_1)\), with
\(\rho_1\in\{0.1,\ldots,0.9\}\).  In both experiments, \(G(0.3)\),
\(G(0.5)\), and \(G(0.7)\) have nearly overlapping power curves.  The maximum
setting-level difference among the three target-transformed rejection rates is
0.011 in the \(\rho_0=0\) experiment and 0.021 in the \(\rho_0=0.5\)
experiment.  Relative to the identity implementation, the transformation is
especially beneficial near the identity baseline: when \(\rho_0=0\), \(I_p\)
has essentially no rejection before \(\rho_1=0.03\), whereas \(G(0.5)\) already
has rejection rates 0.387 and 0.922 at \(\rho_1=0.002\) and
\(\rho_1=0.005\), respectively.  We therefore use \(G(0.5)\) as the default
target in the remaining comparisons: it has stable size, strong power, and is
representative of the target-transformed family rather than a specially tuned
choice.

\subsection{Simulation 2: default target versus benchmark tests}
\label{sec:simulation_benchmark}

We now fix the proposed method at \(G(0.5)\) and compare it with the existing
tests.  Under the same null grid used in Simulation 1, the max-bootstrap and
Wang-NR tests are often conservative, with mean empirical sizes 0.042 and
0.025, respectively, while Li--Chen is liberal in several correlated
non-Gaussian settings, with mean empirical size 0.105 and maximum empirical
size 0.269.  These size results provide the calibration context for the power
comparisons below: rejection rates away from the null are most meaningful when
the corresponding method is also well calibrated near the null.

The first benchmark power comparison uses the two equicorrelation designs
described above.  The results are reported graphically because the scientific
question is the shape of the power curves as the alternative strengthens,
rather than isolated rejection rates at a single point.

\begin{figure}[!htbp]
\centering
\begin{subfigure}{0.49\textwidth}
\centering
\includegraphics[width=\textwidth]{figures/study_power_target_eq0_final_power.png}
\caption{\(\rho_0=0\)}
\label{fig:power_eq0}
\end{subfigure}\hfill
\begin{subfigure}{0.49\textwidth}
\centering
\includegraphics[width=\textwidth]{figures/study_power_target_eq05_final_power.png}
\caption{\(\rho_0=0.5\)}
\label{fig:power_eq05}
\end{subfigure}
\caption{Empirical rejection rates for dense equicorrelation alternatives.
Panel (a) compares \(\Sigma_{\rm Eq}(0)\) with
\(\Sigma_{\rm Eq}(\rho_1)\).  Panel (b) compares
\(\Sigma_{\rm Eq}(0.5)\) with \(\Sigma_{\rm Eq}(\rho_1)\).  The proposed method
uses the default target \(G(0.5)\) and is compared with max-bootstrap, Wang-NR,
and Li--Chen.  Results are based on \(R=1000\) Monte Carlo repetitions and
\(B=500\) multiplier draws for bootstrap methods.}
\label{fig:eqcor_power}
\end{figure}

Figure~\ref{fig:eqcor_power} shows that \(G(0.5)\) is highly sensitive to dense
covariance alternatives while preserving the calibration pattern observed under
the null.  When \(\rho_0=0\), even very small positive equicorrelations produce
a covariance-coordinate signal spread over many entries: \(G(0.5)\) has power
0.387 at \(\rho_1=0.002\), 0.922 at \(\rho_1=0.005\), and reaches 1.000 by
\(\rho_1=0.010\).  Over the same range, max-bootstrap remains near the nominal
level, and Li--Chen increases more slowly.  When \(\rho_0=0.5\), the
\(G(0.5)\) curve is lowest at the null point \(\rho_1=0.5\), with rejection
rate 0.053, and increases as \(\rho_1\) moves away from the correlated
baseline.  The benchmark methods also gain power for large departures, but the
default target-transformed test is especially competitive for weak and moderate
dense alternatives while maintaining the size behavior established in
Simulation 1.

The second benchmark power comparison is designed to isolate the statistical
role of studentization.  Both groups have independent coordinates under the
null.  The first 80 coordinates form a low-variance block with marginal
variance one, and the remaining 120 coordinates are high-variance nuisance
coordinates with marginal variance 25.  Under the alternative, only the low-variance block
changes: its within-block correlation is increased from zero to
\(\rho_{\rm block}\), while all other correlations and all marginal variances
remain unchanged.  Thus \(\rho_{\rm block}=0\) is the null case.  This
construction creates a dense covariance signal in coordinates whose raw
Frobenius contribution can be masked by high-variance nuisance coordinates.

\begin{figure}[!htbp]
\centering
\includegraphics[width=0.86\textwidth]{figures/study_power_studentized_targets_final_power.png}
\caption{Empirical rejection rates for the variance-heterogeneous low-variance
block alternative.  The first 80 coordinates carry the dense block-correlation
signal and have marginal variance one; the remaining 120 coordinates have
marginal variance 25 and act as high-variance nuisance coordinates.  Results
are based on \(R=1000\) Monte Carlo repetitions and \(B=500\) multiplier draws
for bootstrap methods.}
\label{fig:studentized_target_power}
\end{figure}

Figure~\ref{fig:studentized_target_power} demonstrates that studentization is
not merely a technical normalization.  The raw \(\ell_2\) statistic, which
uses the unstudentized quadratic norm of the covariance-coordinate difference,
has essentially no power in this design because the high-variance nuisance
coordinates dominate the raw scale.  By contrast, \(G(0.5)\) is close to
nominal size at \(\rho_{\rm block}=0\) and becomes powerful for weak
within-block correlations: its rejection rate is 0.589 at
\(\rho_{\rm block}=0.05\), 0.804 at \(\rho_{\rm block}=0.10\), and 0.903 at
\(\rho_{\rm block}=0.15\).  The max-bootstrap test becomes powerful only at
larger block correlations, while Wang-NR and Li--Chen have little power in this
particular variance-heterogeneous design.  This comparison shows that the
default target-transformed test is not only calibrated in the main size study
but also effective against dense covariance signals that can be hidden by
heterogeneous marginal scales.

\section{Real Data Application}\label{sec:realdata}

\subsection{TCGA-BRCA covariance analysis}

We apply the proposed covariance-testing framework to TCGA-BRCA RNA-seq data, comparing Basal-like and Luminal~A breast-cancer subtypes.  This comparison is biologically well motivated because TCGA molecular profiling shows that these breast-cancer subtypes have distinct molecular and transcriptional profiles \citep{tcga_brca_molecular_portraits_2012}.  In contrast to conventional analyses that focus primarily on marginal expression differences, our objective is to examine whether a biologically coherent transcriptional program exhibits subtype-specific covariance organization.

Expression and clinical annotation files were obtained from the UCSC Xena TCGA-BRCA cohort \citep{goldman_visualizing_2020}, using the \texttt{HiSeqV2.gz} expression matrix and \texttt{BRCA\_clinicalMatrix} clinical annotation file.  We retained primary tumor samples with PAM50 RNA-seq subtype labels Basal-like or Luminal~A.  The resulting groups contain \(n_1=141\) Basal-like tumors and \(n_2=421\) Luminal~A tumors.  Expression values are analyzed on the Xena log-transformed RNA-seq scale.  Genes are matched by symbol to the MSigDB Hallmark file \texttt{h.all.v2025.1.Hs.symbols.gmt} \citep{liberzon_molecular_2015}; duplicate gene symbols are averaged, and genes with negligible variance are removed.

The application focuses on a Hallmark cell-cycle core panel formed by the union of the complete Hallmark E2F targets and G2M checkpoint gene sets.  These two Hallmark components represent closely related transcriptional programs governing cell-cycle entry, DNA replication, and mitotic progression, and are therefore directly relevant to the proliferative phenotype that distinguishes Basal-like tumors from Luminal~A tumors.  After matching to the expression matrix and applying the same low-variance filtering described above, the panel contains \(p=304\) genes.  Since \(p\ge \min(n_1,n_2)\), this application remains in the high-dimensional regime targeted by the proposed methodology.

We compare the proposed bootstrap statistic with \(A=I_p\), the target-transformed implementation with \(A=G\), max-bootstrap, Wang-NR, and Li--Chen.  Bootstrap methods use \(B=2000\) multiplier draws.  The table reports program-level \(p\)-values for this focused biological application; these values are interpreted as evidence for covariance reorganization within the defined cell-cycle program rather than as a genome-wide discovery claim.

\begin{table}[!htbp]
\centering
\footnotesize
\caption{TCGA-BRCA Basal-like versus Luminal~A covariance comparison for the Hallmark cell-cycle core panel.  The table reports the observed test statistic, program-level \(p\)-value, and rejection decision.  For bootstrap methods, \(<0.0005\) means that no bootstrap draw exceeded the observed statistic among \(B=2000\) draws.}
\label{tab:tcga_realdata_programs}
\begin{tabular}{@{}lrlc@{}}
\toprule
method & statistic & \(p\)-value & reject \\
\midrule
\(A=I_p\) & 4.239 & \(<0.0005\) & yes \\
\(A=G\) & 7.232 & \(<0.0005\) & yes \\
Max & 11.980 & \(<0.0005\) & yes \\
Wang-NR & 1777.476 & \(5.1\times 10^{-13}\) & yes \\
Li--Chen & 27.004 & \(6.6\times 10^{-161}\) & yes \\
\bottomrule
\end{tabular}
\end{table}

Table~\ref{tab:tcga_realdata_programs} shows strong evidence of covariance reorganization in the Hallmark cell-cycle core panel.  Both proposed implementations reject the null hypothesis: the untransformed statistic with \(A=I_p\) and the target-transformed implementation with \(A=G\) each yield bootstrap \(p<0.0005\).  The same conclusion is supported by the benchmark procedures, with max-bootstrap giving \(p<0.0005\), Wang-NR giving \(p=5.1\times 10^{-13}\), and Li--Chen giving \(p=6.6\times 10^{-161}\).  Thus, the empirical finding is not an artifact of a particular calibration scheme.  The agreement between \(A=I_p\) and \(A=G\) is also consistent with the dense program-level covariance signal emphasized in the simulation studies.

The selected panel has a direct biological interpretation.  It contains 187 matched E2F-target genes and 183 matched G2M-checkpoint genes, with 304 unique genes after accounting for overlap between the two Hallmark components.  These components capture coordinated transcriptional activity associated with DNA replication, cell-cycle progression, and mitotic control, all of which are central to the proliferative biology of Basal-like breast tumors relative to Luminal~A tumors.  As a descriptive check on the source of the signal, the mean absolute difference between the two subtype-specific sample correlations is 0.145 within the E2F-target component, 0.137 within the G2M-checkpoint component, and 0.141 across the two components.  The covariance-level rejection therefore has a natural program-level interpretation: the two subtypes differ not only in the expression levels of proliferation-related genes, but also in the pattern of gene-gene coordination within this cell-cycle program.

\subsection{Real-data-calibrated simulation}

As an additional calibration check tied directly to this application, we conducted a real-data-calibrated Gaussian simulation.  The sample sizes and dimension were fixed at the observed values \(n_1=141\), \(n_2=421\), and \(p=304\).  For the size experiment, both samples were generated from a common covariance matrix equal to the pooled empirical covariance from the selected panel.  For the power experiment, the two samples were generated from the subtype-specific empirical covariance matrices estimated from the Basal-like and Luminal~A groups, respectively.  This design preserves the scale and correlation geometry of the application while making the null and alternative data-generating mechanisms explicit.

\begin{table}[!htbp]
\centering
\caption{Real-data-calibrated Gaussian simulation for the TCGA-BRCA cell-cycle core panel.  Size uses the pooled empirical covariance as the common covariance matrix; power uses the two subtype-specific empirical covariance matrices.  Results are based on \(R=1000\) Monte Carlo repetitions and \(B=2000\) bootstrap draws for bootstrap methods.}
\label{tab:tcga_realdata_calibrated_sim}
\begin{tabular}{lrr}
\toprule
method & size & power \\
\midrule
\(A=I_p\) & 0.071 & 1.000 \\
\(A=G\) & 0.065 & 1.000 \\
Max-bootstrap & 0.134 & 1.000 \\
Wang-NR & 0.062 & 1.000 \\
Li--Chen & 0.112 & 1.000 \\
\bottomrule
\end{tabular}
\end{table}

Table~\ref{tab:tcga_realdata_calibrated_sim} supports the interpretation of the application.  Under the pooled-covariance null, the two proposed bootstrap tests have empirical rejection rates 0.071 for \(A=I_p\) and 0.065 for \(A=G\), close to the nominal level.  Wang-NR is similar, whereas max-bootstrap and Li--Chen are more liberal in this high-dimensional empirical-covariance design.  Under the subtype-specific covariance alternative, all methods reject in every repetition, indicating that the covariance contrast estimated from the real data is strong under Gaussian resampling.
\section{Conclusion and Discussion}\label{sec:discussion}

This paper develops a studentized \(\ell_2\) approach to two-sample covariance
testing.  The method compares covariance matrices through standardized
upper-triangular covariance-coordinate differences and aggregates them by a
squared Euclidean norm.  This construction gives a transparent scale-normalized
global statistic, avoids estimating the full covariance matrix of covariance
entries, and leads naturally to centered-ball Gaussian approximation.  We prove
oracle Gaussian approximation, feasible transfer from oracle to sample-centered
studentized statistics, multiplier-bootstrap validity, and power consistency
under accumulated standardized covariance signal.  The same quadratic structure
also yields a direct simplification in highly unbalanced designs, where the
larger sample acts as a covariance reference and the smaller sample governs the
stochastic fluctuation.

The simulations support the intended use of the procedure.  The proposed
bootstrap tests have stable empirical size across the null settings considered,
and the target-transformed implementation \(A=G\) is particularly effective for
dense equicorrelation alternatives when compared among methods with acceptable
size behavior.  The untransformed statistic \(A=I_p\) remains a useful baseline,
while max-type procedures are less sensitive to diffuse covariance changes and
some quadratic-form benchmarks can exhibit size distortion in correlated
designs.  In the TCGA-BRCA application, both proposed implementations detect a
clear covariance reorganization in the Hallmark cell-cycle core panel between
Basal-like and Luminal~A tumors, and the real-data-calibrated simulation
supports the interpretation that the observed covariance contrast is strong
rather than a calibration artifact.

The theory is deliberately explicit, but the resulting high-dimensional
sufficient conditions are conservative.  In particular, the proof uses
uniform covariance-coordinate controls and Gaussian comparison bounds that are
not expected to be sharp for all covariance geometries.  The simulation results
suggest that the bootstrap statistic can perform well in regimes beyond those
directly captured by the current growth conditions.  A natural next step is to
replace worst-case dimension dependence with sharper effective-rank,
spectral-decay, or geometry-adaptive quantities, which may better reflect the
behavior observed in finite samples.

Another important direction is the choice of the transformation \(A\).  The
power analysis shows that \(A\) affects the standardized signal by changing how
the covariance difference is represented before aggregation.  This raises a
broader question: whether one can characterize, estimate, or approximate
transformations that are globally more powerful for classes of alternatives
while retaining reliable bootstrap calibration.  The dense equicorrelation
target used here is one practical default, but it is unlikely to be optimal in
all scientific settings.  Developing principled rules for selecting \(A\), and
understanding the tradeoff between signal enhancement and estimation stability,
would make the framework more adaptive and potentially more powerful.

\section{Disclosure statement}\label{disclosure-statement}

A disclosure statement should be completed before submission.

\section{Data Availability Statement}\label{data-availability-statement}

The simulation and real-data analysis code are available at
\url{https://github.com/shuhaozhen1/covtestsim}.  The TCGA-BRCA expression and
clinical annotations used in the application are publicly available from UCSC
Xena, and the Hallmark gene-set definitions are publicly available from MSigDB.
The real-data pipeline in the code repository downloads these public resources
and reproduces the processed analysis files.

\phantomsection\label{supplementary-material}
\bigskip

\begin{center}
{\large\bf SUPPLEMENTARY MATERIAL}
\end{center}

\begin{description}
\item[Technical supplement:]
Auxiliary lemmas and proofs are included below in the appendix of this
combined manuscript.
\item[Code repository:]
The simulation and real-data analysis code are organized as a standalone code
repository, with source code, command-line runners, tests, and a README for
reproducing the empirical results, at
\url{https://github.com/shuhaozhen1/covtestsim}.
\end{description}

\clearpage
\appendix
\spacingset{1.8}

\begin{center}
{\LARGE\bf Supplementary Material}\\[0.5em]
{\large Simple Studentized \texorpdfstring{$\ell_2$}{l2} Statistics and Gaussian Bootstrap for Two-Sample Covariance Testing}\\[0.5em]
Haozhen Shu and Tianming Zhu
\end{center}

\bigskip

This appendix contains the auxiliary lemmas and proofs for the main manuscript.
Equation, theorem, proposition, corollary, and assumption numbers refer to the
main manuscript unless otherwise stated.

\section{Auxiliary Lemmas and Proofs}

\subsection{Auxiliary lemmas}
\label{subsec:auxiliary_lemmas}

The proofs in this supplement rely on four external results. We collect them
here for ease of reference. The constants appearing in these lemmas are
universal (i.e., independent of \(n_1,n_2,d\)) unless stated otherwise.

\begin{lemma}[Gaussian approximation on Euclidean balls; \citealp{fang_large-dimensional_2021}, Theorem~3.2]
\label{lem:fang_koike_balls}
Let \(\zeta_1,\ldots,\zeta_N\) be independent centered random vectors in
\(\mathbb R^d\) with finite fourth moments, and set
\(W_\zeta=\sum_{\ell=1}^N\zeta_\ell\) and
\(\Sigma_W=\operatorname{Var}(W_\zeta)\). Let \(Z\sim N(0,\Sigma)\) and assume
\(\Lambda_2(\Sigma)>0\). Then
\[
\sup_{r\ge 0}
\bigl|
\mathbb P(\|W_\zeta\|_2\le r)
-\mathbb P(\|Z\|_2\le r)
\bigr|
\le
K\bigl\{
\kappa(\Sigma)^{3/4}\delta_1(W_\zeta)^{1/4}
+\kappa(\Sigma)^{2/3}\delta_2(W_\zeta)^{1/3}
\]
\[
\qquad\qquad
+\kappa(\Sigma)^{2/3}\delta_0(\Sigma_W,\Sigma)^{1/3}
+\kappa(\Sigma)^{1/2}\delta_0'(\Sigma_W,\Sigma)^{1/2}
\bigr\},
\]
where \(K_0>0\) is a universal absolute constant and
\[
\delta_0(\Sigma_W,\Sigma)
=
\sqrt{
(\operatorname{tr}(\Sigma_W)+\operatorname{tr}(\Sigma))
(\|\Sigma_W\|_{\mathrm{op}}+\|\Sigma\|_{\mathrm{op}})
}\,
\|\Sigma-\Sigma_W\|_{\mathrm{HS}},
\]
\[
\delta_0'(\Sigma_W,\Sigma)
=
\sum_{j=1}^d|\Sigma_{jj}-\Sigma_{W,jj}|,
\]
\[
\delta_1(W_\zeta)
=
\|\Sigma_W\|_{\mathrm{HS}}\sum_{\ell=1}^N\mathbb E\|\zeta_\ell\|_2^4
+
\|\Sigma_W\|_{\mathrm{op}}^{3/2}\sum_{\ell=1}^N\mathbb E\|\zeta_\ell\|_2^3,
\]
\[
\delta_2(W_\zeta)
=
\|\Sigma_W\|_{\mathrm{op}}^{1/2}\sum_{\ell=1}^N\mathbb E\|\zeta_\ell\|_2^3
+
\sum_{\ell=1}^N\mathbb E\|\zeta_\ell\|_2^4.
\]
\end{lemma}

\begin{lemma}[Anti-concentration for the squared norm of a Gaussian vector;
\citealp{fang_large-dimensional_2021}, Lemma~5.10, originating from
\citealp{gotze_large_2019}, Theorem~2.7]
\label{lem:gaussian_anticoncentration}
Let \(Z\sim N(0,\Sigma)\) with \(\Lambda_2(\Sigma)>0\). Then for every
\(\mu\in\mathbb R^d\) and every \(\varepsilon>0\),
\[
\sup_{a\ge 0}
\mathbb P\bigl(a\le \|Z+\mu\|_2^2\le a+\varepsilon\bigr)
\le
K\kappa(\Sigma)\varepsilon,
\]
where \(K_0>0\) is a universal absolute constant.
\end{lemma}

\begin{lemma}[Gaussian comparison for squared Euclidean norms]
\label{lem:gaussian_comparison_quadratic}
Let \(R_1,R_2\) be \(d\times d\) symmetric positive semidefinite matrices,
and let \(Z_m\sim N(0,R_m)\) for \(m=1,2\).  For each \(m\), let
\[
\lambda_1(R_m)\ge\cdots\ge\lambda_d(R_m)\ge0
\]
be the eigenvalues of \(R_m\), and define
\[
\Lambda_s(R_m)
=
\left(\sum_{\ell=s}^d\lambda_\ell(R_m)^2\right)^{1/2},
\qquad s=1,2.
\]
Assume
\[
\Lambda_2(R_1)>0,
\qquad
\Lambda_2(R_2)>0,
\]
and set
\[
\kappa_*(R_m)
=
\{\Lambda_1(R_m)\Lambda_2(R_m)\}^{-1/2}.
\]
Then there exists a universal constant \(K>0\) such that
\[
\sup_{r\ge0}
\left|
\mathbb P(\|Z_1\|_2^2\le r)
-
\mathbb P(\|Z_2\|_2^2\le r)
\right|
\le
K\{\kappa_*(R_1)+\kappa_*(R_2)\}
\|R_1-R_2\|_{S_1},
\]
where \(\|\cdot\|_{S_1}\) denotes the Schatten-\(1\), or nuclear, norm.
\end{lemma}

\begin{proof}
We apply Theorem~2.1 and Corollary~2.3 of
\citet{gotze_large_2019}.  Their Gaussian comparison result gives, for two
centered Gaussian elements with covariance operators \(R_1\) and \(R_2\),
\[
\sup_{x\ge0}
\left|
\mathbb P(\|Z_1\|_2\le x)
-
\mathbb P(\|Z_2\|_2\le x)
\right|
\le
K
\{\kappa_G(R_1)+\kappa_G(R_2)\}
\|R_1-R_2\|_{S_1},
\]
where \(\kappa_G(\cdot)\) denotes the geometry factor used in
\citet{gotze_large_2019}.  Corollary~2.3 of that paper is precisely the
nuclear-norm version of the comparison bound.

Their \(\kappa_G(R)\) is equivalent, up to universal constants, to
\[
\{\Lambda_1(R)\Lambda_2(R)\}^{-1/2}.
\]
Therefore, after absorbing this equivalence into the universal constant \(K\),
we obtain
\[
\sup_{x\ge0}
\left|
\mathbb P(\|Z_1\|_2\le x)
-
\mathbb P(\|Z_2\|_2\le x)
\right|
\le
K
\{\kappa_*(R_1)+\kappa_*(R_2)\}
\|R_1-R_2\|_{S_1}.
\]

Finally, the map \(r\mapsto\sqrt r\) is a bijection from \([0,\infty)\) to
\([0,\infty)\).  Hence
\[
\sup_{r\ge0}
\left|
\mathbb P(\|Z_1\|_2^2\le r)
-
\mathbb P(\|Z_2\|_2^2\le r)
\right|
=
\sup_{x\ge0}
\left|
\mathbb P(\|Z_1\|_2\le x)
-
\mathbb P(\|Z_2\|_2\le x)
\right|.
\]
This proves the claim.
\end{proof}

\begin{lemma}[Hoffman--Wielandt inequality]
\label{lem:hoffman_wielandt}
Let \(A,B\) be \(d\times d\) real symmetric matrices with eigenvalues
\(\lambda_1(A)\ge\cdots\ge\lambda_d(A)\) and
\(\lambda_1(B)\ge\cdots\ge\lambda_d(B)\). Then
\[
\sum_{\ell=1}^d
(\lambda_\ell(A)-\lambda_\ell(B))^2
\le
\|A-B\|_{\mathrm F}^2.
\]
\end{lemma}

\begin{lemma}[Euclidean moment consequences of coordinate moments]
\label{lem:coordinate_moments_imply_euclidean}
Suppose Assumption~\ref{ass:coordinate_high_moment} holds.  Then there exists a
constant \(C<\infty\), depending only on \(K_8\), such that
\[
\operatorname{tr}\{\operatorname{Cov}(U_i)\}
+
\operatorname{tr}\{\operatorname{Cov}(V_j)\}
\le Cp,
\]
\[
\mathbb E\|U_i\|_2^4+\mathbb E\|V_j\|_2^4\le Cp^2,
\qquad
\mathbb E\|U_i\|_2^8+\mathbb E\|V_j\|_2^8\le Cp^4 .
\]
Moreover, for independent centered observations satisfying the same coordinate
moment bound,
\[
\mathbb E\|\bar U\|_2^4\le C\left(\frac p{n_1}\right)^2,
\qquad
\mathbb E\|\bar U\|_2^8\le C\left(\frac p{n_1}\right)^4,
\]
and the analogous bounds hold for \(\bar V\).  Finally,
\[
\mathbb E
\left(
\frac1{n_1}\sum_{i=1}^{n_1}\|U_i-\bar U\|_2^2
\right)^4
\le Cp^4,
\]
with the same conclusion for the second sample.
\end{lemma}

\begin{proof}
We prove the statements for \(U_i\); the argument for \(V_j\) is identical.
By Lyapunov's inequality,
\(\mathbb E U_{ia}^2\le K_8^{1/4}\) and
\(\mathbb E U_{ia}^4\le K_8^{1/2}\) uniformly in \(a\).  Hence
\[
\operatorname{tr}\{\operatorname{Cov}(U_i)\}
\le
\sum_{a=1}^p \mathbb E U_{ia}^2
\le Cp.
\]
For the fourth moment, with \(z_a=U_{ia}^2\),
\[
\|U_i\|_2^4
=
\left(\sum_{a=1}^p z_a\right)^2
\le
p\sum_{a=1}^p z_a^2
=
p\sum_{a=1}^p U_{ia}^4,
\]
so \(\mathbb E\|U_i\|_2^4\le Cp^2\).  Similarly,
\[
\|U_i\|_2^8
=
\left(\sum_{a=1}^p z_a\right)^4
\le
p^3\sum_{a=1}^p z_a^4
=
p^3\sum_{a=1}^p U_{ia}^8,
\]
and therefore \(\mathbb E\|U_i\|_2^8\le Cp^4\).

For the sample mean, the standard scalar Rosenthal inequality gives
\[
\mathbb E|\bar U_a|^4\le Cn_1^{-2},
\qquad
\mathbb E|\bar U_a|^8\le Cn_1^{-4}
\]
uniformly in \(a\).  Applying the preceding sum inequalities to
\(\bar U=(\bar U_1,\ldots,\bar U_p)^\top\) yields
\[
\mathbb E\|\bar U\|_2^4\le C\left(\frac p{n_1}\right)^2,
\qquad
\mathbb E\|\bar U\|_2^8\le C\left(\frac p{n_1}\right)^4.
\]
Finally,
\[
\frac1{n_1}\sum_{i=1}^{n_1}\|U_i-\bar U\|_2^2
\le
\frac2{n_1}\sum_{i=1}^{n_1}\|U_i\|_2^2+2\|\bar U\|_2^2.
\]
Using \((x+y)^4\le 8(x^4+y^4)\), Jensen's inequality, and the eighth-moment
bound just proved,
\[
\mathbb E
\left(
\frac1{n_1}\sum_{i=1}^{n_1}\|U_i-\bar U\|_2^2
\right)^4
\le
C\mathbb E\|U_i\|_2^8
+
C\mathbb E\|\bar U\|_2^8
\le Cp^4 .
\]
This proves the lemma.
\end{proof}

\subsection{Proofs}
\label{subsec:proofs}

\begin{proof}[Proof of Theorem~\ref{thm:oracle_gaussian}]
Throughout the proof, we work in the transformed model and suppress the
superscript \((A)\). Under \(H_0:\Sigma_1=\Sigma_2\), the transformed covariance
matrices are equal.

For \(1\le a\le b\le p\), let \(k=k(a,b)\) denote the corresponding coordinate
of the half-vectorization. Define
\[
H_i^{(1)}
=
\operatorname{vech}
\{(X_i-\mu_1)(X_i-\mu_1)^\top-\Sigma_1\},
\]
and
\[
H_j^{(2)}
=
\operatorname{vech}
\{(Y_j-\mu_2)(Y_j-\mu_2)^\top-\Sigma_2\}.
\]
Equivalently, in the notation of the primitive assumptions,
\[
H_{i,k}^{(1)}
=
U_{ia}U_{ib}-\mathbb E(U_{ia}U_{ib}),
\qquad
H_{j,k}^{(2)}
=
V_{ja}V_{jb}-\mathbb E(V_{ja}V_{jb}).
\]

We first record the standardized moment bounds for these covariance-coordinate
summands. By Assumption~\ref{ass:coordinate_high_moment}, for every
\(1\le a\le b\le p\),
\[
\mathbb E |U_{ia}U_{ib}|^4
\le
\{\mathbb E|U_{ia}|^8\mathbb E|U_{ib}|^8\}^{1/2}
\le K_8,
\]
and similarly for \(V_j\). Hence,
\[
\mathbb E
\left|
U_{ia}U_{ib}-\mathbb E(U_{ia}U_{ib})
\right|^4
\le
16K_8.
\]
By Lyapunov's inequality,
\[
\mathbb E
\left|
U_{ia}U_{ib}-\mathbb E(U_{ia}U_{ib})
\right|^3
\le
(16K_8)^{3/4}.
\]
The same bounds hold for the second sample. Together with
Assumption~\ref{ass:coordinate_nondegeneracy}, this implies that, for
\(q=3,4\),
\[
\max_{m=1,2}
\max_{1\le k\le d}
\frac{
\mathbb E|H_{1,k}^{(m)}|^q
}{
\operatorname{Var}(H_{1,k}^{(m)})^{q/2}
}
\le M_q,
\]
where \(M_q<\infty\) depends only on \(K_8\) and \(c_H\).

Next, under \(H_0\),
\[
\widetilde W
=
\sqrt{n_{\mathrm{eff}}}
\left\{
\frac1{n_1}\sum_{i=1}^{n_1}H_i^{(1)}
-
\frac1{n_2}\sum_{j=1}^{n_2}H_j^{(2)}
\right\}.
\]
For each coordinate \(k\),
\[
\widetilde D_{kk}
=
\operatorname{Var}(\widetilde W_k)
=
n_{\mathrm{eff}}
\left\{
\frac{\operatorname{Var}(H_{1,k}^{(1)})}{n_1}
+
\frac{\operatorname{Var}(H_{1,k}^{(2)})}{n_2}
\right\}.
\]
By Assumption~\ref{ass:coordinate_nondegeneracy},
\[
\operatorname{Var}(H_{1,k}^{(1)})\ge c_H,
\qquad
\operatorname{Var}(H_{1,k}^{(2)})\ge c_H.
\]
Therefore,
\[
\widetilde D_{kk}
\ge
n_{\mathrm{eff}}c_H
\left(
\frac1{n_1}+\frac1{n_2}
\right)
=
c_H,
\]
because
\[
n_{\mathrm{eff}}
=
\frac{n_1n_2}{n_1+n_2}.
\]
Thus \(\widetilde D^{-1/2}\) is well defined.

Define independent centered random vectors
\[
\xi_i
=
\frac{\sqrt{n_{\mathrm{eff}}}}{n_1}
\widetilde D^{-1/2}H_i^{(1)},
\qquad
1\le i\le n_1,
\]
and
\[
\xi_{n_1+j}
=
-
\frac{\sqrt{n_{\mathrm{eff}}}}{n_2}
\widetilde D^{-1/2}H_j^{(2)},
\qquad
1\le j\le n_2.
\]
Let
\[
S_n
=
\sum_{\ell=1}^{n_1+n_2}\xi_\ell.
\]
Then
\begin{equation}
\label{eq:supp-oracle-standardized-sum}
S_n
=
\widetilde D^{-1/2}\widetilde W.
\end{equation}
Therefore,
\begin{equation}
\label{eq:supp-oracle-statistic}
\widetilde T_n
=
\frac1d\|S_n\|_2^2.
\end{equation}
Moreover, by independence of the \(\xi_\ell\)'s,
\[
\operatorname{Cov}(S_n)
=
\sum_{\ell=1}^{n_1+n_2}\operatorname{Cov}(\xi_\ell)
=
\sum_{i=1}^{n_1}\frac{n_{\mathrm{eff}}}{n_1^2}
\widetilde D^{-1/2}\Gamma_1\widetilde D^{-1/2}
+
\sum_{j=1}^{n_2}\frac{n_{\mathrm{eff}}}{n_2^2}
\widetilde D^{-1/2}\Gamma_2\widetilde D^{-1/2},
\]
where \(\Gamma_1=\operatorname{Cov}(H_i^{(1)})\) and
\(\Gamma_2=\operatorname{Cov}(H_j^{(2)})\). Simplifying,
\[
\operatorname{Cov}(S_n)
=
\widetilde D^{-1/2}
\left\{
n_{\mathrm{eff}}
\left(
\frac{\Gamma_1}{n_1}
+
\frac{\Gamma_2}{n_2}
\right)
\right\}
\widetilde D^{-1/2}
=
\widetilde D^{-1/2}
\operatorname{Cov}(\widetilde W)
\widetilde D^{-1/2}
=
R.
\]
Let
\[
Z\sim N(0,R).
\]

We next bound the third and fourth Euclidean moments of the summands
\(\xi_\ell\). For \(1\le i\le n_1\),
\[
\xi_{i,k}
=
\frac{\sqrt{n_{\mathrm{eff}}}}{n_1}
\frac{H_{i,k}^{(1)}}{\sqrt{\widetilde D_{kk}}}.
\]
Since
\[
\widetilde D_{kk}
\ge
\frac{n_{\mathrm{eff}}}{n_1}
\operatorname{Var}(H_{1,k}^{(1)}),
\]
we have, for \(q=3,4\),
\[
\mathbb E|\xi_{i,k}|^q
=
\left(\frac{\sqrt{n_{\mathrm{eff}}}}{n_1}\right)^q
\frac{\mathbb E|H_{i,k}^{(1)}|^q}{\widetilde D_{kk}^{q/2}}
\le
n_1^{-q/2}
\frac{\mathbb E|H_{i,k}^{(1)}|^q}
{\operatorname{Var}(H_{1,k}^{(1)})^{q/2}}
\le
M_q n_1^{-q/2}.
\]
Similarly, for \(1\le j\le n_2\),
\[
\mathbb E|\xi_{n_1+j,k}|^q
\le
M_q n_2^{-q/2},
\qquad q=3,4.
\]

For any \(x\in\mathbb R^d\),
\[
\|x\|_2^3
=
\left(\sum_{k=1}^d x_k^2\right)^{3/2}
\le
d^{1/2}\sum_{k=1}^d |x_k|^3,
\]
and
\[
\|x\|_2^4
=
\left(\sum_{k=1}^d x_k^2\right)^2
\le
d\sum_{k=1}^d |x_k|^4.
\]
Therefore,
\[
\sum_{\ell=1}^{n_1+n_2}
\mathbb E\|\xi_\ell\|_2^3
\le
d^{1/2}
\sum_{\ell=1}^{n_1+n_2}
\sum_{k=1}^d
\mathbb E|\xi_{\ell,k}|^3.
\]
Using the coordinatewise bounds,
\begin{equation}
\label{eq:supp-oracle-third-moment}
\sum_{\ell=1}^{n_1+n_2}
\mathbb E\|\xi_\ell\|_2^3
\le
K_0 d^{3/2}
\left(
n_1^{-1/2}+n_2^{-1/2}
\right).
\end{equation}
Similarly,
\begin{equation}
\label{eq:supp-oracle-fourth-moment}
\sum_{\ell=1}^{n_1+n_2}
\mathbb E\|\xi_\ell\|_2^4
\le
d
\sum_{\ell=1}^{n_1+n_2}
\sum_{k=1}^d
\mathbb E|\xi_{\ell,k}|^4
\le
K_0 d^2
\left(
n_1^{-1}+n_2^{-1}
\right).
\end{equation}

We now apply Lemma~\ref{lem:fang_koike_balls} (the
\citealp{fang_large-dimensional_2021} fourth-moment Gaussian
approximation on Euclidean balls) to the standardized sum
\(S_n=\sum_{\ell=1}^{n_1+n_2}\xi_\ell\), which satisfies
\eqref{eq:supp-oracle-standardized-sum} and \(\operatorname{Cov}(S_n)=R\),
with reference covariance \(\Sigma=R\) (so the centered Gaussian reference
\(Z\sim N(0,R)\) has the same covariance matrix as \(S_n\)). In the notation of
Lemma~\ref{lem:fang_koike_balls}, this corresponds to taking
\(\Sigma_W=\operatorname{Var}(W_\zeta)=\Sigma=R\). Consequently, the two
covariance-comparison terms vanish identically:
\[
\delta_0(\Sigma_W,\Sigma)
=
\sqrt{2\operatorname{tr}(R)\cdot 2\|R\|_{\mathrm{op}}}\,\|R-R\|_{\mathrm{HS}}
=
0,
\qquad
\delta_0'(\Sigma_W,\Sigma)
=
\sum_{j=1}^d|R_{jj}-R_{jj}|
=0.
\]
Lemma~\ref{lem:fang_koike_balls} therefore reduces, under the present setting,
to
\[
\sup_{B\in\mathcal B}
\left|
\mathbb P(S_n\in B)-\mathbb P(Z\in B)
\right|
\le
K
\left\{
\kappa(R)^{3/4}\Delta_1^{1/4}
+
\kappa(R)^{2/3}\Delta_2^{1/3}
\right\},
\]
where \(\mathcal B\) denotes the class of Euclidean balls in \(\mathbb R^d\),
and
\[
\Delta_1
\le
\|R\|_{\mathrm{HS}}
\sum_{\ell=1}^{n_1+n_2}\mathbb E\|\xi_\ell\|_2^4
+
\|R\|_{\mathrm{op}}^{3/2}
\sum_{\ell=1}^{n_1+n_2}\mathbb E\|\xi_\ell\|_2^3,
\]
\[
\Delta_2
\le
\|R\|_{\mathrm{op}}^{1/2}
\sum_{\ell=1}^{n_1+n_2}\mathbb E\|\xi_\ell\|_2^3
+
\sum_{\ell=1}^{n_1+n_2}\mathbb E\|\xi_\ell\|_2^4.
\]
Substituting \eqref{eq:supp-oracle-third-moment} and
\eqref{eq:supp-oracle-fourth-moment} yields
\[
\Delta_1
\le
K_0
\left[
\|R\|_{\mathrm{HS}}
d^2
\left(n_1^{-1}+n_2^{-1}\right)
+
\|R\|_{\mathrm{op}}^{3/2}
d^{3/2}
\left(n_1^{-1/2}+n_2^{-1/2}\right)
\right],
\]
and
\[
\Delta_2
\le
K_0
\left[
\|R\|_{\mathrm{op}}^{1/2}
d^{3/2}
\left(n_1^{-1/2}+n_2^{-1/2}\right)
+
d^2
\left(n_1^{-1}+n_2^{-1}\right)
\right].
\]
Absorbing constants into \(K\), define
\[
\delta_1
=
\|R\|_{\mathrm{HS}}
d^2
\left(n_1^{-1}+n_2^{-1}\right)
+
\|R\|_{\mathrm{op}}^{3/2}
d^{3/2}
\left(n_1^{-1/2}+n_2^{-1/2}\right),
\]
and
\[
\delta_2
=
\|R\|_{\mathrm{op}}^{1/2}
d^{3/2}
\left(n_1^{-1/2}+n_2^{-1/2}\right)
+
d^2
\left(n_1^{-1}+n_2^{-1}\right).
\]
Thus
\begin{equation}
\label{eq:supp-oracle-ball-approx}
\sup_{B\in\mathcal B}
\left|
\mathbb P(S_n\in B)-\mathbb P(Z\in B)
\right|
\le
K
\left\{
\kappa(R)^{3/4}\delta_1^{1/4}
+
\kappa(R)^{2/3}\delta_2^{1/3}
\right\}.
\end{equation}

Finally, by \eqref{eq:supp-oracle-statistic}, for any \(r\ge0\),
\[
\{\widetilde T_n\le r\}
=
\left\{
\frac1d\|S_n\|_2^2\le r
\right\}
=
\left\{
\|S_n\|_2\le \sqrt{dr}
\right\},
\]
which is a Euclidean ball centered at the origin. Similarly,
\[
\{T_Z\le r\}
=
\left\{
\|Z\|_2\le \sqrt{dr}
\right\}.
\]
Therefore,
\[
\left|
\mathbb P(\widetilde T_n\le r)
-
\mathbb P(T_Z\le r)
\right|
\le
\sup_{B\in\mathcal B}
\left|
\mathbb P(S_n\in B)-\mathbb P(Z\in B)
\right|.
\]
Taking the supremum over \(r\ge0\) and using
\eqref{eq:supp-oracle-ball-approx} gives
\[
\sup_{r\ge0}
\left|
\mathbb P(\widetilde T_n\le r)
-
\mathbb P(T_Z\le r)
\right|
\le
K
\left\{
\kappa(R)^{3/4}\delta_1^{1/4}
+
\kappa(R)^{2/3}\delta_2^{1/3}
\right\}.
\]
This proves the theorem.
\end{proof}


The oracle approximation concerns the ideal statistic
\(\widetilde T_n\), which is defined using the true means and the oracle
diagonal variance matrix \(\widetilde D\). The statistic used in practice,
however, replaces the population means by sample means and replaces
\(\widetilde D\) by the feasible diagonal estimator \(\widehat D\). We first
show that the diagonal studentizer is uniformly consistent. This result is then
combined with a sample-centering bound to transfer the oracle approximation to
the feasible statistic.

\begin{proposition}[Uniform consistency of the diagonal studentizer]
\label{prop:studentizer_consistency}
Suppose Assumptions~\ref{ass:coordinate_high_moment} and
\ref{ass:coordinate_nondegeneracy} hold. Let
\[
  \widehat D_{kk}
  =
  n_{\mathrm{eff}}
  \left(
    \frac1{n_1^2}\sum_{i=1}^{n_1}
    \bigl(\widehat h_{i,k}^{(1)}\bigr)^2
    +
    \frac1{n_2^2}\sum_{j=1}^{n_2}
    \bigl(\widehat h_{j,k}^{(2)}\bigr)^2
  \right),
  \qquad 1\le k\le d.
\]
Then there exists a constant \(K_0>0\), depending only on the constants in
Assumptions~\ref{ass:coordinate_high_moment} and
\ref{ass:coordinate_nondegeneracy}, such that
\begin{equation}
\label{eq:supp-studentizer-mse}
\max_{1\le k\le d}
\mathbb E
\left[
\left(
\frac{\widehat D_{kk}}{\widetilde D_{kk}}-1
\right)^2
\right]
\le
\frac{K_0}{n_0},
\qquad
n_0=n_1\wedge n_2.
\end{equation}
Consequently, for every \(\eta\in(0,1)\),
\begin{equation}
\label{eq:supp-studentizer-hp}
\mathbb P
\left(
\max_{1\le k\le d}
\left|
\frac{\widehat D_{kk}}{\widetilde D_{kk}}-1
\right|
>
K_0\sqrt{\frac{d}{n_0\eta}}
\right)
\le \eta .
\end{equation}
\end{proposition}

\begin{proof}[Proof of Proposition \ref{prop:studentizer_consistency}]
We prove the result for a fixed coordinate \(k=k(a,b)\), \(1\le a\le b\le p\).
The bound is uniform in \(k\), and the final high-probability statement follows
from a union bound.

We first consider the first sample. Let
\[
G_i=U_{ia}U_{ib},
\qquad
\mu_G=\mathbb E G_i,
\qquad
H_i=G_i-\mu_G,
\qquad
\sigma_{1,k}^2=\operatorname{Var}(G_i)=\mathbb E H_i^2.
\]
Let
\[
\bar G=\frac1{n_1}\sum_{i=1}^{n_1}G_i,
\qquad
\bar H=\frac1{n_1}\sum_{i=1}^{n_1}H_i.
\]
Define the oracle sample variance of the covariance-coordinate variable by
\[
\check\sigma_{1,k}^2
=
\frac1{n_1}\sum_{i=1}^{n_1}(G_i-\bar G)^2
=
\frac1{n_1}\sum_{i=1}^{n_1}(H_i-\bar H)^2.
\]
Since
\[
\check\sigma_{1,k}^2
=
\frac1{n_1}\sum_{i=1}^{n_1}H_i^2-\bar H^2,
\]
we have
\[
\check\sigma_{1,k}^2-\sigma_{1,k}^2
=
\frac1{n_1}\sum_{i=1}^{n_1}(H_i^2-\sigma_{1,k}^2)-\bar H^2.
\]
Therefore,
\[
\mathbb E
\left(
\check\sigma_{1,k}^2-\sigma_{1,k}^2
\right)^2
\le
2\mathbb E
\left[
\frac1{n_1}\sum_{i=1}^{n_1}(H_i^2-\sigma_{1,k}^2)
\right]^2
+
2\mathbb E\bar H^4.
\]
By independence,
\[
\mathbb E
\left[
\frac1{n_1}\sum_{i=1}^{n_1}(H_i^2-\sigma_{1,k}^2)
\right]^2
=
\frac1{n_1}\operatorname{Var}(H_i^2)
\le
\frac1{n_1}\mathbb E H_i^4.
\]
Assumption~\ref{ass:coordinate_high_moment} implies
\[
\mathbb E H_i^4\le K_0,
\]
uniformly over \(k\). Moreover,
\[
\mathbb E\bar H^4
\le
K_0 n_1^{-2}\mathbb E H_i^4
\le
\frac{K_0}{n_1^2}.
\]
Hence
\[
\mathbb E
\left(
\check\sigma_{1,k}^2-\sigma_{1,k}^2
\right)^2
\le
\frac{K_0}{n_1}.
\]

We now compare the feasible residuals with the oracle residuals. Recall that
\[
\widehat h_{i,k}^{(1)}
=
(U_{ia}-\bar U_a)(U_{ib}-\bar U_b)-\widehat\Sigma_{1,ab},
\]
where
\[
\widehat\Sigma_{1,ab}
=
\frac1{n_1}\sum_{\ell=1}^{n_1}
(U_{\ell a}-\bar U_a)(U_{\ell b}-\bar U_b).
\]
Expanding the centered product gives
\[
(U_{ia}-\bar U_a)(U_{ib}-\bar U_b)
=
U_{ia}U_{ib}
-\bar U_aU_{ib}
-\bar U_bU_{ia}
+\bar U_a\bar U_b.
\]
Also,
\[
\widehat\Sigma_{1,ab}
=
\bar G-\bar U_a\bar U_b.
\]
Therefore,
\[
\widehat h_{i,k}^{(1)}
=
G_i-\bar G
-
\bar U_a(U_{ib}-\bar U_b)
-
\bar U_b(U_{ia}-\bar U_a).
\]
Write
\[
\widehat h_{i,k}^{(1)}
=
(G_i-\bar G)+R_i,
\]
where
\[
R_i
=
-\bar U_a(U_{ib}-\bar U_b)
-
\bar U_b(U_{ia}-\bar U_a).
\]

Define
\[
\widehat\sigma_{1,k}^2
=
\frac1{n_1}\sum_{i=1}^{n_1}
\bigl(\widehat h_{i,k}^{(1)}\bigr)^2.
\]
Then
\[
\widehat\sigma_{1,k}^2-\check\sigma_{1,k}^2
=
\frac1{n_1}\sum_{i=1}^{n_1}
\{2(G_i-\bar G)R_i+R_i^2\}.
\]
By Cauchy's inequality,
\[
\left|
\widehat\sigma_{1,k}^2-\check\sigma_{1,k}^2
\right|
\le
2
\left[
\frac1{n_1}\sum_{i=1}^{n_1}(G_i-\bar G)^2
\right]^{1/2}
\left[
\frac1{n_1}\sum_{i=1}^{n_1}R_i^2
\right]^{1/2}
+
\frac1{n_1}\sum_{i=1}^{n_1}R_i^2.
\]
Under Assumption~\ref{ass:coordinate_high_moment}, the fourth moment of
\(G_i=U_{ia}U_{ib}\) is uniformly bounded. Hence
\[
\mathbb E
\left[
\frac1{n_1}\sum_{i=1}^{n_1}(G_i-\bar G)^2
\right]^2
\le K_0.
\]

Next,
\[
R_i^2
\le
2\bar U_a^2(U_{ib}-\bar U_b)^2
+
2\bar U_b^2(U_{ia}-\bar U_a)^2.
\]
Thus
\[
\frac1{n_1}\sum_{i=1}^{n_1}R_i^2
\le
2\bar U_a^2
\frac1{n_1}\sum_{i=1}^{n_1}(U_{ib}-\bar U_b)^2
+
2\bar U_b^2
\frac1{n_1}\sum_{i=1}^{n_1}(U_{ia}-\bar U_a)^2.
\]
The coordinate eighth-moment bound implies
\[
\mathbb E\bar U_a^4\le \frac{K_0}{n_1^2},
\]
and
\[
\mathbb E
\left[
\frac1{n_1}\sum_{i=1}^{n_1}(U_{ib}-\bar U_b)^2
\right]^2
\le K_0,
\]
uniformly over \(a,b\). Hence
\[
\mathbb E
\left(
\frac1{n_1}\sum_{i=1}^{n_1}R_i^2
\right)^2
\le
\frac{K_0}{n_1^2}.
\]
Combining the preceding displays yields
\[
\mathbb E
\left(
\widehat\sigma_{1,k}^2-\check\sigma_{1,k}^2
\right)^2
\le
\frac{K_0}{n_1}.
\]
Therefore,
\[
\mathbb E
\left(
\widehat\sigma_{1,k}^2-\sigma_{1,k}^2
\right)^2
\le
\frac{K_0}{n_1}.
\]

The same argument applied to the second sample gives
\[
\mathbb E
\left(
\widehat\sigma_{2,k}^2-\sigma_{2,k}^2
\right)^2
\le
\frac{K_0}{n_2},
\]
where
\[
\sigma_{2,k}^2=\operatorname{Var}(V_{ja}V_{jb}),
\qquad
\widehat\sigma_{2,k}^2
=
\frac1{n_2}\sum_{j=1}^{n_2}
\bigl(\widehat h_{j,k}^{(2)}\bigr)^2.
\]

Now observe that
\[
\widehat D_{kk}
=
n_{\mathrm{eff}}
\left(
\frac{\widehat\sigma_{1,k}^2}{n_1}
+
\frac{\widehat\sigma_{2,k}^2}{n_2}
\right),
\]
whereas
\[
\widetilde D_{kk}
=
n_{\mathrm{eff}}
\left(
\frac{\sigma_{1,k}^2}{n_1}
+
\frac{\sigma_{2,k}^2}{n_2}
\right).
\]
Therefore,
\[
\widehat D_{kk}-\widetilde D_{kk}
=
n_{\mathrm{eff}}
\left[
\frac{\widehat\sigma_{1,k}^2-\sigma_{1,k}^2}{n_1}
+
\frac{\widehat\sigma_{2,k}^2-\sigma_{2,k}^2}{n_2}
\right].
\]
Using \((x+y)^2\le 2x^2+2y^2\), we get
\[
\mathbb E(\widehat D_{kk}-\widetilde D_{kk})^2
\le
2n_{\mathrm{eff}}^2
\left[
\frac1{n_1^2}
\mathbb E(\widehat\sigma_{1,k}^2-\sigma_{1,k}^2)^2
+
\frac1{n_2^2}
\mathbb E(\widehat\sigma_{2,k}^2-\sigma_{2,k}^2)^2
\right].
\]
Hence
\[
\mathbb E(\widehat D_{kk}-\widetilde D_{kk})^2
\le
K_0 n_{\mathrm{eff}}^2
\left(
\frac1{n_1^3}
+
\frac1{n_2^3}
\right).
\]
By Assumption~\ref{ass:coordinate_nondegeneracy},
\[
\widetilde D_{kk}
=
n_{\mathrm{eff}}
\left(
\frac{\sigma_{1,k}^2}{n_1}
+
\frac{\sigma_{2,k}^2}{n_2}
\right)
\ge
n_{\mathrm{eff}}c_H
\left(
\frac1{n_1}+\frac1{n_2}
\right)
=
c_H.
\]
Thus
\[
\mathbb E
\left[
\left(
\frac{\widehat D_{kk}}{\widetilde D_{kk}}-1
\right)^2
\right]
\le
K_0 n_{\mathrm{eff}}^2
\left(
\frac1{n_1^3}
+
\frac1{n_2^3}
\right).
\]
Since
\[
n_{\mathrm{eff}}\le n_0=n_1\wedge n_2,
\]
we have
\[
n_{\mathrm{eff}}^2
\left(
\frac1{n_1^3}
+
\frac1{n_2^3}
\right)
\le
\frac{K_0}{n_0}.
\]
Therefore,
\[
\max_{1\le k\le d}
\mathbb E
\left[
\left(
\frac{\widehat D_{kk}}{\widetilde D_{kk}}-1
\right)^2
\right]
\le
\frac{K_0}{n_0}.
\]

This proves \eqref{eq:supp-studentizer-mse}.  Finally, for any \(t>0\),
Markov's inequality and the union bound give
\[
\mathbb P
\left(
\max_{1\le k\le d}
\left|
\frac{\widehat D_{kk}}{\widetilde D_{kk}}-1
\right|
>t
\right)
\le
\sum_{k=1}^d
\frac{
\mathbb E
\left[
\left(
\widehat D_{kk}/\widetilde D_{kk}-1
\right)^2
\right]
}{t^2}
\le
\frac{Cd}{n_0t^2}.
\]
Taking
\[
t=K_0\sqrt{\frac{d}{n_0\eta}}
\]
with \(K_0\) sufficiently large yields
\[
\mathbb P
\left(
\max_{1\le k\le d}
\left|
\frac{\widehat D_{kk}}{\widetilde D_{kk}}-1
\right|
>
K_0\sqrt{\frac{d}{n_0\eta}}
\right)
\le \eta.
\]
which proves \eqref{eq:supp-studentizer-hp}.  The proposition follows.
\end{proof}

Next we state the proof of Proposition \ref{prop:sample_oracle_transfer}.
\begin{proof}[Proof of Proposition~\ref{prop:sample_oracle_transfer}]
Throughout the proof, we work in the transformed model and suppress the
superscript \((A)\). Recall that
\[
U_i=X_i-\mu_1,
\qquad
V_j=Y_j-\mu_2.
\]
Let
\[
\widehat T_n
=
\frac1d
\left\|
\widehat D^{-1/2}\widehat W
\right\|_2^2,
\qquad
\widetilde T_n
=
\frac1d
\left\|
\widetilde D^{-1/2}\widetilde W
\right\|_2^2.
\]
We decompose the difference between the feasible statistic and the oracle
statistic into the sample-centering error and the diagonal-studentization error.

Define
\[
a
=
d^{-1/2}\widetilde D^{-1/2}\widetilde W.
\]
Then
\[
\|a\|_2^2
=
\widetilde T_n.
\]
Also define
\[
b
=
d^{-1/2}
\widetilde D^{-1/2}
(\widehat W-\widetilde W).
\]
Thus
\[
d^{-1/2}\widetilde D^{-1/2}\widehat W
=
a+b.
\]

We first control \(\|a\|_2\). Let
\[
S_n=\widetilde D^{-1/2}\widetilde W.
\]
Then
\[
\widetilde T_n=\frac1d\|S_n\|_2^2.
\]
As in the proof of Theorem~\ref{thm:oracle_gaussian}, under
Assumptions~\ref{ass:coordinate_high_moment} and
\ref{ass:coordinate_nondegeneracy}, the standardized covariance-coordinate
summands have uniformly bounded fourth moments. Consequently,
\[
\max_{1\le k\le d}\mathbb E S_{n,k}^4\le K_0.
\]
Therefore, by Cauchy--Schwarz,
\[
\mathbb E\widetilde T_n^2
=
\mathbb E
\left(
\frac1d\sum_{k=1}^d S_{n,k}^2
\right)^2
=
\frac1{d^2}
\sum_{k=1}^d\sum_{\ell=1}^d
\mathbb E(S_{n,k}^2S_{n,\ell}^2)
\]
\[
\le
\frac1{d^2}
\sum_{k=1}^d\sum_{\ell=1}^d
\left(\mathbb E S_{n,k}^4\right)^{1/2}
\left(\mathbb E S_{n,\ell}^4\right)^{1/2}
\le K_0.
\]
Hence, by Markov's inequality,
\[
\mathbb P
\left(
\|a\|_2>B_T(\eta)
\right)
=
\mathbb P
\left(
\sqrt{\widetilde T_n}>K_0\eta^{-1/4}
\right)
\]
\[
=
\mathbb P
\left(
\widetilde T_n^2>K_0^4\eta^{-1}
\right)
\le
\frac{K_0}{K_0^4\eta^{-1}}
\le
\frac{\eta}{3},
\]
after increasing the constant \(K_0\) if necessary. Thus, with probability at
least \(1-\eta/3\),
\[
\|a\|_2\le B_T(\eta),
\]
where
\[
B_T(\eta)=K_0\eta^{-1/4}.
\]

We next control the sample-centering error \(b\). Let
\[
\bar U=\frac1{n_1}\sum_{i=1}^{n_1}U_i,
\qquad
\bar V=\frac1{n_2}\sum_{j=1}^{n_2}V_j.
\]
The oracle covariance matrices are
\[
\widetilde\Sigma_1
=
\frac1{n_1}\sum_{i=1}^{n_1}U_iU_i^\top,
\qquad
\widetilde\Sigma_2
=
\frac1{n_2}\sum_{j=1}^{n_2}V_jV_j^\top.
\]
The feasible covariance matrices are
\[
\widehat\Sigma_1
=
\frac1{n_1}\sum_{i=1}^{n_1}(U_i-\bar U)(U_i-\bar U)^\top,
\]
and
\[
\widehat\Sigma_2
=
\frac1{n_2}\sum_{j=1}^{n_2}(V_j-\bar V)(V_j-\bar V)^\top.
\]
Using the identity
\[
\frac1n\sum_{i=1}^n(Z_i-\bar Z)(Z_i-\bar Z)^\top
=
\frac1n\sum_{i=1}^n Z_iZ_i^\top-\bar Z\bar Z^\top,
\]
we obtain
\[
\widehat\Sigma_1-\widetilde\Sigma_1
=
-\bar U\bar U^\top,
\qquad
\widehat\Sigma_2-\widetilde\Sigma_2
=
-\bar V\bar V^\top.
\]
Therefore,
\[
\widehat W-\widetilde W
=
\sqrt{n_{\mathrm{eff}}}
\operatorname{vech}
\left\{
-\bar U\bar U^\top+\bar V\bar V^\top
\right\}.
\]

By Assumption~\ref{ass:coordinate_nondegeneracy}, the oracle diagonal
studentizer is uniformly nondegenerate. Indeed, for \(k=k(a,b)\),
\[
\widetilde D_{kk}
=
n_{\mathrm{eff}}
\left\{
\frac{\operatorname{Var}(U_{ia}U_{ib})}{n_1}
+
\frac{\operatorname{Var}(V_{ja}V_{jb})}{n_2}
\right\}
\ge
n_{\mathrm{eff}}c_H
\left(
\frac1{n_1}+\frac1{n_2}
\right)
=
c_H.
\]
Hence
\[
\|\widetilde D^{-1/2}\|_{\mathrm{op}}\le c_H^{-1/2}.
\]
Moreover, for every symmetric matrix \(A\),
\[
\|\operatorname{vech}(A)\|_2\le \|A\|_{\mathrm F}.
\]
Thus
\[
\left\|
\widetilde D^{-1/2}(\widehat W-\widetilde W)
\right\|_2
\le
K_0\sqrt{n_{\mathrm{eff}}}
\left\|
-\bar U\bar U^\top+\bar V\bar V^\top
\right\|_{\mathrm F}.
\]
Using
\[
\|uu^\top\|_{\mathrm F}=\|u\|_2^2,
\]
we get
\[
\left\|
\widetilde D^{-1/2}(\widehat W-\widetilde W)
\right\|_2
\le
K_0\sqrt{n_{\mathrm{eff}}}
\left(
\|\bar U\|_2^2+\|\bar V\|_2^2
\right).
\]
Therefore,
\[
\|b\|_2
=
d^{-1/2}
\left\|
\widetilde D^{-1/2}(\widehat W-\widetilde W)
\right\|_2
\le
K_0
\frac{\sqrt{n_{\mathrm{eff}}}}{\sqrt d}
\left(
\|\bar U\|_2^2+\|\bar V\|_2^2
\right).
\]

We now bound the transformed sample means. By
Lemma~\ref{lem:coordinate_moments_imply_euclidean},
\[
\mathbb E\|\bar U\|_2^4
\le
K_0\left(\frac p{n_1}\right)^2,
\qquad
\mathbb E\|\bar V\|_2^4
\le
K_0\left(\frac p{n_2}\right)^2.
\]

It follows that
\[
\mathbb E
\left(
\|\bar U\|_2^2+\|\bar V\|_2^2
\right)^2
\le
2\mathbb E\|\bar U\|_2^4
+
2\mathbb E\|\bar V\|_2^4
\]
\[
\le
K_0
\left[
\left(\frac p{n_1}\right)^2
+
\left(\frac p{n_2}\right)^2
\right]
\le
K_0
\left(
\frac p{n_1}
+
\frac p{n_2}
\right)^2.
\]
By Markov's inequality, with probability at least \(1-\eta/3\),
\[
\|\bar U\|_2^2+\|\bar V\|_2^2
\le
K_0\eta^{-1/2}
\left(
\frac p{n_1}
+
\frac p{n_2}
\right).
\]
Consequently, with probability at least \(1-\eta/3\),
\[
\|b\|_2
\le
K_0
\frac{\sqrt{n_{\mathrm{eff}}}}{\sqrt d}
\eta^{-1/2}
\left(
\frac p{n_1}
+
\frac p{n_2}
\right)
=
B_W(\eta).
\]

We next control the diagonal-studentization error. Define
\[
r_D
=
\max_{1\le k\le d}
\left|
\frac{\widehat D_{kk}}{\widetilde D_{kk}}-1
\right|.
\]
By the studentizer concentration bound \eqref{eq:supp-studentizer-hp}, after
increasing the constant \(K_0\) if necessary,
\[
\mathbb P
\left(
r_D>
K_0\sqrt{\frac{d}{n_0\eta}}
\right)
\le
\frac{\eta}{3}.
\]
Therefore, with probability at least \(1-\eta/3\),
\[
r_D\le B_D(\eta),
\]
where
\[
B_D(\eta)=K_0\left(\frac{d}{n_0\eta}\right)^{1/2}.
\]

Define the event
\begin{equation}
\label{eq:supp-transfer-event}
\mathcal E_\eta
=
\left\{
\|a\|_2\le B_T(\eta)
\right\}
\cap
\left\{
\|b\|_2\le B_W(\eta)
\right\}
\cap
\left\{
r_D\le B_D(\eta)
\right\}.
\end{equation}
By the preceding bounds and the union bound,
\begin{equation}
\label{eq:supp-transfer-event-prob}
\mathbb P(\mathcal E_\eta)\ge 1-\eta.
\end{equation}

We now prove the deterministic perturbation bound on the event
\(\mathcal E_\eta\) in \eqref{eq:supp-transfer-event}.
Define
\[
E
=
\widetilde D^{1/2}\widehat D^{-1}\widetilde D^{1/2}.
\]
Since \(E\) is diagonal,
\[
E_{kk}
=
\frac{\widetilde D_{kk}}{\widehat D_{kk}}.
\]
Moreover,
\[
\widehat T_n
=
\frac1d\widehat W^\top\widehat D^{-1}\widehat W
=
(a+b)^\top E(a+b),
\]
whereas
\[
\widetilde T_n=a^\top a.
\]

On \(\mathcal E_\eta\), if
\[
B_D(\eta)\le \frac12,
\]
then
\[
\frac12
\le
\frac{\widehat D_{kk}}{\widetilde D_{kk}}
\le
\frac32
\]
for every \(1\le k\le d\). Hence \(\widehat D\) is positive definite, and
\[
|E_{kk}-1|
=
\left|
\frac{\widetilde D_{kk}}{\widehat D_{kk}}-1
\right|
=
\frac{
\left|
\widehat D_{kk}/\widetilde D_{kk}-1
\right|
}{
\widehat D_{kk}/\widetilde D_{kk}
}
\le
2r_D.
\]
Thus
\[
\|E-I_d\|_{\mathrm{op}}\le 2r_D.
\]

We decompose
\[
\widehat T_n-\widetilde T_n
=
(a+b)^\top E(a+b)-a^\top a
\]
\[
=
(a+b)^\top(E-I_d)(a+b)
+
\left(
\|a+b\|_2^2-\|a\|_2^2
\right).
\]
Therefore,
\[
|\widehat T_n-\widetilde T_n|
\le
\left|
(a+b)^\top(E-I_d)(a+b)
\right|
+
\left|
\|a+b\|_2^2-\|a\|_2^2
\right|.
\]
The first term satisfies
\[
\left|
(a+b)^\top(E-I_d)(a+b)
\right|
\le
\|E-I_d\|_{\mathrm{op}}\|a+b\|_2^2
\le
2r_D(\|a\|_2+\|b\|_2)^2.
\]
For the second term,
\[
\|a+b\|_2^2-\|a\|_2^2
=
2a^\top b+\|b\|_2^2,
\]
so
\[
\left|
\|a+b\|_2^2-\|a\|_2^2
\right|
\le
2\|a\|_2\|b\|_2+\|b\|_2^2.
\]
Combining the two inequalities gives
\begin{equation}
\label{eq:supp-transfer-perturbation}
|\widehat T_n-\widetilde T_n|
\le
2r_D(\|a\|_2+\|b\|_2)^2
+
2\|a\|_2\|b\|_2
+
\|b\|_2^2.
\end{equation}
On \(\mathcal E_\eta\),
\[
\|a\|_2\le B_T(\eta),
\qquad
\|b\|_2\le B_W(\eta),
\qquad
r_D\le B_D(\eta).
\]
Hence \eqref{eq:supp-transfer-perturbation} gives
\[
|\widehat T_n-\widetilde T_n|
\le
2B_D(\eta)\{B_T(\eta)+B_W(\eta)\}^2
+
2B_T(\eta)B_W(\eta)
+
B_W^2(\eta).
\]
By the definition of \(\Delta(\eta)\),
\[
|\widehat T_n-\widetilde T_n|
\le
\Delta(\eta)
\]
on \(\mathcal E_\eta\). Since \eqref{eq:supp-transfer-event-prob} holds, the
first claim follows.

It remains to verify the explicit rate in the final statement. Take
\[
\eta_n=n_0^{-1/8}.
\]
Then
\[
B_T(\eta_n)
=
K_0\eta_n^{-1/4}
=
K_0 n_0^{1/32},
\]
and
\[
B_D(\eta_n)
=
K_0
\left(
\frac{d}{n_0\eta_n}
\right)^{1/2}
=
K_0
\left(
\frac{d}{n_0^{7/8}}
\right)^{1/2}.
\]
Moreover, since
\[
n_{\mathrm{eff}}\le n_0,
\qquad
n_1,n_2\ge n_0,
\]
we have
\[
\frac p{n_1}+\frac p{n_2}
\le
\frac{2p}{n_0}.
\]
Because
\[
d=\frac{p(p+1)}2,
\]
we also have
\[
\frac p{\sqrt d}\le K_0.
\]
Therefore,
\[
B_W(\eta_n)
=
K_0
\frac{\sqrt{n_{\mathrm{eff}}}}{\sqrt d}
n_0^{1/16}
\left(
\frac p{n_1}
+
\frac p{n_2}
\right)
\le
K_0
\frac{\sqrt{n_0}}{\sqrt d}
n_0^{1/16}
\frac p{n_0}
\le
K_0 n_0^{-7/16}.
\]

Substituting these bounds into
\[
\Delta(\eta_n)
=
2B_D(\eta_n)\{B_T(\eta_n)+B_W(\eta_n)\}^2
+
2B_T(\eta_n)B_W(\eta_n)
+
B_W^2(\eta_n),
\]
we obtain
\[
\Delta(\eta_n)
\le
K_0 B_D(\eta_n)B_T^2(\eta_n)
+
K_0 B_D(\eta_n)B_T(\eta_n)B_W(\eta_n)
\]
\[
\qquad
+
K_0 B_D(\eta_n)B_W^2(\eta_n)
+
K_0 B_T(\eta_n)B_W(\eta_n)
+
K_0 B_W^2(\eta_n).
\]
The first leading term is
\[
B_D(\eta_n)B_T^2(\eta_n)
=
K_0
\left(
\frac{d}{n_0^{7/8}}
\right)^{1/2}
n_0^{1/16}
=
K_0
\left(
\frac{d}{n_0^{3/4}}
\right)^{1/2}.
\]
The second leading term is
\[
B_T(\eta_n)B_W(\eta_n)
\le
K_0 n_0^{1/32}n_0^{-7/16}
=
K_0 n_0^{-13/32}.
\]
For the remaining terms, using
\[
B_D(\eta_n)\le \frac12
\]
and \(B_T(\eta_n)\ge 1\) for large \(n_0\), we have
\[
B_D(\eta_n)B_T(\eta_n)B_W(\eta_n)
\le
K_0 B_T(\eta_n)B_W(\eta_n)
\le
K_0 n_0^{-13/32},
\]
and
\[
B_D(\eta_n)B_W^2(\eta_n)
\le
K_0 B_W^2(\eta_n)
\le
K_0 n_0^{-7/8}.
\]
Also,
\[
B_W^2(\eta_n)\le K_0 n_0^{-7/8}.
\]
Since
\[
n_0^{-7/8}\le n_0^{-13/32}
\]
for \(n_0\ge1\), all remaining terms are dominated by
\[
K_0
\left[
\left(
\frac{d}{n_0^{3/4}}
\right)^{1/2}
+
n_0^{-13/32}
\right].
\]
Therefore,
\[
\Delta(\eta_n)
\le
K_0
\left[
\left(
\frac{d}{n_0^{3/4}}
\right)^{1/2}
+
n_0^{-13/32}
\right].
\]
Since the first part of the proposition gives
\[
\mathbb P
\left(
|\widehat T_n-\widetilde T_n|
\le
\Delta(\eta_n)
\right)
\ge
1-\eta_n
=
1-n_0^{-1/8},
\]
we conclude that, with probability at least \(1-n_0^{-1/8}\),
\[
|\widehat T_n-\widetilde T_n|
\le
K_0
\left[
\left(
\frac{d}{n_0^{3/4}}
\right)^{1/2}
+
n_0^{-13/32}
\right].
\]
This completes the proof.
\end{proof}

\begin{proof}[Proof of Corollary~\ref{cor:feasible_gaussian}]
Let
\[
F_{\widehat T}(r)
=
\mathbb P(\widehat T_n\le r),
\qquad
F_{\widetilde T}(r)
=
\mathbb P(\widetilde T_n\le r),
\]
and
\[
G(r)
=
\mathbb P(T_Z\le r).
\]
By Theorem~\ref{thm:oracle_gaussian},
\[
\sup_{r\ge0}
\left|
F_{\widetilde T}(r)-G(r)
\right|
\le
\varepsilon_{\mathrm{or}},
\]
where
\[
\varepsilon_{\mathrm{or}}
=
K
\left\{
\kappa(R)^{3/4}\delta_1^{1/4}
+
\kappa(R)^{2/3}\delta_2^{1/3}
\right\}.
\]

Next, by Proposition~\ref{prop:sample_oracle_transfer}, if
\[
B_D(\eta)\le \frac12,
\]
then the event
\[
\mathcal E_\eta
=
\left\{
|\widehat T_n-\widetilde T_n|\le \Delta(\eta)
\right\}
\]
satisfies
\[
\mathbb P(\mathcal E_\eta)\ge 1-\eta.
\]

We first derive an upper bound for
\[
F_{\widehat T}(r)-G(r).
\]
For any fixed \(r\ge0\),
\[
\mathbb P(\widehat T_n\le r)
=
\mathbb P(\widehat T_n\le r,\mathcal E_\eta)
+
\mathbb P(\widehat T_n\le r,\mathcal E_\eta^c).
\]
Since
\[
\mathbb P(\mathcal E_\eta^c)\le \eta,
\]
we have
\[
\mathbb P(\widehat T_n\le r)
\le
\mathbb P(\widehat T_n\le r,\mathcal E_\eta)
+
\eta.
\]
On the event \(\mathcal E_\eta\),
\[
|\widehat T_n-\widetilde T_n|\le \Delta(\eta).
\]
Therefore,
\[
\widehat T_n\le r
\quad\Longrightarrow\quad
\widetilde T_n\le r+\Delta(\eta).
\]
Hence
\[
\mathbb P(\widehat T_n\le r)
\le
\mathbb P(\widetilde T_n\le r+\Delta(\eta))
+
\eta.
\]
Using the oracle Gaussian approximation,
\[
\mathbb P(\widetilde T_n\le r+\Delta(\eta))
\le
\mathbb P(T_Z\le r+\Delta(\eta))
+
\varepsilon_{\mathrm{or}}.
\]
Thus
\[
\mathbb P(\widehat T_n\le r)
\le
\mathbb P(T_Z\le r+\Delta(\eta))
+
\varepsilon_{\mathrm{or}}
+
\eta.
\]
By the definition of the anti-concentration modulus
\[
\omega_Z(a)
=
\sup_{s\ge0}\mathbb P(s<T_Z\le s+a),
\]
we have
\[
\mathbb P(T_Z\le r+\Delta(\eta))
\le
\mathbb P(T_Z\le r)
+
\omega_Z\{\Delta(\eta)\}.
\]
Therefore,
\[
\mathbb P(\widehat T_n\le r)-\mathbb P(T_Z\le r)
\le
\varepsilon_{\mathrm{or}}
+
\eta
+
\omega_Z\{\Delta(\eta)\}.
\]

We next derive the reverse inequality. On \(\mathcal E_\eta\),
\[
\widetilde T_n\le r-\Delta(\eta)
\quad\Longrightarrow\quad
\widehat T_n\le r.
\]
If \(r-\Delta(\eta)<0\), the event
\[
\{\widetilde T_n\le r-\Delta(\eta)\}
\]
is interpreted as the empty event. Hence, for every \(r\ge0\),
\[
\mathbb P(\widehat T_n\le r)
\ge
\mathbb P(\widetilde T_n\le r-\Delta(\eta),\mathcal E_\eta).
\]
Therefore,
\[
\mathbb P(\widehat T_n\le r)
\ge
\mathbb P(\widetilde T_n\le r-\Delta(\eta))
-
\mathbb P(\mathcal E_\eta^c).
\]
Since
\[
\mathbb P(\mathcal E_\eta^c)\le \eta,
\]
we get
\[
\mathbb P(\widehat T_n\le r)
\ge
\mathbb P(\widetilde T_n\le r-\Delta(\eta))
-
\eta.
\]
Using the oracle Gaussian approximation again,
\[
\mathbb P(\widetilde T_n\le r-\Delta(\eta))
\ge
\mathbb P(T_Z\le r-\Delta(\eta))
-
\varepsilon_{\mathrm{or}}.
\]
Thus
\[
\mathbb P(\widehat T_n\le r)
\ge
\mathbb P(T_Z\le r-\Delta(\eta))
-
\varepsilon_{\mathrm{or}}
-
\eta.
\]
It follows that
\[
\mathbb P(T_Z\le r)-\mathbb P(\widehat T_n\le r)
\le
\mathbb P(T_Z\le r)
-
\mathbb P(T_Z\le r-\Delta(\eta))
+
\varepsilon_{\mathrm{or}}
+
\eta.
\]
By the definition of \(\omega_Z\),
\[
\mathbb P(T_Z\le r)
-
\mathbb P(T_Z\le r-\Delta(\eta))
\le
\omega_Z\{\Delta(\eta)\}.
\]
Therefore,
\[
\mathbb P(T_Z\le r)-\mathbb P(\widehat T_n\le r)
\le
\varepsilon_{\mathrm{or}}
+
\eta
+
\omega_Z\{\Delta(\eta)\}.
\]

Combining the two one-sided bounds, we obtain, for every \(r\ge0\),
\[
\left|
\mathbb P(\widehat T_n\le r)
-
\mathbb P(T_Z\le r)
\right|
\le
\varepsilon_{\mathrm{or}}
+
\eta
+
\omega_Z\{\Delta(\eta)\}.
\]
Taking the supremum over \(r\ge0\) gives
\[
\sup_{r\ge0}
\left|
\mathbb P(\widehat T_n\le r)
-
\mathbb P(T_Z\le r)
\right|
\le
\varepsilon_{\mathrm{or}}
+
\eta
+
\omega_Z\{\Delta(\eta)\}.
\]

It remains to use the Gaussian quadratic-form anti-concentration bound. Under
Assumption~\ref{ass:gaussian_reference_geometry},
\[
T_Z
=
\frac1d\|Z\|_2^2,
\qquad
Z\sim N(0,R).
\]
Applying Lemma~\ref{lem:gaussian_anticoncentration} (the
Götze--Naumov--Spokoiny--Ulyanov anti-concentration bound, as recorded in
\citealp{fang_large-dimensional_2021}, Lemma~5.10) with \(\Sigma=R\) and
\(\mu=0\), we obtain, for any \(\varepsilon>0\),
\[
\sup_{a\ge 0}\mathbb P(a\le \|Z\|_2^2\le a+\varepsilon)\le K\kappa(R)\varepsilon.
\]
Translating to the rescaled statistic \(T_Z=d^{-1}\|Z\|_2^2\) via the change of
variable \(a=rd\), \(\varepsilon=ad\), gives
\[
\omega_Z(a)
=
\sup_{r\ge0}\mathbb P(r<T_Z\le r+a)
=
\sup_{a'\ge0}\mathbb P(a'<\|Z\|_2^2\le a'+ad)
\le
K\kappa(R)\cdot ad
=
K\kappa(R)d a.
\]
Taking \(a=\Delta(\eta)\), we obtain
\[
\omega_Z\{\Delta(\eta)\}
\le
K\kappa(R)d\Delta(\eta).
\]
Hence
\[
\sup_{r\ge0}
\left|
\mathbb P(\widehat T_n\le r)
-
\mathbb P(T_Z\le r)
\right|
\le
\varepsilon_{\mathrm{or}}
+
\eta
+
K\kappa(R)d\Delta(\eta).
\]
This completes the proof.
\end{proof}


The multiplier bootstrap statistic is conditionally Gaussian given the data,
with conditional covariance matrix \(\widehat R\). To compare the bootstrap law
with the oracle Gaussian reference \(T_Z=d^{-1}\|Z\|_2^2\), \(Z\sim N(0,R)\),
we need to control the discrepancy between \(\widehat R\) and \(R\). The next
proposition provides a conservative nuclear-norm bound for this discrepancy.
The proof is technical and is given in the appendix.

\begin{proposition}[Bootstrap covariance consistency]
\label{prop:bootstrap_covariance_consistency}
Suppose Assumptions~\ref{ass:coordinate_high_moment} and
\ref{ass:coordinate_nondegeneracy} hold. Let
\[
\Gamma=\operatorname{Cov}(\widetilde W),
\qquad
R=\widetilde D^{-1/2}\Gamma\widetilde D^{-1/2}.
\]
Let
\[
\widehat\Gamma
=
n_{\mathrm{eff}}
\left\{
\frac1{n_1^2}\sum_{i=1}^{n_1}
\widehat h_i^{(1)}(\widehat h_i^{(1)})^\top
+
\frac1{n_2^2}\sum_{j=1}^{n_2}
\widehat h_j^{(2)}(\widehat h_j^{(2)})^\top
\right\},
\]
and
\[
\widehat R
=
\widehat D^{-1/2}\widehat\Gamma\widehat D^{-1/2}.
\]
Then there exists a constant \(K_R>0\), depending only on the constants in
the assumptions, such that for every \(\eta\in(0,1)\),
\[
\mathbb P
\left(
\|\widehat R-R\|_{S_1}
>
K_R\sqrt{\frac{d^3}{n_0\eta}}
\right)
\le \eta,
\qquad
n_0=n_1\wedge n_2.
\]
Equivalently,
\[
\|\widehat R-R\|_{S_1}
=
O_{\mathbb P}
\left(
\sqrt{\frac{d^3}{n_0}}
\right).
\]
\end{proposition}

\begin{proof}
Throughout the proof, we work in the transformed model and suppress the
superscript \((A)\). Let
\[
U_i=X_i-\mu_1,
\qquad
V_j=Y_j-\mu_2.
\]
Define the oracle covariance-coordinate vectors
\[
H_i^{(1)}
=
\operatorname{vech}
\{U_iU_i^\top-\mathbb E(U_iU_i^\top)\},
\]
and
\[
H_j^{(2)}
=
\operatorname{vech}
\{V_jV_j^\top-\mathbb E(V_jV_j^\top)\}.
\]
Let
\[
\Gamma_1=\operatorname{Cov}(H_i^{(1)}),
\qquad
\Gamma_2=\operatorname{Cov}(H_j^{(2)}).
\]
Then
\[
\Gamma
=
\operatorname{Cov}(\widetilde W)
=
n_{\mathrm{eff}}
\left(
\frac{\Gamma_1}{n_1}
+
\frac{\Gamma_2}{n_2}
\right).
\]

We first control the covariance estimator before diagonal standardization.
For the first sample, define
\[
\bar H^{(1)}
=
\frac1{n_1}\sum_{i=1}^{n_1}H_i^{(1)}
\]
and
\[
\check\Gamma_1
=
\frac1{n_1}
\sum_{i=1}^{n_1}
(H_i^{(1)}-\bar H^{(1)})(H_i^{(1)}-\bar H^{(1)})^\top.
\]
Also define
\[
\widehat\Gamma_1
=
\frac1{n_1}
\sum_{i=1}^{n_1}
\widehat h_i^{(1)}(\widehat h_i^{(1)})^\top.
\]
The second-sample analogues \(\check\Gamma_2\) and \(\widehat\Gamma_2\) are
defined similarly.

We claim that
\[
\mathbb E\|\widehat\Gamma_1-\Gamma_1\|_{\mathrm F}^2
\le
\frac{Cd^2}{n_1},
\qquad
\mathbb E\|\widehat\Gamma_2-\Gamma_2\|_{\mathrm F}^2
\le
\frac{Cd^2}{n_2}.
\]
We prove the first bound; the second is identical.

Write
\[
\Gamma_{1,0}
=
\frac1{n_1}\sum_{i=1}^{n_1}
H_i^{(1)}(H_i^{(1)})^\top.
\]
Then
\[
\check\Gamma_1
=
\Gamma_{1,0}
-
\bar H^{(1)}(\bar H^{(1)})^\top.
\]
Therefore,
\[
\|\check\Gamma_1-\Gamma_1\|_{\mathrm F}^2
\le
2\|\Gamma_{1,0}-\Gamma_1\|_{\mathrm F}^2
+
2\|\bar H^{(1)}(\bar H^{(1)})^\top\|_{\mathrm F}^2.
\]
Since the \(H_i^{(1)}\)'s are i.i.d. centered,
\[
\mathbb E\|\Gamma_{1,0}-\Gamma_1\|_{\mathrm F}^2
=
\frac1{n_1}
\mathbb E
\left\|
H_i^{(1)}(H_i^{(1)})^\top-\Gamma_1
\right\|_{\mathrm F}^2.
\]
Moreover,
\[
\left\|
H_i^{(1)}(H_i^{(1)})^\top
\right\|_{\mathrm F}^2
=
\|H_i^{(1)}\|_2^4.
\]
By Assumption~\ref{ass:coordinate_high_moment}, for every coordinate \(k\),
\[
\mathbb E|H_{i,k}^{(1)}|^4\le K_0.
\]
Hence
\[
\mathbb E\|H_i^{(1)}\|_2^4
=
\mathbb E
\left(
\sum_{k=1}^d (H_{i,k}^{(1)})^2
\right)^2
\le
d\sum_{k=1}^d
\mathbb E|H_{i,k}^{(1)}|^4
\le
Cd^2.
\]
It follows that
\[
\mathbb E\|\Gamma_{1,0}-\Gamma_1\|_{\mathrm F}^2
\le
\frac{Cd^2}{n_1}.
\]
Also,
\[
\|\bar H^{(1)}(\bar H^{(1)})^\top\|_{\mathrm F}^2
=
\|\bar H^{(1)}\|_2^4.
\]
Using the same fourth-moment calculation for the sample mean of the
\(d\)-dimensional vector \(H_i^{(1)}\), together with
\(\mathbb E\|H_i^{(1)}\|_2^4\le Cd^2\), gives
\[
\mathbb E\|\bar H^{(1)}\|_2^4
\le
\frac{Cd^2}{n_1^2}.
\]
Thus
\[
\mathbb E\|\check\Gamma_1-\Gamma_1\|_{\mathrm F}^2
\le
\frac{Cd^2}{n_1}.
\]

We next compare the feasible residuals with the oracle centered residuals.
Let
\[
g_i^{(1)}
=
H_i^{(1)}-\bar H^{(1)}.
\]
Then
\[
\check\Gamma_1
=
\frac1{n_1}\sum_{i=1}^{n_1}
g_i^{(1)}(g_i^{(1)})^\top.
\]
By the sample-centering expansion,
\[
\widehat h_i^{(1)}
=
g_i^{(1)}+r_i^{(1)},
\]
where
\[
r_i^{(1)}
=
\operatorname{vech}
\left\{
-\bar U(U_i-\bar U)^\top
-
(U_i-\bar U)\bar U^\top
\right\}.
\]
Indeed, this follows from expanding
\[
(U_i-\bar U)(U_i-\bar U)^\top-\widehat\Sigma_1
\]
and comparing it with
\[
U_iU_i^\top-\frac1{n_1}\sum_{\ell=1}^{n_1}U_\ell U_\ell^\top.
\]

Therefore,
\[
\widehat\Gamma_1-\check\Gamma_1
=
\frac1{n_1}\sum_{i=1}^{n_1}
\left\{
g_i^{(1)}(r_i^{(1)})^\top
+
r_i^{(1)}(g_i^{(1)})^\top
+
r_i^{(1)}(r_i^{(1)})^\top
\right\}.
\]
Using
\[
\|xy^\top\|_{\mathrm F}=\|x\|_2\|y\|_2,
\]
we obtain
\[
\|\widehat\Gamma_1-\check\Gamma_1\|_{\mathrm F}
\le
2
\left(
\frac1{n_1}\sum_{i=1}^{n_1}\|g_i^{(1)}\|_2^2
\right)^{1/2}
\left(
\frac1{n_1}\sum_{i=1}^{n_1}\|r_i^{(1)}\|_2^2
\right)^{1/2}
+
\frac1{n_1}\sum_{i=1}^{n_1}\|r_i^{(1)}\|_2^2.
\]
Set
\[
L_1=
\frac1{n_1}\sum_{i=1}^{n_1}\|g_i^{(1)}\|_2^2,
\qquad
M_1=
\frac1{n_1}\sum_{i=1}^{n_1}\|r_i^{(1)}\|_2^2.
\]
Then
\[
\|\widehat\Gamma_1-\check\Gamma_1\|_{\mathrm F}^2
\le
K_0(L_1M_1+M_1^2).
\]

We bound \(L_1\) and \(M_1\). Since
\[
\mathbb E\|H_i^{(1)}\|_2^4\le Cd^2,
\]
we have
\[
\mathbb E L_1^2\le Cd^2.
\]
For \(M_1\), note that
\[
\|r_i^{(1)}\|_2
\le
\left\|
-\bar U(U_i-\bar U)^\top
-
(U_i-\bar U)\bar U^\top
\right\|_{\mathrm F}
\le
2\|\bar U\|_2\|U_i-\bar U\|_2.
\]
Therefore,
\[
M_1
\le
4\|\bar U\|_2^2
\left(
\frac1{n_1}\sum_{i=1}^{n_1}\|U_i-\bar U\|_2^2
\right).
\]
By Lemma~\ref{lem:coordinate_moments_imply_euclidean},
\[
\mathbb E\|\bar U\|_2^8
\le
K_0\left(\frac{p}{n_1}\right)^4,
\]
and
\[
\mathbb E
\left(
\frac1{n_1}\sum_{i=1}^{n_1}\|U_i-\bar U\|_2^2
\right)^4
\le
Cp^4.
\]
Consequently, by Cauchy--Schwarz,
\[
\mathbb E M_1^2
\le
K_0
\left(\mathbb E\|\bar U\|_2^8\right)^{1/2}
\left[
\mathbb E
\left(
\frac1{n_1}\sum_{i=1}^{n_1}\|U_i-\bar U\|_2^2
\right)^4
\right]^{1/2}
\le
K_0\frac{p^4}{n_1^2}.
\]
Since
\[
d=\frac{p(p+1)}2\asymp p^2,
\]
we have \(p^4\le Cd^2\), and hence
\[
\mathbb E M_1^2
\le
\frac{Cd^2}{n_1^2}.
\]
Therefore,
\[
\mathbb E(L_1M_1)
\le
(\mathbb E L_1^2)^{1/2}
(\mathbb E M_1^2)^{1/2}
\le
Cd\cdot \frac d{n_1}
=
\frac{Cd^2}{n_1}.
\]
Together with \(\mathbb E M_1^2\le Cd^2/n_1^2\), this gives
\[
\mathbb E\|\widehat\Gamma_1-\check\Gamma_1\|_{\mathrm F}^2
\le
\frac{Cd^2}{n_1}.
\]
Combining this with the bound for
\(\check\Gamma_1-\Gamma_1\), we conclude that
\[
\mathbb E\|\widehat\Gamma_1-\Gamma_1\|_{\mathrm F}^2
\le
\frac{Cd^2}{n_1}.
\]
Similarly,
\[
\mathbb E\|\widehat\Gamma_2-\Gamma_2\|_{\mathrm F}^2
\le
\frac{Cd^2}{n_2}.
\]

Now recall that
\[
\widehat\Gamma
=
n_{\mathrm{eff}}
\left(
\frac{\widehat\Gamma_1}{n_1}
+
\frac{\widehat\Gamma_2}{n_2}
\right),
\]
and
\[
\Gamma
=
n_{\mathrm{eff}}
\left(
\frac{\Gamma_1}{n_1}
+
\frac{\Gamma_2}{n_2}
\right).
\]
Thus
\[
\widehat\Gamma-\Gamma
=
n_{\mathrm{eff}}
\left[
\frac{\widehat\Gamma_1-\Gamma_1}{n_1}
+
\frac{\widehat\Gamma_2-\Gamma_2}{n_2}
\right].
\]
Using \((x+y)^2\le 2x^2+2y^2\),
\[
\|\widehat\Gamma-\Gamma\|_{\mathrm F}^2
\le
2n_{\mathrm{eff}}^2
\left[
\frac{\|\widehat\Gamma_1-\Gamma_1\|_{\mathrm F}^2}{n_1^2}
+
\frac{\|\widehat\Gamma_2-\Gamma_2\|_{\mathrm F}^2}{n_2^2}
\right].
\]
Taking expectations gives
\[
\mathbb E\|\widehat\Gamma-\Gamma\|_{\mathrm F}^2
\le
K_0 n_{\mathrm{eff}}^2 d^2
\left(
\frac1{n_1^3}
+
\frac1{n_2^3}
\right).
\]
Since \(n_{\mathrm{eff}}\le n_0=n_1\wedge n_2\),
\[
n_{\mathrm{eff}}^2
\left(
\frac1{n_1^3}
+
\frac1{n_2^3}
\right)
\le
\frac{K_0}{n_0}.
\]
Hence
\[
\mathbb E\|\widehat\Gamma-\Gamma\|_{\mathrm F}^2
\le
\frac{Cd^2}{n_0}.
\]
By Markov's inequality,
\begin{equation}
\label{eq:supp-bootstrap-gamma-hp}
\mathbb P
\left(
\|\widehat\Gamma-\Gamma\|_{\mathrm F}
>
K_0\frac{d}{\sqrt{n_0\eta}}
\right)
\le
\frac{\eta}{2}
\end{equation}
after increasing \(K_0\) if necessary.

We now pass from \(\widehat\Gamma-\Gamma\) to
\(\widehat R-R\). Let
\[
r_D
=
\max_{1\le k\le d}
\left|
\frac{\widehat D_{kk}}{\widetilde D_{kk}}-1
\right|.
\]
By the studentizer concentration bound \eqref{eq:supp-studentizer-hp},
\begin{equation}
\label{eq:supp-bootstrap-rd-hp}
\mathbb P
\left(
r_D>
K_0\sqrt{\frac{d}{n_0\eta}}
\right)
\le
\frac{\eta}{2}.
\end{equation}
On the event \(r_D\le 1/2\), the matrix \(\widehat D\) is positive definite and
\[
\|\widehat D^{-1/2}\|_{\mathrm{op}}\le K_0,
\qquad
\|\widetilde D^{-1/2}\|_{\mathrm{op}}\le K_0.
\]
Moreover,
\[
\left\|
\widehat D^{-1/2}\widetilde D^{1/2}-I_d
\right\|_{\mathrm{op}}
\le
Cr_D.
\]
Define
\[
B=\widehat D^{-1/2}\widetilde D^{1/2}.
\]
Then
\[
\widehat D^{-1/2}\Gamma\widehat D^{-1/2}
=
B
\left(
\widetilde D^{-1/2}\Gamma\widetilde D^{-1/2}
\right)
B
=
BRB.
\]
Therefore,
\begin{equation}
\label{eq:supp-bootstrap-R-decomposition}
\widehat R-R
=
\widehat D^{-1/2}(\widehat\Gamma-\Gamma)\widehat D^{-1/2}
+
(BRB-R).
\end{equation}

For the first term in \eqref{eq:supp-bootstrap-R-decomposition},
\[
\left\|
\widehat D^{-1/2}(\widehat\Gamma-\Gamma)\widehat D^{-1/2}
\right\|_{S_1}
\le
\|\widehat D^{-1/2}\|_{\mathrm{op}}^2
\|\widehat\Gamma-\Gamma\|_{S_1}.
\]
Since
\[
\|M\|_{S_1}\le \sqrt d\,\|M\|_{\mathrm F},
\]
we have
\[
\left\|
\widehat D^{-1/2}(\widehat\Gamma-\Gamma)\widehat D^{-1/2}
\right\|_{S_1}
\le
K_0\sqrt d\,\|\widehat\Gamma-\Gamma\|_{\mathrm F}.
\]

For the second term in \eqref{eq:supp-bootstrap-R-decomposition},
\[
BRB-R
=
(B-I_d)RB+R(B-I_d).
\]
Hence
\[
\|BRB-R\|_{S_1}
\le
\|B-I_d\|_{\mathrm{op}}\|R\|_{S_1}\|B\|_{\mathrm{op}}
+
\|R\|_{S_1}\|B-I_d\|_{\mathrm{op}}.
\]
Because \(\Gamma=\operatorname{Cov}(\widetilde W)\) is positive semidefinite
and \(\widetilde D^{-1/2}\) is symmetric, the matrix
\(R=\widetilde D^{-1/2}\Gamma\widetilde D^{-1/2}\) is positive semidefinite.
Moreover, by definition of \(\widetilde D\), the diagonal entries of \(R\) are
all equal to one. Therefore,
\[
\|R\|_{S_1}=\operatorname{tr}(R)=d.
\]
Also,
\[
\|B-I_d\|_{\mathrm{op}}\le K_0r_D,
\qquad
\|B\|_{\mathrm{op}}\le K_0.
\]
Therefore,
\[
\|BRB-R\|_{S_1}
\le
Cd r_D.
\]
Combining the two bounds,
\begin{equation}
\label{eq:supp-bootstrap-R-combined}
\|\widehat R-R\|_{S_1}
\le
K_0\sqrt d\,\|\widehat\Gamma-\Gamma\|_{\mathrm F}
+
Cd r_D.
\end{equation}

On the event where
\[
\|\widehat\Gamma-\Gamma\|_{\mathrm F}
\le
K_0\frac d{\sqrt{n_0\eta}}
\]
and
\[
r_D\le K_0\sqrt{\frac d{n_0\eta}},
\]
we therefore have, by \eqref{eq:supp-bootstrap-R-combined},
\[
\|\widehat R-R\|_{S_1}
\le
K_0\sqrt d\cdot \frac d{\sqrt{n_0\eta}}
+
Cd\sqrt{\frac d{n_0\eta}}
\le
K_0\sqrt{\frac{d^3}{n_0\eta}}.
\]
The probability of this event is at least \(1-\eta\). Hence
\[
\mathbb P
\left(
\|\widehat R-R\|_{S_1}
>
K_R\sqrt{\frac{d^3}{n_0\eta}}
\right)
\le
\eta
\]
for a sufficiently large constant \(K_R>0\). This proves the proposition.
\end{proof}

\begin{proof}[Proof of Theorem~\ref{thm:bootstrap_explicit}]
We first identify the conditional distribution of the multiplier statistic.
Recall that
\[
\widehat W^{\mathrm{mb}}
=
\sqrt{n_{\mathrm{eff}}}
\left\{
\frac1{n_1}\sum_{i=1}^{n_1}e_i\widehat h_i^{(1)}
-
\frac1{n_2}\sum_{j=1}^{n_2}e_j'\widehat h_j^{(2)}
\right\},
\]
where, conditionally on the data, the multiplier variables
\[
e_1,\ldots,e_{n_1},e_1',\ldots,e_{n_2}'
\]
are independent standard normal random variables. Therefore,
conditionally on the data,
\[
\widehat W^{\mathrm{mb}}\mid\mathcal X
\sim
N(0,\widehat\Gamma),
\]
where
\[
\widehat\Gamma
=
n_{\mathrm{eff}}
\left\{
\frac1{n_1^2}\sum_{i=1}^{n_1}
\widehat h_i^{(1)}(\widehat h_i^{(1)})^\top
+
\frac1{n_2^2}\sum_{j=1}^{n_2}
\widehat h_j^{(2)}(\widehat h_j^{(2)})^\top
\right\}.
\]
Consequently,
\[
\widehat Z
:=
\widehat D^{-1/2}\widehat W^{\mathrm{mb}}
\mid\mathcal X
\sim
N(0,\widehat R),
\]
where
\[
\widehat R
=
\widehat D^{-1/2}\widehat\Gamma\widehat D^{-1/2}.
\]
Thus, conditionally on the data,
\[
\widehat T_n^{\mathrm{mb}}
=
\frac1d
\left\|
\widehat D^{-1/2}\widehat W^{\mathrm{mb}}
\right\|_2^2
=
\frac1d\|\widehat Z\|_2^2.
\]

The oracle Gaussian reference is
\[
T_Z
=
\frac1d\|Z\|_2^2,
\qquad
Z\sim N(0,R),
\]
where
\[
R=
\widetilde D^{-1/2}
\operatorname{Cov}(\widetilde W)
\widetilde D^{-1/2}.
\]
Hence the bootstrap approximation problem reduces to comparing two Gaussian
quadratic forms:
\[
\frac1d\|\widehat Z\|_2^2,
\qquad
\frac1d\|Z\|_2^2,
\]
with covariance matrices \(\widehat R\) and \(R\), respectively.

By Proposition~\ref{prop:bootstrap_covariance_consistency}, there exists a
constant \(K_R>0\) such that, for every \(\eta_2\in(0,1)\),
\[
\mathbb P
\left(
\|\widehat R-R\|_{S_1}
>
K_R\sqrt{\frac{d^3}{n_0\eta_2}}
\right)
\le
\eta_2.
\]
Define the event
\[
\mathcal B_{\eta_2}
=
\left\{
\|\widehat R-R\|_{S_1}
\le
K_R\sqrt{\frac{d^3}{n_0\eta_2}}
\right\}.
\]
Then
\[
\mathbb P(\mathcal B_{\eta_2})\ge 1-\eta_2.
\]

We now show that on \(\mathcal B_{\eta_2}\), the Gaussian reference geometry of
\(\widehat R\) is controlled by that of \(R\). Since
\[
\|\widehat R-R\|_{\mathrm{op}}
\le
\|\widehat R-R\|_{S_1}
\]
and
\[
\|\widehat R-R\|_{\mathrm F}
\le
\|\widehat R-R\|_{S_1},
\]
on \(\mathcal B_{\eta_2}\) we have
\[
\|\widehat R-R\|_{\mathrm F}
\le
K_R\sqrt{\frac{d^3}{n_0\eta_2}}.
\]
Let
\[
\widehat\lambda_1\ge\cdots\ge\widehat\lambda_d\ge0
\]
be the eigenvalues of \(\widehat R\), and define
\[
\widehat\Lambda_s
=
\left(
\sum_{\ell=s}^d\widehat\lambda_\ell^2
\right)^{1/2},
\qquad s=1,2.
\]
By Lemma~\ref{lem:hoffman_wielandt} (the Hoffman--Wielandt inequality applied to
\(\widehat R\) and \(R\)),
\[
\left[
\sum_{\ell=1}^d
(\widehat\lambda_\ell-\lambda_\ell)^2
\right]^{1/2}
\le
\|\widehat R-R\|_{\mathrm F}.
\]
Therefore, on \(\mathcal B_{\eta_2}\),
\[
|\widehat\Lambda_s-\Lambda_s|
\le
\|\widehat R-R\|_{\mathrm F}
\le
K_R\sqrt{\frac{d^3}{n_0\eta_2}},
\qquad s=1,2.
\]
By the assumed geometry stability condition,
\[
K_R\sqrt{\frac{d^3}{n_0\eta_2}}
\le
\frac12\Lambda_2(R).
\]
Since \(\Lambda_1(R)\ge\Lambda_2(R)\), the same bound also implies
\[
K_R\sqrt{\frac{d^3}{n_0\eta_2}}
\le
\frac12\Lambda_1(R).
\]
Hence, on \(\mathcal B_{\eta_2}\),
\[
\widehat\Lambda_1
\ge
\Lambda_1(R)
-
K_R\sqrt{\frac{d^3}{n_0\eta_2}}
\ge
\frac12\Lambda_1(R),
\]
and
\[
\widehat\Lambda_2
\ge
\Lambda_2(R)
-
K_R\sqrt{\frac{d^3}{n_0\eta_2}}
\ge
\frac12\Lambda_2(R).
\]
In particular,
\[
\widehat\Lambda_2>0
\]
on \(\mathcal B_{\eta_2}\). Moreover,
\[
\kappa(\widehat R)
=
(\widehat\Lambda_1\widehat\Lambda_2)^{-1/2}
\le
2(\Lambda_1(R)\Lambda_2(R))^{-1/2}
=
2\kappa(R).
\]

We now apply Lemma~\ref{lem:gaussian_comparison_quadratic} (the Gaussian
comparison inequality for squared norms). Conditional on the data and on the
event \(\mathcal B_{\eta_2}\), both covariance matrices \(\widehat R\) and
\(R\) satisfy the required nondegeneracy condition (\(\Lambda_2(\cdot)>0\) by
the preceding step). Lemma~\ref{lem:gaussian_comparison_quadratic} therefore
gives, after the rescaling \(\widehat T_n^{\mathrm{mb}}=d^{-1}\|\widehat Z\|_2^2\)
and \(T_Z=d^{-1}\|Z\|_2^2\),
\[
\sup_{r\ge0}
\left|
\mathbb P_\xi
\left(
\widehat T_n^{\mathrm{mb}}\le r\mid\mathcal X
\right)
-
\mathbb P(T_Z\le r)
\right|
\le
K
\{\kappa(\widehat R)+\kappa(R)\}
\|\widehat R-R\|_{S_1}.
\]
Using
\[
\kappa(\widehat R)\le 2\kappa(R)
\]
and
\[
\|\widehat R-R\|_{S_1}
\le
K_R\sqrt{\frac{d^3}{n_0\eta_2}}
\]
on \(\mathcal B_{\eta_2}\), we obtain
\[
\sup_{r\ge0}
\left|
\mathbb P_\xi
\left(
\widehat T_n^{\mathrm{mb}}\le r\mid\mathcal X
\right)
-
\mathbb P(T_Z\le r)
\right|
\le
K\kappa(R)
\sqrt{\frac{d^3}{n_0\eta_2}}.
\]
By the definition
\[
\rho_{\mathrm{mb}}(\eta_2)
=
K\kappa(R)
\sqrt{\frac{d^3}{n_0\eta_2}},
\]
this becomes
\[
\sup_{r\ge0}
\left|
\mathbb P_\xi
\left(
\widehat T_n^{\mathrm{mb}}\le r\mid\mathcal X
\right)
-
\mathbb P(T_Z\le r)
\right|
\le
\rho_{\mathrm{mb}}(\eta_2)
\]
on \(\mathcal B_{\eta_2}\).

Since
\[
\mathbb P(\mathcal B_{\eta_2})\ge 1-\eta_2,
\]
the asserted high-probability bootstrap approximation follows.
\end{proof}

\begin{proof}[Proof of Corollary~\ref{cor:size_validity}]
Let
\[
F_n(r)
=
\mathbb P_{H_0}(\widehat T_n\le r)
\]
be the unconditional distribution function of the feasible statistic under
\(H_0\), and let
\[
G(r)
=
\mathbb P(T_Z\le r)
\]
be the distribution function of the oracle Gaussian reference
\[
T_Z=\frac1d\|Z\|_2^2,
\qquad
Z\sim N(0,R).
\]
Also define the conditional bootstrap distribution function
\[
\widehat F_{\mathrm{mb}}(r)
=
\mathbb P_\xi
\left(
\widehat T_n^{\mathrm{mb}}\le r\mid\mathcal X
\right).
\]

By Corollary~\ref{cor:feasible_gaussian} and the assumptions of
Corollary~\ref{cor:size_validity}, there exists a deterministic error
\(\varepsilon_n\to0\), given by the bound in
\eqref{eq:feasible-gaussian-error}, such that
\begin{equation}
\label{eq:supp-size-feasible-kolmogorov}
\sup_{r\ge0}
|F_n(r)-G(r)|
\le
\varepsilon_n.
\end{equation}

Next, by Theorem~\ref{thm:bootstrap_explicit} and the assumptions of
Corollary~\ref{cor:size_validity}, there exists a deterministic error
\(\rho_n\to0\), given by the bound in \eqref{eq:bootstrap-error}, such that
with probability tending to one,
\begin{equation}
\label{eq:supp-size-bootstrap-kolmogorov}
\sup_{r\ge0}
\left|
\widehat F_{\mathrm{mb}}(r)-G(r)
\right|
\le
\rho_n.
\end{equation}
Let
\begin{equation}
\label{eq:supp-size-bootstrap-event}
\mathcal B_n
=
\left\{
\sup_{r\ge0}
\left|
\widehat F_{\mathrm{mb}}(r)-G(r)
\right|
\le \rho_n
\right\}.
\end{equation}
Then
\[
\mathbb P(\mathcal B_n)\to1.
\]
Let \(\tau_n=\mathbb P(\mathcal B_n^c)\).  Then \(\tau_n\to0\).

We first verify that the oracle Gaussian reference has well-behaved quantiles.
Since
\[
T_Z=\frac1d\|Z\|_2^2,
\qquad
Z\sim N(0,R),
\]
we may diagonalize \(R\) and write
\[
T_Z
\overset{d}{=}
\frac1d
\sum_{\ell=1}^d
\lambda_\ell \xi_\ell^2,
\]
where
\(\xi_1,\ldots,\xi_d\) are independent standard normal random variables and
\(\lambda_1,\ldots,\lambda_d\) are the eigenvalues of \(R\). By
Assumption~\ref{ass:gaussian_reference_geometry},
\[
\Lambda_2
=
\left(
\sum_{\ell=2}^d\lambda_\ell^2
\right)^{1/2}
>0.
\]
Thus the Gaussian quadratic form is nondegenerate. In particular, \(G\) is
continuous on \([0,\infty)\). Moreover, for every \(u\in(0,1)\), if
\[
q_u=\inf\{r\ge0:G(r)\ge u\},
\]
then
\[
G(q_u)=u.
\]

Let
\[
\widehat q_u
=
\inf\{r\ge0:\widehat F_{\mathrm{mb}}(r)\ge u\}
\]
be the conditional bootstrap \(u\)-quantile. By definition,
\[
\widehat c_{1-\alpha}^{\mathrm{mb}}
=
\widehat q_{1-\alpha}.
\]
On the event \(\mathcal B_n\) in \eqref{eq:supp-size-bootstrap-event}, the
Kolmogorov bound \eqref{eq:supp-size-bootstrap-kolmogorov} implies the
quantile sandwich
\begin{equation}
\label{eq:supp-size-quantile-sandwich}
q_{1-\alpha-\rho_n}
\le
\widehat c_{1-\alpha}^{\mathrm{mb}}
\le
q_{1-\alpha+\rho_n},
\end{equation}
for all sufficiently large \(n\), since \(\rho_n\to0\).

Indeed, if
\[
r<q_{1-\alpha-\rho_n},
\]
then
\[
G(r)<1-\alpha-\rho_n.
\]
Hence, on \(\mathcal B_n\),
\[
\widehat F_{\mathrm{mb}}(r)
\le
G(r)+\rho_n
<
1-\alpha.
\]
Therefore \(r\) cannot reach the bootstrap \((1-\alpha)\)-quantile, giving
\[
q_{1-\alpha-\rho_n}
\le
\widehat c_{1-\alpha}^{\mathrm{mb}}.
\]
Similarly, since
\[
G(q_{1-\alpha+\rho_n})
=
1-\alpha+\rho_n,
\]
we have, on \(\mathcal B_n\),
\[
\widehat F_{\mathrm{mb}}(q_{1-\alpha+\rho_n})
\ge
G(q_{1-\alpha+\rho_n})-\rho_n
=
1-\alpha.
\]
Therefore
\[
\widehat c_{1-\alpha}^{\mathrm{mb}}
\le
q_{1-\alpha+\rho_n}.
\]

We now bound the rejection probability from above.  By the left inequality in
\eqref{eq:supp-size-quantile-sandwich}, on \(\mathcal B_n\),
\[
\left\{
\widehat T_n>\widehat c_{1-\alpha}^{\mathrm{mb}}
\right\}
\subset
\left\{
\widehat T_n>q_{1-\alpha-\rho_n}
\right\}
\]
Hence
\[
\mathbb P_{H_0}
\left(
\widehat T_n>\widehat c_{1-\alpha}^{\mathrm{mb}}
\right)
\le
\mathbb P_{H_0}
\left(
\widehat T_n>q_{1-\alpha-\rho_n}
\right)
+
\mathbb P(\mathcal B_n^c).
\]
Since \(\mathbb P(\mathcal B_n^c)=\tau_n\),
\[
\mathbb P_{H_0}
\left(
\widehat T_n>\widehat c_{1-\alpha}^{\mathrm{mb}}
\right)
\le
1-F_n(q_{1-\alpha-\rho_n})
+
\tau_n.
\]
Using \eqref{eq:supp-size-feasible-kolmogorov}, we obtain
\[
1-F_n(q_{1-\alpha-\rho_n})
\le
1-G(q_{1-\alpha-\rho_n})
+
\varepsilon_n.
\]
Since
\[
G(q_{1-\alpha-\rho_n})
=
1-\alpha-\rho_n,
\]
we get
\[
1-G(q_{1-\alpha-\rho_n})
=
\alpha+\rho_n.
\]
Therefore,
\begin{equation}
\label{eq:supp-size-upper-bound}
\mathbb P_{H_0}
\left(
\widehat T_n>\widehat c_{1-\alpha}^{\mathrm{mb}}
\right)
\le
\alpha+\rho_n+\varepsilon_n+\tau_n.
\end{equation}
Because
\[
\rho_n\to0,
\qquad
\varepsilon_n\to0,
\qquad
\tau_n\to0,
\]
we have
\[
\limsup_{n\to\infty}
\mathbb P_{H_0}
\left(
\widehat T_n>\widehat c_{1-\alpha}^{\mathrm{mb}}
\right)
\le
\alpha.
\]

We next derive the lower bound.  By the right inequality in
\eqref{eq:supp-size-quantile-sandwich}, on \(\mathcal B_n\),
\[
\left\{
\widehat T_n>q_{1-\alpha+\rho_n}
\right\}
\subset
\left\{
\widehat T_n>\widehat c_{1-\alpha}^{\mathrm{mb}}
\right\}
\]
Therefore,
\[
\mathbb P_{H_0}
\left(
\widehat T_n>\widehat c_{1-\alpha}^{\mathrm{mb}}
\right)
\ge
\mathbb P_{H_0}
\left(
\widehat T_n>q_{1-\alpha+\rho_n}
\right)
-
\mathbb P(\mathcal B_n^c).
\]
Since \(\mathbb P(\mathcal B_n^c)=\tau_n\),
\[
\mathbb P_{H_0}
\left(
\widehat T_n>\widehat c_{1-\alpha}^{\mathrm{mb}}
\right)
\ge
1-F_n(q_{1-\alpha+\rho_n})
-
\tau_n.
\]
Again using \eqref{eq:supp-size-feasible-kolmogorov}, we get
\[
1-F_n(q_{1-\alpha+\rho_n})
\ge
1-G(q_{1-\alpha+\rho_n})
-
\varepsilon_n.
\]
Since
\[
G(q_{1-\alpha+\rho_n})
=
1-\alpha+\rho_n,
\]
we have
\[
1-G(q_{1-\alpha+\rho_n})
=
\alpha-\rho_n.
\]
Thus
\begin{equation}
\label{eq:supp-size-lower-bound}
\mathbb P_{H_0}
\left(
\widehat T_n>\widehat c_{1-\alpha}^{\mathrm{mb}}
\right)
\ge
\alpha-\rho_n-\varepsilon_n-\tau_n.
\end{equation}
Letting \(n\to\infty\), and using
\[
\rho_n\to0,
\qquad
\varepsilon_n\to0,
\qquad
\tau_n\to0,
\]
gives
\[
\liminf_{n\to\infty}
\mathbb P_{H_0}
\left(
\widehat T_n>\widehat c_{1-\alpha}^{\mathrm{mb}}
\right)
\ge
\alpha.
\]

Combining the upper and lower bounds, we conclude that
\[
\mathbb P_{H_0}
\left(
\widehat T_n>\widehat c_{1-\alpha}^{\mathrm{mb}}
\right)
\to
\alpha.
\]
This proves the corollary.
\end{proof}

\begin{proof}[Proof of Theorem~\ref{thm:power_consistency}]
Throughout the proof, we work in the transformed model and suppress the
superscript \((A)\).  Under \(H_1\), write
\[
\Delta_\Sigma
=
\operatorname{vech}\{A(\Sigma_1-\Sigma_2)A^\top\},
\qquad
\theta_n
=
\widetilde D^{-1/2}\Delta_\Sigma .
\]
The oracle covariance-difference vector can be decomposed as
\[
\widetilde W
=
\sqrt{n_{\mathrm{eff}}}\Delta_\Sigma
+
\widetilde W_0,
\]
where
\[
\widetilde W_0
=
\sqrt{n_{\mathrm{eff}}}
\left\{
\operatorname{vech}(\widetilde\Sigma_1-\Sigma_1)
-
\operatorname{vech}(\widetilde\Sigma_2-\Sigma_2)
\right\}
\]
is centered.  Define
\[
S_n=\widetilde D^{-1/2}\widetilde W_0.
\]
Then
\[
\widetilde D^{-1/2}\widetilde W
=
\sqrt{n_{\mathrm{eff}}}\theta_n+S_n.
\]

By construction, the diagonal entries of
\(\operatorname{Cov}(S_n)\) are equal to one.  Hence
\[
\mathbb E_{H_1}
\left(
\frac1d\|S_n\|_2^2
\right)
=
\frac1d
\operatorname{tr}\{\operatorname{Cov}(S_n)\}
=
1.
\]
Therefore,
\[
\frac1{\sqrt d}\|S_n\|_2=O_{\mathbb P}(1).
\]

By the primitive alternative condition, there exists
\(\mathcal I_n\subset\{1,\ldots,d\}\), with
\(|\mathcal I_n|=m_n\), such that
\[
\min_{k\in\mathcal I_n}|\theta_{n,k}|\ge a_n.
\]
Consequently,
\[
\frac1d
\left\|
\sqrt{n_{\mathrm{eff}}}\theta_n
\right\|_2^2
=
\frac{n_{\mathrm{eff}}}{d}\|\theta_n\|_2^2
\ge
n_{\mathrm{eff}}\frac{m_n}{d}a_n^2
=
g_n.
\]
Since \(g_n\to\infty\), we have
\[
\frac1{\sqrt d}
\left\|
\sqrt{n_{\mathrm{eff}}}\theta_n
\right\|_2
\ge
\sqrt{g_n}
\to\infty.
\]
Combining this deterministic lower bound with
\(d^{-1/2}\|S_n\|_2=O_{\mathbb P}(1)\), the triangle inequality gives
\[
\frac1{\sqrt d}
\left\|
\widetilde D^{-1/2}\widetilde W
\right\|_2
=
\frac1{\sqrt d}
\left\|
\sqrt{n_{\mathrm{eff}}}\theta_n+S_n
\right\|_2
\ge
\sqrt{g_n}-O_{\mathbb P}(1).
\]
Thus
\[
\widetilde T_n
=
\frac1d
\left\|
\widetilde D^{-1/2}\widetilde W
\right\|_2^2
\to\infty
\]
in probability under \(H_1\).

It remains to transfer this divergence from the oracle statistic to the feasible
statistic.  Let
\[
E_n
=
\widetilde D^{-1/2}(\widehat W-\widetilde W).
\]
The same sample-centering argument used in the proof of
Proposition~\ref{prop:sample_oracle_transfer} gives, under
Assumption~\ref{ass:coordinate_high_moment},
\[
\frac1{\sqrt d}\|E_n\|_2
=
O_{\mathbb P}
\left[
\sqrt{\frac{n_{\mathrm{eff}}}{d}}
\left(
\frac{p}{n_1}+\frac{p}{n_2}
\right)
\right].
\]
By the assumed growth condition,
\[
\sqrt{\frac{n_{\mathrm{eff}}}{d}}
\left(
\frac{p}{n_1}+\frac{p}{n_2}
\right)
=
o(\sqrt{g_n}),
\]
we obtain
\[
\frac1{\sqrt d}\|E_n\|_2=o_{\mathbb P}(\sqrt{g_n}).
\]

Next, by the coordinatewise studentizer consistency established in the
sample-to-oracle analysis,
\[
\max_{1\le k\le d}
\left|
\frac{\widehat D_k}{\widetilde D_k}-1
\right|
=
O_{\mathbb P}
\left(
\sqrt{\frac{d}{n_0}}
\right)
=
o_{\mathbb P}(1),
\]
where the last equality follows from \(d/n_0\to0\).  Hence, with probability
tending to one,
\[
\max_{1\le k\le d}
\left|
\frac{\widehat D_k}{\widetilde D_k}-1
\right|
\le \frac12.
\]
On this event,
\[
\frac{\widetilde D_k}{\widehat D_k}\ge \frac23
\qquad
\text{for all } k=1,\ldots,d.
\]
Therefore,
\[
\widehat T_n
=
\frac1d
\widehat W^\top \widehat D^{-1}\widehat W
=
\frac1d
\sum_{k=1}^d
\frac{\widetilde D_k}{\widehat D_k}
\left[
\widetilde D^{-1/2}\widehat W
\right]_k^2
\ge
\frac23
\frac1d
\left\|
\widetilde D^{-1/2}\widehat W
\right\|_2^2.
\]
Since
\[
\widetilde D^{-1/2}\widehat W
=
\widetilde D^{-1/2}\widetilde W+E_n
=
\sqrt{n_{\mathrm{eff}}}\theta_n+S_n+E_n,
\]
another application of the triangle inequality yields
\[
\frac1{\sqrt d}
\left\|
\widetilde D^{-1/2}\widehat W
\right\|_2
\ge
\sqrt{g_n}
-
\frac1{\sqrt d}\|S_n\|_2
-
\frac1{\sqrt d}\|E_n\|_2.
\]
The second term is \(O_{\mathbb P}(1)\), while the third term is
\(o_{\mathbb P}(\sqrt{g_n})\).  Since \(g_n\to\infty\),
\[
\frac1{\sqrt d}
\left\|
\widetilde D^{-1/2}\widehat W
\right\|_2
\to\infty
\]
in probability.  Consequently,
\[
\widehat T_n\to\infty
\]
in probability under \(H_1\).

It remains only to show that the bootstrap critical value does not diverge.
Conditionally on the data, the multiplier-bootstrap statistic has the form
\[
\widehat T_n^{\mathrm{mb}}
=
\frac1d
\left\|
\widehat D^{-1/2}\widehat W^{\mathrm{mb}}
\right\|_2^2.
\]
Let
\[
\widehat Z
=
\widehat D^{-1/2}\widehat W^{\mathrm{mb}}.
\]
Then, conditionally on the data,
\[
\widehat Z\mid\mathcal X\sim N(0,\widehat R),
\]
where \(\widehat R\) has diagonal entries equal to one by construction of
\(\widehat D\).  Hence
\[
\operatorname{tr}(\widehat R)=d
\]
and
\[
\mathbb E_\xi
\left(
\widehat T_n^{\mathrm{mb}}\mid\mathcal X
\right)
=
\frac1d\operatorname{tr}(\widehat R)
=
1.
\]
By conditional Markov's inequality, for any \(M>0\),
\[
\mathbb P_\xi
\left(
\widehat T_n^{\mathrm{mb}}>M
\mid\mathcal X
\right)
\le
\frac1M.
\]
Choose \(M>\alpha^{-1}\).  Then
\[
\mathbb P_\xi
\left(
\widehat T_n^{\mathrm{mb}}\le M
\mid\mathcal X
\right)
\ge
1-\frac1M
>
1-\alpha.
\]
Therefore, by the definition of the conditional \((1-\alpha)\)-quantile,
\[
\widehat c_{1-\alpha}^{\mathrm{mb}}\le M
\]
with probability tending to one.  Thus
\[
\widehat c_{1-\alpha}^{\mathrm{mb}}=O_{\mathbb P}(1).
\]

We have shown that
\[
\widehat T_n\to\infty
\quad\text{in probability},
\qquad
\widehat c_{1-\alpha}^{\mathrm{mb}}=O_{\mathbb P}(1).
\]
It follows that
\[
\mathbb P_{H_1}
\left(
\widehat T_n>\widehat c_{1-\alpha}^{\mathrm{mb}}
\right)
\to1.
\]

Finally, the right-tail multiplier-bootstrap \(p\)-value is
\[
p_n^{\mathrm{mb}}
=
\mathbb P_\xi
\left(
\widehat T_n^{\mathrm{mb}}\ge \widehat T_n
\mid\mathcal X
\right).
\]
By conditional Markov's inequality,
\[
p_n^{\mathrm{mb}}
\le
\frac{
\mathbb E_\xi(\widehat T_n^{\mathrm{mb}}\mid\mathcal X)
}{
\widehat T_n
}
=
\frac1{\widehat T_n}.
\]
Since \(\widehat T_n\to\infty\) in probability, we conclude that
\[
p_n^{\mathrm{mb}}\to0
\]
in probability.  This completes the proof.
\end{proof}

\begin{proof}[Proof of Proposition~\ref{prop:unbalanced_reduction}]
We work under \(H_0\) and suppress the superscript \((A)\).  Recall that
\[
\widetilde W
=
\sqrt{n_{\mathrm{eff}}}(\bar H_1-\bar H_2),
\qquad
W_{2|1}
=
-\sqrt{n_2}\bar H_2.
\]
Since the two samples are independent,
\[
\operatorname{Cov}(\widetilde W)
=
\frac{n_{\mathrm{eff}}}{n_1}\Gamma_1
+
\frac{n_{\mathrm{eff}}}{n_2}\Gamma_2.
\]
Therefore,
\[
\widetilde D
=
\alpha_nD_1+\beta_nD_2,
\]
where
\[
\alpha_n
=
\frac{n_{\mathrm{eff}}}{n_1}
=
\frac{n_2}{n_1+n_2},
\qquad
\beta_n
=
\frac{n_{\mathrm{eff}}}{n_2}
=
\frac{n_1}{n_1+n_2}.
\]
Writing
\[
\varepsilon_n=\frac{n_2}{n_1},
\]
we have
\[
\alpha_n=\frac{\varepsilon_n}{1+\varepsilon_n}=O(\varepsilon_n),
\qquad
\beta_n=\frac1{1+\varepsilon_n}=1+O(\varepsilon_n).
\]
Hence
\[
\frac{n_{\mathrm{eff}}}{n_2}=\beta_n=1+O(\varepsilon_n).
\]

By Assumption~\ref{ass:coordinate_nondegeneracy},
\[
D_{g,k}\ge c_H
\qquad
\text{for all }g=1,2,\quad k=1,\ldots,d.
\]
By Assumption~\ref{ass:coordinate_high_moment}, there exists a constant
\(C<\infty\) such that
\[
D_{g,k}\le C
\qquad
\text{for all }g=1,2,\quad k=1,\ldots,d.
\]
Indeed, for a covariance coordinate \(k=(a,b)\),
\[
D_{g,k}
=
\operatorname{Var}(Q_{ga}Q_{gb})
\le
\mathbb E(Q_{ga}^2Q_{gb}^2)
\le
\left(\mathbb E|Q_{ga}|^4\right)^{1/2}
\left(\mathbb E|Q_{gb}|^4\right)^{1/2}
\le C,
\]
where \(Q_g\) denotes a generic transformed centered observation from group
\(g\), that is, \(Q_1=U_i\) and \(Q_2=V_j\).

Consequently,
\[
\max_{1\le k\le d}
\left|
\frac{\widetilde D_k}{D_{2,k}}-1
\right|
=
\max_{1\le k\le d}
\left|
\alpha_n
\left(
\frac{D_{1,k}}{D_{2,k}}-1
\right)
\right|
\le
C\alpha_n
=
O(\varepsilon_n).
\]
It follows that
\[
\left\|
\widetilde D^{-1/2}D_2^{1/2}
\right\|_{\mathrm{op}}
=
O(1),
\qquad
\left\|
\widetilde D^{-1/2}D_2^{1/2}-I_d
\right\|_{\mathrm{op}}
=
O(\varepsilon_n).
\]

Next,
\[
\widetilde W-W_{2|1}
=
\sqrt{n_{\mathrm{eff}}}\bar H_1
+
(\sqrt{n_2}-\sqrt{n_{\mathrm{eff}}})\bar H_2.
\]
Since \(\bar H_1\) and \(\bar H_2\) are independent and centered,
\[
\begin{aligned}
\mathbb E
\left\|
D_2^{-1/2}(\widetilde W-W_{2|1})
\right\|_2^2
&=
\frac{n_{\mathrm{eff}}}{n_1}
\operatorname{tr}
\left(
D_2^{-1/2}\Gamma_1D_2^{-1/2}
\right) \\
&\quad+
\frac{(\sqrt{n_2}-\sqrt{n_{\mathrm{eff}}})^2}{n_2}
\operatorname{tr}
\left(
D_2^{-1/2}\Gamma_2D_2^{-1/2}
\right).
\end{aligned}
\]
Using the uniform bounds on \(D_{2,k}\) and \(\Gamma_{1,kk}\),
\[
\operatorname{tr}
\left(
D_2^{-1/2}\Gamma_1D_2^{-1/2}
\right)
=
\sum_{k=1}^d
\frac{\Gamma_{1,kk}}{D_{2,k}}
\le Cd.
\]
Moreover,
\[
\operatorname{tr}
\left(
D_2^{-1/2}\Gamma_2D_2^{-1/2}
\right)
=
\sum_{k=1}^d
\frac{\Gamma_{2,kk}}{D_{2,k}}
=
d.
\]
Therefore,
\[
\frac1d
\mathbb E
\left\|
D_2^{-1/2}(\widetilde W-W_{2|1})
\right\|_2^2
\le
C\alpha_n+(1-\sqrt{\beta_n})^2.
\]
Since
\[
\alpha_n=O(\varepsilon_n),
\qquad
1-\sqrt{\beta_n}=O(\varepsilon_n),
\]
we obtain
\[
\frac1d
\mathbb E
\left\|
D_2^{-1/2}(\widetilde W-W_{2|1})
\right\|_2^2
\le
C\varepsilon_n.
\]
By Markov's inequality,
\[
\frac1{\sqrt d}
\left\|
D_2^{-1/2}(\widetilde W-W_{2|1})
\right\|_2
=
O_{\mathbb P}(\varepsilon_n^{1/2}).
\]

Now decompose
\[
\widetilde D^{-1/2}\widetilde W
-
D_2^{-1/2}W_{2|1}
=
\widetilde D^{-1/2}(\widetilde W-W_{2|1})
+
(\widetilde D^{-1/2}-D_2^{-1/2})W_{2|1}.
\]
For the first term,
\[
\begin{aligned}
\frac1{\sqrt d}
\left\|
\widetilde D^{-1/2}(\widetilde W-W_{2|1})
\right\|_2
&\le
\left\|
\widetilde D^{-1/2}D_2^{1/2}
\right\|_{\mathrm{op}}
\frac1{\sqrt d}
\left\|
D_2^{-1/2}(\widetilde W-W_{2|1})
\right\|_2 \\
&=
O_{\mathbb P}(\varepsilon_n^{1/2}).
\end{aligned}
\]
For the second term,
\[
(\widetilde D^{-1/2}-D_2^{-1/2})W_{2|1}
=
(\widetilde D^{-1/2}D_2^{1/2}-I_d)D_2^{-1/2}W_{2|1}.
\]
Furthermore,
\[
\mathbb E
\left[
\frac1d
\left\|
D_2^{-1/2}W_{2|1}
\right\|_2^2
\right]
=
\frac1d
\operatorname{tr}
\left(
D_2^{-1/2}\Gamma_2D_2^{-1/2}
\right)
=
1.
\]
Thus
\[
\frac1{\sqrt d}
\left\|
D_2^{-1/2}W_{2|1}
\right\|_2
=
O_{\mathbb P}(1),
\]
and the second term is \(O_{\mathbb P}(\varepsilon_n)\) after normalization by
\(d^{-1/2}\).  Combining the two terms gives
\[
\frac1{\sqrt d}
\left\|
\widetilde D^{-1/2}\widetilde W
-
D_2^{-1/2}W_{2|1}
\right\|_2
=
O_{\mathbb P}(\varepsilon_n^{1/2}).
\]

Finally, let
\[
A_n=\widetilde D^{-1/2}\widetilde W,
\qquad
B_n=D_2^{-1/2}W_{2|1}.
\]
We have shown that
\[
d^{-1/2}\|A_n-B_n\|_2
=
O_{\mathbb P}(\varepsilon_n^{1/2}).
\]
Also,
\[
d^{-1/2}\|B_n\|_2
=
O_{\mathbb P}(1),
\]
and hence
\[
d^{-1/2}\|A_n\|_2
=
O_{\mathbb P}(1).
\]
Using
\[
\left|
\|A_n\|_2^2-\|B_n\|_2^2
\right|
\le
\|A_n-B_n\|_2
\left(
\|A_n\|_2+\|B_n\|_2
\right),
\]
we obtain
\[
\left|
\widetilde T_n-T_{2|1}^{\mathrm{or}}
\right|
=
O_{\mathbb P}(\varepsilon_n^{1/2}).
\]
This proves the proposition.
\end{proof}

\begin{proof}[Proof of Corollary~\ref{cor:feasible_unbalanced}]
We prove that the feasible simplified statistic is close to its oracle version.
Write
\[
\widehat W_{2|1}
=
\sqrt{n_2}\operatorname{vech}(\widehat\Sigma_1-\widehat\Sigma_2),
\qquad
W_{2|1}
=
-\sqrt{n_2}\bar H_2.
\]
Under \(H_0\), \(\Sigma_1=\Sigma_2\), and
\[
W_{2|1}
=
\sqrt{n_2}\operatorname{vech}(\Sigma_1-\widetilde\Sigma_2).
\]
Therefore,
\[
\widehat W_{2|1}-W_{2|1}
=
\sqrt{n_2}
\operatorname{vech}
\left[
(\widehat\Sigma_1-\Sigma_1)
+
(\widetilde\Sigma_2-\widehat\Sigma_2)
\right].
\]
Using the sample-centering identity
\[
\widehat\Sigma_g-\widetilde\Sigma_g
=
-(\bar X_g-\mu_g)(\bar X_g-\mu_g)^\top,
\qquad g=1,2,
\]
we obtain
\[
\widehat W_{2|1}-W_{2|1}
=
\sqrt{n_2}
\operatorname{vech}
\left[
(\widetilde\Sigma_1-\Sigma_1)
-
(\bar X-\mu_1)(\bar X-\mu_1)^\top
+
(\bar Y-\mu_2)(\bar Y-\mu_2)^\top
\right].
\]

We first control this vector after \(D_2^{-1/2}\)-standardization.  The leading
term is
\[
\sqrt{n_2}
\operatorname{vech}(\widetilde\Sigma_1-\Sigma_1)
=
\sqrt{n_2}\bar H_1.
\]
As in the proof of Proposition~\ref{prop:unbalanced_reduction},
\[
\mathbb E
\left[
\frac1d
\left\|
D_2^{-1/2}\sqrt{n_2}\bar H_1
\right\|_2^2
\right]
=
\frac{n_2}{n_1}
\frac1d
\operatorname{tr}
(D_2^{-1/2}\Gamma_1D_2^{-1/2})
\le
C\varepsilon_n.
\]
Hence
\[
\frac1{\sqrt d}
\left\|
D_2^{-1/2}\sqrt{n_2}\bar H_1
\right\|_2
=
O_{\mathbb P}(\varepsilon_n^{1/2}).
\]

For the sample-centering terms, use
\[
\|\operatorname{vech}(aa^\top)\|_2
\le
\|aa^\top\|_{\mathrm F}
=
\|a\|_2^2 .
\]
By Lemma~\ref{lem:coordinate_moments_imply_euclidean},
\[
\|\bar X-\mu_1\|_2^2
=
O_{\mathbb P}\left(\frac p{n_1}\right),
\qquad
\|\bar Y-\mu_2\|_2^2
=
O_{\mathbb P}\left(\frac p{n_2}\right).
\]
Since \(D_{2,k}\ge c_H\) and \(d=p(p+1)/2\), we have
\[
\frac{\sqrt{n_2}}{\sqrt d}
\left\|
D_2^{-1/2}
\operatorname{vech}
\{(\bar X-\mu_1)(\bar X-\mu_1)^\top\}
\right\|_2
=
O_{\mathbb P}
\left(
\frac{\sqrt{n_2}}{\sqrt d}\frac p{n_1}
\right)
=
O_{\mathbb P}
\left(
\frac{\sqrt{n_2}}{n_1}
\right),
\]
and
\[
\frac{\sqrt{n_2}}{\sqrt d}
\left\|
D_2^{-1/2}
\operatorname{vech}
\{(\bar Y-\mu_2)(\bar Y-\mu_2)^\top\}
\right\|_2
=
O_{\mathbb P}
\left(
\frac{\sqrt{n_2}}{\sqrt d}\frac p{n_2}
\right)
=
O_{\mathbb P}(n_2^{-1/2}).
\]
The first term is \(O_{\mathbb P}(\varepsilon_n^{1/2})\), and the second term
is absorbed by \(O_{\mathbb P}\{(d/n_2)^{1/2}\}\).  Therefore,
\[
\frac1{\sqrt d}
\left\|
D_2^{-1/2}(\widehat W_{2|1}-W_{2|1})
\right\|_2
=
O_{\mathbb P}
\left(
\varepsilon_n^{1/2}
+
\sqrt{\frac d{n_2}}
\right).
\]

Next, the one-sample version of the diagonal studentizer consistency gives
\[
\max_{1\le k\le d}
\left|
\frac{\widehat D_{2,kk}}{D_{2,k}}-1
\right|
=
O_{\mathbb P}
\left(
\sqrt{\frac d{n_2}}
\right).
\]
Since \(d/n_2\to0\), with probability tending to one,
\[
\left\|
\widehat D_2^{-1/2}D_2^{1/2}-I_d
\right\|_{\mathrm{op}}
=
O_{\mathbb P}
\left(
\sqrt{\frac d{n_2}}
\right),
\qquad
\left\|
\widehat D_2^{-1/2}D_2^{1/2}
\right\|_{\mathrm{op}}
=
O_{\mathbb P}(1).
\]

Now decompose
\[
\widehat D_2^{-1/2}\widehat W_{2|1}
-
D_2^{-1/2}W_{2|1}
=
\widehat D_2^{-1/2}D_2^{1/2}
D_2^{-1/2}(\widehat W_{2|1}-W_{2|1})
+
(\widehat D_2^{-1/2}D_2^{1/2}-I_d)
D_2^{-1/2}W_{2|1}.
\]
The first term is
\[
O_{\mathbb P}
\left(
\varepsilon_n^{1/2}
+
\sqrt{\frac d{n_2}}
\right)
\]
after normalization by \(d^{-1/2}\).  For the second term, recall from the proof
of Proposition~\ref{prop:unbalanced_reduction} that
\[
\frac1{\sqrt d}
\left\|
D_2^{-1/2}W_{2|1}
\right\|_2
=
O_{\mathbb P}(1).
\]
Hence the second term is
\[
O_{\mathbb P}
\left(
\sqrt{\frac d{n_2}}
\right)
\]
after normalization by \(d^{-1/2}\).  Therefore,
\[
\frac1{\sqrt d}
\left\|
\widehat D_2^{-1/2}\widehat W_{2|1}
-
D_2^{-1/2}W_{2|1}
\right\|_2
=
O_{\mathbb P}
\left(
\varepsilon_n^{1/2}
+
\sqrt{\frac d{n_2}}
\right).
\]

Let
\[
A_n=\widehat D_2^{-1/2}\widehat W_{2|1},
\qquad
B_n=D_2^{-1/2}W_{2|1}.
\]
The preceding display gives
\[
d^{-1/2}\|A_n-B_n\|_2
=
O_{\mathbb P}
\left(
\varepsilon_n^{1/2}
+
\sqrt{\frac d{n_2}}
\right).
\]
Moreover,
\[
d^{-1/2}\|B_n\|_2=O_{\mathbb P}(1),
\]
and the same is true for \(A_n\).  Therefore,
\[
\left|
\widehat T_{2|1}-T_{2|1}^{\mathrm{or}}
\right|
\le
\frac1{\sqrt d}\|A_n-B_n\|_2
\left(
\frac1{\sqrt d}\|A_n\|_2
+
\frac1{\sqrt d}\|B_n\|_2
\right),
\]
which yields
\[
\widehat T_{2|1}
-
T_{2|1}^{\mathrm{or}}
=
O_{\mathbb P}
\left(
\varepsilon_n^{1/2}
+
\sqrt{\frac d{n_2}}
\right).
\]

Finally, by the triangle inequality,
\[
|\widehat T_n-\widehat T_{2|1}|
\le
|\widehat T_n-\widetilde T_n|
+
|\widetilde T_n-T_{2|1}^{\mathrm{or}}|
+
|T_{2|1}^{\mathrm{or}}-\widehat T_{2|1}|.
\]
Using Proposition~\ref{prop:unbalanced_reduction} and the first part of this
corollary gives
\[
|\widehat T_n-\widehat T_{2|1}|
\le
|\widehat T_n-\widetilde T_n|
+
O_{\mathbb P}
\left(
\varepsilon_n^{1/2}
+
\sqrt{\frac d{n_2}}
\right).
\]
If \(\widehat T_n-\widetilde T_n=o_{\mathbb P}(1)\), \(\varepsilon_n\to0\),
and \(d/n_2\to0\), then
\[
\widehat T_n-\widehat T_{2|1}=o_{\mathbb P}(1).
\]
This completes the proof.
\end{proof}

\section{Additional Simulation Diagnostic}
\label{sec:supp_additional_simulation}

For completeness, we report the simulation diagnostic for the simplified
implementation in highly unbalanced designs.  This diagnostic is not used in
the main simulation section, which focuses on empirical size across dimensions
and power under dense alternatives.  The setting uses \(p=50\),
\(n_1=5000\), \(n_2=100\), and a Toeplitz null covariance with \(\rho=0.8\).
The purpose is to compare the full two-sample statistic with the simplified
statistic that treats the larger sample as a stable reference covariance.

\begin{table}[!htbp]
\centering
\footnotesize
\caption{Unbalanced-sample diagnostic with \(p=50\) and \(R=1000\).  The full and simplified columns report empirical size.  The relative statistic difference is \(|T_{\rm full}-T_{\rm simp}|/\{(|T_{\rm full}|+|T_{\rm simp}|)/2\}\), averaged over Monte Carlo replications; agreement is the percentage of replications in which the two procedures make the same rejection decision.}
\label{tab:supp_unbalanced_restart}
\begin{tabular}{lrrrrrrr}
\toprule
covariance & \(n_1\) & \(n_2\) & \(n_1/n_2\) & full & simplified & mean rel. \(\Delta_T\) & agreement \\
\midrule
Toeplitz(0.8) & 5000 & 100 & 50 & 0.049 & 0.056 & 2.4\% & 99.1\% \\
\bottomrule
\end{tabular}
\end{table}

The simplified statistic is close to the full statistic in this highly
unbalanced setting, and both rejection rates are near the nominal level.  This
supports the asymptotic simplification as a useful diagnostic approximation
when the information in the larger sample is comparatively stable.

We also examined finite-sample size as the sample-size ratio moves toward the
balanced case.  The design fixes \(n_2=60\), uses
\((n_1,n_2)\in\{(300,60),(180,60),(120,60),(90,60),(60,60)\}\), and sets
\(\rho_0=0.5\).  For each \(p\in\{100,200\}\), the simulation averages over
equicorrelation and Toeplitz covariance families and over Gaussian,
standardized centered \(\chi^2_1\), and standardized Laplace innovations.  The
full setting-level output, including all benchmark methods, is retained in the
formal simulation archive.  Table~\ref{tab:supp_unbalanced_ratio} reports the
average empirical sizes for the proposed target-transformed tests and the
identity implementation.

\begin{table}[!htbp]
\centering
\footnotesize
\caption{Average empirical size in the unbalanced-ratio diagnostic with
\(\rho_0=0.5\), \(n_2=60\), \(R=1000\), and \(B=500\).  Each row averages over
the two covariance families and the three innovation distributions.}
\label{tab:supp_unbalanced_ratio}
\begin{tabular}{@{}rrrrrrrr@{}}
\toprule
$p$ & $n_1$ & $n_2$ & $n_1/n_2$ & $G(0.3)$ & $G(0.5)$ & $G(0.7)$ & $I_p$ \\
\midrule
100 & 300 & 60 & 5.0 & 0.117 & 0.111 & 0.098 & 0.057 \\
100 & 180 & 60 & 3.0 & 0.095 & 0.097 & 0.090 & 0.050 \\
100 & 120 & 60 & 2.0 & 0.065 & 0.065 & 0.061 & 0.035 \\
100 & 90 & 60 & 1.5 & 0.057 & 0.060 & 0.059 & 0.030 \\
100 & 60 & 60 & 1.0 & 0.050 & 0.057 & 0.057 & 0.027 \\
\midrule
200 & 300 & 60 & 5.0 & 0.105 & 0.101 & 0.092 & 0.054 \\
200 & 180 & 60 & 3.0 & 0.089 & 0.089 & 0.083 & 0.048 \\
200 & 120 & 60 & 2.0 & 0.063 & 0.064 & 0.061 & 0.030 \\
200 & 90 & 60 & 1.5 & 0.057 & 0.055 & 0.055 & 0.029 \\
200 & 60 & 60 & 1.0 & 0.054 & 0.055 & 0.054 & 0.028 \\
\bottomrule
\end{tabular}

\end{table}

The ratio-grid diagnostic is consistent with the balanced and mildly
unbalanced size results when \(n_1/n_2\le 2\).  For stronger imbalance,
especially \(n_1/n_2\in\{3,5\}\) with \(n_2=60\), the target-transformed
procedures become liberal on average.  We therefore treat this experiment as a
finite-sample warning for severe imbalance rather than as part of the main
calibration evidence.

\bibliographystyle{agsm}
\bibliography{references}

\end{document}
